\documentclass[12pt,a4paper]{article}
\usepackage[utf8]{inputenc}
\usepackage[T1]{fontenc}
\usepackage[spanish]{babel}
\usepackage{geometry}
\geometry{margin=2.2cm}
\usepackage{enumitem}
\usepackage{booktabs}
\usepackage{xcolor}
\usepackage{fancyhdr}
\usepackage{titlesec}
\usepackage{mdframed}
\usepackage{longtable}
\usepackage{parskip}

\definecolor{azul}{RGB}{37,99,235}
\definecolor{verde}{RGB}{22,163,74}
\definecolor{rojo}{RGB}{220,38,38}
\definecolor{ambar}{RGB}{180,83,9}
\definecolor{gris}{RGB}{90,100,128}
\definecolor{fondogris}{RGB}{240,243,248}
\definecolor{fondoverde}{RGB}{240,253,244}

\pagestyle{fancy}
\fancyhf{}
\rhead{\small\color{gris} CAT · Métodos de Investigación · León \& Montero · 4.ª ed.}
\lhead{\small\color{gris} Banco de 100 ítems}
\cfoot{\small\thepage}
\renewcommand{\headrulewidth}{0.4pt}

\titleformat{\section}{\large\bfseries\color{azul}}{}{0em}{}[\vspace{-4pt}\rule{\linewidth}{0.8pt}\vspace{4pt}]
\titleformat{\subsection}{\normalsize\bfseries\color{gris}}{}{0em}{}

\setlist[enumerate,1]{label=\textbf{\Alph*.},leftmargin=*,itemsep=3pt,topsep=4pt}
\setlist[itemize,1]{leftmargin=*,itemsep=2pt,topsep=3pt}

\begin{document}

\begin{center}
{\LARGE\bfseries\color{azul} Test Adaptativo Informatizado (CAT)}\\[6pt]
{\Large Métodos de Investigación · León \& Montero (4.ª ed.)}\\[10pt]
{\large\bfseries Banco completo de 100 ítems}\\[4pt]
{\normalsize\color{gris} 7 categorías diagnósticas · 3 niveles de dificultad}
\end{center}
\vspace{8pt}
\section*{Distribución del banco}

\begin{center}
\begin{tabular}{llccccc}
\toprule
\textbf{Categoría} & \textbf{Código} & \textbf{Nivel 1} & \textbf{Nivel 2} & \textbf{Nivel 3} & \textbf{Total} & \textbf{Temas}\\
\midrule
Experimental mismo grupo & EXP-MG & 5 & 4 & 3 & 12 & T.7\\
Factorial & FACT & 5 & 5 & 3 & 13 & T.8\\
Cuasiexperimental & CUASI & 6 & 4 & 3 & 13 & T.10\\
Descriptivos & DES & 8 & 4 & 2 & 14 & T.4--5\\
Cualitativo & CUAL & 6 & 5 & 3 & 14 & T.12--15\\
Ex post facto & EPF & 9 & 5 & 2 & 16 & T.11\\
Experimental entre grupos & EXP-EG & 8 & 7 & 3 & 18 & T.6\\
\midrule
\textbf{Total} & & \textbf{47} & \textbf{34} & \textbf{19} & \textbf{100} &\\
\bottomrule
\end{tabular}
\end{center}

\bigskip
\noindent\textbf{Leyenda de niveles:} \textcolor{azul}{Nivel 1 = Básico} (identificación y reconocimiento) · \textcolor{ambar}{Nivel 2 = Intermedio} (comprensión de características específicas) · \textcolor{rojo}{Nivel 3 = Avanzado} (integración y aplicación)

\newpage

\section{Experimental mismo grupo (EXP-MG) · T.7}
{\small\color{gris} 12 ítems · Niveles: L1=5 · L2=4 · L3=3}

\begin{mdframed}[backgroundcolor=fondogris,linecolor=azul,linewidth=1pt,innerleftmargin=10pt,innerrightmargin=10pt,innertopmargin=7pt,innerbottommargin=7pt]
\textbf{Ítem 1 / m1} \hfill {\small\bfseries\color{azul} Nivel 1 · Básico}
\end{mdframed}

Un investigador estudia el efecto de dos estilos musicales (clásica vs. jazz) sobre la concentración. Todos los participantes completan la misma tarea de concentración dos veces: una con música clásica y otra con jazz. El orden en que se presentan los dos estilos varía entre participantes. ¿Cómo se denomina este tipo de diseño?

\begin{enumerate}
  \item[A.] Diseño entre grupos distintos, porque se comparan dos tipos de música
  \item[\textbf{\color{verde}B.}] \textbf{\color{verde} Diseño intrasujeto o de medidas repetidas, porque el mismo participante pasa por todas las condiciones} \hfill $\checkmark$
  \item[C.] Diseño cuasiexperimental, porque los grupos no están formados por aleatorización
  \item[D.] Diseño ex post facto, porque los estilos musicales son variables ya existentes
\end{enumerate}

\noindent\textbf{\color{azul}Explicación:} El diseño intrasujeto o de medidas repetidas (within-subjects design) se caracteriza porque cada participante es expuesto a todas las condiciones experimentales, actuando como su propio control. Esto elimina las diferencias individuales entre condiciones y aumenta la potencia estadística con el mismo número de participantes.

\noindent\textbf{\color{gris}Por qué son incorrectas las otras opciones:}
\begin{itemize}
  \item {\small A — Incorrecto: el diseño entre grupos implica que cada participante solo pertenece a una condición; aquí todos pasan por ambas condiciones.}
  \item {\small C — Incorrecto: el cuasiexperimental implica grupos preexistentes sin aleatorización; aquí hay manipulación experimental con todos los participantes.}
  \item {\small D — Incorrecto: en el ex post facto no hay manipulación deliberada; aquí el investigador asigna deliberadamente el orden de las condiciones.}
\end{itemize}
\vspace{8pt}\hrule\vspace{10pt}

\begin{mdframed}[backgroundcolor=fondogris,linecolor=azul,linewidth=1pt,innerleftmargin=10pt,innerrightmargin=10pt,innertopmargin=7pt,innerbottommargin=7pt]
\textbf{Ítem 2 / n05} \hfill {\small\bfseries\color{azul} Nivel 1 · Básico}
\end{mdframed}

Una investigadora estudia si escuchar música clásica o pop afecta la concentración. Todos los participantes realizan la misma tarea de concentración dos veces: una con música clásica y otra con pop. La mitad empieza con clásica y luego pop; la otra mitad al revés. ¿Qué tipo de diseño es este?

\begin{enumerate}
  \item[A.] Diseño entre grupos, porque hay dos tipos de música distintos
  \item[\textbf{\color{verde}B.}] \textbf{\color{verde} Diseño intrasujeto o de medidas repetidas, porque el mismo participante pasa por las dos condiciones} \hfill $\checkmark$
  \item[C.] Diseño factorial 2×2, porque se cruzan dos tipos de música con dos grupos de participantes
  \item[D.] Diseño ex post facto, porque los gustos musicales de los participantes ya existían antes
\end{enumerate}

\noindent\textbf{\color{azul}Explicación:} En el diseño intrasujeto cada participante pasa por todas las condiciones. Aquí todos los sujetos escuchan tanto música clásica como pop. Alternar el orden (contrabalanceo) sirve para controlar los efectos de práctica o fatiga.

\noindent\textbf{\color{gris}Por qué son incorrectas las otras opciones:}
\begin{itemize}
  \item {\small A — Incorrecto: en el diseño entre grupos cada participante pertenece a una sola condición; aquí todos pasan por las dos condiciones de música.}
  \item {\small C — Incorrecto: el diseño factorial cruza dos variables independientes distintas; aquí la única VI es el tipo de música.}
  \item {\small D — Incorrecto: en el ex post facto no hay manipulación; aquí la investigadora asigna deliberadamente el orden de las condiciones musicales.}
\end{itemize}
\vspace{8pt}\hrule\vspace{10pt}

\begin{mdframed}[backgroundcolor=fondogris,linecolor=azul,linewidth=1pt,innerleftmargin=10pt,innerrightmargin=10pt,innertopmargin=7pt,innerbottommargin=7pt]
\textbf{Ítem 3 / n12} \hfill {\small\bfseries\color{azul} Nivel 1 · Básico}
\end{mdframed}

Una investigadora estudia si la postura corporal (erguida vs. encorvada) afecta el estado de ánimo. Cada participante adopta primero una postura durante cinco minutos y luego la otra, y se mide el estado de ánimo después de cada una. El orden se contrabalancea. ¿Qué diseño es este?

\begin{enumerate}
  \item[A.] Diseño entre grupos, porque hay dos condiciones de postura
  \item[\textbf{\color{verde}B.}] \textbf{\color{verde} Diseño intrasujeto, porque cada participante pasa por las dos condiciones de postura} \hfill $\checkmark$
  \item[C.] Diseño factorial, porque hay dos posturas y dos momentos de medida
  \item[D.] Diseño ex post facto, porque la postura habitual de cada persona ya existía antes
\end{enumerate}

\noindent\textbf{\color{azul}Explicación:} Cuando el mismo participante pasa por todas las condiciones experimentales, el diseño es intrasujeto (medidas repetidas). El contrabalanceo del orden garantiza que los efectos de práctica o fatiga no favorezcan sistemáticamente a ninguna condición.

\noindent\textbf{\color{gris}Por qué son incorrectas las otras opciones:}
\begin{itemize}
  \item {\small A — Incorrecto: en el diseño entre grupos cada participante solo pertenece a una condición; aquí todos adoptan las dos posturas.}
  \item {\small C — Incorrecto: el diseño factorial implica cruzar dos variables independientes distintas; aquí solo hay una VI (el tipo de postura).}
  \item {\small D — Incorrecto: en el ex post facto no hay manipulación; aquí la investigadora asigna deliberadamente el orden de las posturas.}
\end{itemize}
\vspace{8pt}\hrule\vspace{10pt}

\begin{mdframed}[backgroundcolor=fondogris,linecolor=azul,linewidth=1pt,innerleftmargin=10pt,innerrightmargin=10pt,innertopmargin=7pt,innerbottommargin=7pt]
\textbf{Ítem 4 / mg06} \hfill {\small\bfseries\color{azul} Nivel 1 · Básico}
\end{mdframed}

Una investigadora quiere comprobar si hacer una pausa activa de cinco minutos mejora el rendimiento en tareas de atención. Todos los participantes realizan la tarea de atención en dos sesiones: una con pausa activa antes y otra sin pausa. El orden se contrabalancea. ¿Cuál es la principal ventaja de usar a los mismos participantes en las dos condiciones?

\begin{enumerate}
  \item[A.] Que no hace falta medir la variable dependiente dos veces
  \item[\textbf{\color{verde}B.}] \textbf{\color{verde} Que las diferencias individuales entre participantes quedan controladas porque cada persona actúa como su propio control} \hfill $\checkmark$
  \item[C.] Que se elimina la necesidad de asignar aleatoriamente a los participantes a las condiciones
  \item[D.] Que los efectos de práctica desaparecen automáticamente al usar el mismo grupo
\end{enumerate}

\noindent\textbf{\color{azul}Explicación:} En el diseño intrasujeto cada participante pasa por todas las condiciones, de modo que las diferencias individuales (inteligencia, motivación, experiencia previa) son constantes en todas las condiciones y no pueden explicar las diferencias observadas entre ellas. Esta es la ventaja fundamental del diseño de medidas repetidas.

\noindent\textbf{\color{gris}Por qué son incorrectas las otras opciones:}
\begin{itemize}
  \item {\small A — Incorrecto: en el diseño intrasujeto sí se mide la VD en cada condición; de hecho se obtienen tantas medidas por participante como condiciones haya.}
  \item {\small C — Incorrecto: la asignación aleatoria sigue siendo relevante en el diseño intrasujeto, aplicándose al orden de presentación de las condiciones.}
  \item {\small D — Incorrecto: los efectos de práctica no desaparecen solos; por eso se necesita el contrabalanceo para distribuirlos equitativamente entre condiciones.}
\end{itemize}
\vspace{8pt}\hrule\vspace{10pt}

\begin{mdframed}[backgroundcolor=fondogris,linecolor=azul,linewidth=1pt,innerleftmargin=10pt,innerrightmargin=10pt,innertopmargin=7pt,innerbottommargin=7pt]
\textbf{Ítem 5 / mg07} \hfill {\small\bfseries\color{azul} Nivel 1 · Básico}
\end{mdframed}

Un investigador estudia el efecto de tres niveles de iluminación (baja, media, alta) sobre la velocidad de lectura. Todos los participantes leen un texto bajo cada nivel de iluminación. Para que ninguna condición salga siempre en primer lugar, utiliza un cuadrado latino de 3×3 que asigna a tres grupos distintos los órdenes BMA, MAB y ABM. ¿Qué problema está controlando con esta estrategia?

\begin{enumerate}
  \item[A.] El sesgo de selección, que podría hacer que los participantes más lentos se concentraran en la condición de iluminación baja
  \item[\textbf{\color{verde}B.}] \textbf{\color{verde} El efecto de orden y el efecto de práctica, que podrían hacer que la condición presentada en segundo o tercer lugar se beneficiara de la experiencia acumulada} \hfill $\checkmark$
  \item[C.] La variabilidad en la longitud de los textos asignados a cada condición de iluminación
  \item[D.] La falta de aleatorización de los participantes a los grupos de lectura
\end{enumerate}

\noindent\textbf{\color{azul}Explicación:} En el diseño intrasujeto con más de dos condiciones, el cuadrado latino garantiza que cada condición aparezca el mismo número de veces en cada posición ordinal. Así, los efectos de orden (fatiga, práctica, aburrimiento) se distribuyen de forma equilibrada entre las condiciones y no favorecen ni perjudican a ninguna de ellas en particular.

\noindent\textbf{\color{gris}Por qué son incorrectas las otras opciones:}
\begin{itemize}
  \item {\small A — Incorrecto: el sesgo de selección es una amenaza en los diseños entre grupos, no en el intrasujeto donde todos los participantes pasan por todas las condiciones.}
  \item {\small C — Incorrecto: la longitud de los textos es un asunto del diseño del material, no de la técnica de contrabalanceo.}
  \item {\small D — Incorrecto: en el diseño intrasujeto no se aleatoriza a los participantes a condiciones distintas; se aleatoriza el orden en que cada participante recibe las condiciones.}
\end{itemize}
\vspace{8pt}\hrule\vspace{10pt}

\begin{mdframed}[backgroundcolor=fondogris,linecolor=ambar,linewidth=1pt,innerleftmargin=10pt,innerrightmargin=10pt,innertopmargin=7pt,innerbottommargin=7pt]
\textbf{Ítem 6 / m2} \hfill {\small\bfseries\color{ambar} Nivel 2 · Intermedio}
\end{mdframed}

Un investigador usa un diseño intrasujeto con tres condiciones (A, B, C). Para controlar los efectos de orden utiliza un cuadrado latino. ¿Qué garantiza exactamente el cuadrado latino en un diseño con tres condiciones?

\begin{enumerate}
  \item[A.] Que cada condición se aplique solo a un tercio de los participantes para evitar fatiga
  \item[\textbf{\color{verde}B.}] \textbf{\color{verde} Que cada condición aparezca exactamente una vez en cada posición ordinal (primera, segunda, tercera), distribuyendo los efectos de orden uniformemente} \hfill $\checkmark$
  \item[C.] Que el orden de las condiciones sea completamente aleatorio para cada participante
  \item[D.] Que los tres grupos de participantes reciban las condiciones en el mismo orden
\end{enumerate}

\noindent\textbf{\color{azul}Explicación:} El cuadrado latino es un diseño de contrabalanceo sistemático en el que, con n condiciones, se crean n órdenes distintos de tal modo que cada condición aparece exactamente una vez en cada posición (primera, segunda, ..., n). Esto garantiza que los efectos de posición se distribuyan de forma equilibrada entre todas las condiciones.

\noindent\textbf{\color{gris}Por qué son incorrectas las otras opciones:}
\begin{itemize}
  \item {\small A — Incorrecto: el cuadrado latino no divide a los participantes entre condiciones; todos los participantes reciben todas las condiciones.}
  \item {\small C — Incorrecto: la aleatorización del orden para cada participante es el contrabalanceo completo; el cuadrado latino es una variante sistemática y más eficiente.}
  \item {\small D — Incorrecto: si todos recibieran el mismo orden, no habría contrabalanceo alguno y los efectos de orden sesgarían los resultados.}
\end{itemize}
\vspace{8pt}\hrule\vspace{10pt}

\begin{mdframed}[backgroundcolor=fondogris,linecolor=ambar,linewidth=1pt,innerleftmargin=10pt,innerrightmargin=10pt,innertopmargin=7pt,innerbottommargin=7pt]
\textbf{Ítem 7 / m3} \hfill {\small\bfseries\color{ambar} Nivel 2 · Intermedio}
\end{mdframed}

En un experimento de aprendizaje verbal intrasujeto, una investigadora sospecha que el efecto de arrastre (carry-over) entre condiciones es irreversible: lo aprendido en la primera condición no puede "desaprenderse" para la segunda. ¿Cuál es la implicación metodológica de este problema?

\begin{enumerate}
  \item[A.] Que debe aumentar el número de participantes para compensar la variabilidad adicional
  \item[\textbf{\color{verde}B.}] \textbf{\color{verde} Que el diseño intrasujeto no es viable y debe optarse por un diseño entre grupos distintos} \hfill $\checkmark$
  \item[C.] Que el cuadrado latino de 4x4 resuelve el problema del arrastre irreversible entre condiciones
  \item[D.] Que debe añadirse un grupo control adicional que no reciba ninguna condición
\end{enumerate}

\noindent\textbf{\color{azul}Explicación:} El efecto de arrastre irreversible (p. ej. aprender algo en la condición A modifica permanentemente al sujeto) invalida el diseño intrasujeto, porque la segunda condición no se aplica a un estado "neutral" del participante. En estos casos, el diseño entre grupos independientes es la única alternativa viable.

\noindent\textbf{\color{gris}Por qué son incorrectas las otras opciones:}
\begin{itemize}
  \item {\small A — Incorrecto: aumentar la muestra no resuelve el problema de fondo: el orden de las condiciones contamina las medidas si el arrastre es irreversible.}
  \item {\small C — Incorrecto: el cuadrado latino distribuye los efectos de posición uniformemente, pero no elimina el arrastre irreversible; solo lo contrabalancea.}
  \item {\small D — Incorrecto: añadir un grupo control no resuelve el problema del arrastre dentro del diseño intrasujeto.}
\end{itemize}
\vspace{8pt}\hrule\vspace{10pt}

\begin{mdframed}[backgroundcolor=fondogris,linecolor=ambar,linewidth=1pt,innerleftmargin=10pt,innerrightmargin=10pt,innertopmargin=7pt,innerbottommargin=7pt]
\textbf{Ítem 8 / mg08} \hfill {\small\bfseries\color{ambar} Nivel 2 · Intermedio}
\end{mdframed}

Una investigadora usa un diseño intrasujeto con cuatro condiciones. Al analizar los datos detecta que los participantes rinden sistemáticamente mejor en la cuarta condición que en la primera, independientemente de cuál sea esa condición. ¿Cómo se llama este fenómeno y qué técnica lo controla?

\begin{enumerate}
  \item[A.] Efecto de fatiga; se controla aumentando el número de participantes en el estudio
  \item[\textbf{\color{verde}B.}] \textbf{\color{verde} Efecto de orden o de práctica; se controla mediante el contrabalanceo, que distribuye las condiciones en distintos órdenes entre los participantes} \hfill $\checkmark$
  \item[C.] Efecto de maduración; se controla añadiendo un grupo control sin ninguna condición experimental
  \item[D.] Sesgo de instrumentación; se controla calibrando el instrumento de medida antes de cada sesión
\end{enumerate}

\noindent\textbf{\color{azul}Explicación:} El efecto de práctica (o de orden) ocurre cuando la posición de una condición en la secuencia —no la condición en sí misma— afecta al rendimiento. El contrabalanceo equilibra estos efectos distribuyendo todas las condiciones en todas las posiciones posibles a lo largo de los grupos de participantes. Si el efecto de práctica es igual para todas las condiciones, no sesga la comparación entre ellas.

\noindent\textbf{\color{gris}Por qué son incorrectas las otras opciones:}
\begin{itemize}
  \item {\small A — Incorrecto: aumentar el número de participantes incrementa la potencia estadística pero no resuelve el problema del efecto de posición dentro del diseño intrasujeto.}
  \item {\small C — Incorrecto: la maduración es un cambio natural de los participantes con el tiempo; el efecto de práctica es un cambio provocado por la exposición repetida a la tarea, que es distinto.}
  \item {\small D — Incorrecto: el sesgo de instrumentación se refiere a cambios en el instrumento de medida entre momentos; aquí el problema es el orden de presentación de las condiciones.}
\end{itemize}
\vspace{8pt}\hrule\vspace{10pt}

\begin{mdframed}[backgroundcolor=fondogris,linecolor=ambar,linewidth=1pt,innerleftmargin=10pt,innerrightmargin=10pt,innertopmargin=7pt,innerbottommargin=7pt]
\textbf{Ítem 9 / mg09} \hfill {\small\bfseries\color{ambar} Nivel 2 · Intermedio}
\end{mdframed}

Una investigadora diseña un experimento intrasujeto para comparar dos condiciones. Sin embargo, al llegar a la segunda condición, muchos participantes ya han aprendido el procedimiento de la tarea y eso mejora su rendimiento independientemente de la condición. Aunque contrabalancea el orden, el efecto de aprendizaje contamina los resultados. ¿Qué solución metodológica es la más adecuada?

\begin{enumerate}
  \item[A.] Reducir el número de ensayos en la segunda condición para compensar el aprendizaje acumulado
  \item[\textbf{\color{verde}B.}] \textbf{\color{verde} Cambiar a un diseño entre grupos independientes, donde cada participante solo recibe una condición y el aprendizaje previo no contamina la comparación} \hfill $\checkmark$
  \item[C.] Añadir una tercera condición placebo para absorber los efectos de práctica antes de las dos condiciones de interés
  \item[D.] Aplicar las dos condiciones en el mismo día para reducir el tiempo entre medidas
\end{enumerate}

\noindent\textbf{\color{azul}Explicación:} Cuando el efecto de práctica o aprendizaje es tan intenso que el contrabalanceo no puede neutralizarlo (porque el aprendizaje acumulado en la primera condición siempre mejora el rendimiento en la segunda), la solución es cambiar al diseño entre grupos. Cada participante recibe solo una condición y no lleva consigo el aprendizaje de las otras.

\noindent\textbf{\color{gris}Por qué son incorrectas las otras opciones:}
\begin{itemize}
  \item {\small A — Incorrecto: reducir los ensayos en la segunda condición haría las condiciones no comparables en duración o carga de trabajo, lo que introduciría un nuevo confundidor.}
  \item {\small C — Incorrecto: añadir una condición placebo añade complejidad pero no elimina el efecto de práctica en las condiciones de interés; solo lo desplaza.}
  \item {\small D — Incorrecto: aplicar las dos condiciones el mismo día puede reducir ciertos efectos de maduración pero intensifica los efectos de fatiga y práctica acumulados dentro de la misma sesión.}
\end{itemize}
\vspace{8pt}\hrule\vspace{10pt}

\begin{mdframed}[backgroundcolor=fondogris,linecolor=rojo,linewidth=1pt,innerleftmargin=10pt,innerrightmargin=10pt,innertopmargin=7pt,innerbottommargin=7pt]
\textbf{Ítem 10 / m4} \hfill {\small\bfseries\color{rojo} Nivel 3 · Avanzado}
\end{mdframed}

Un investigador usa un diseño de medidas repetidas con cuatro condiciones. Analiza los datos con ANOVA de medidas repetidas y obtiene una F significativa. ¿Cuál es la ventaja estadística fundamental del ANOVA de medidas repetidas frente al ANOVA entre grupos, y qué supuesto adicional debe verificarse?

\begin{enumerate}
  \item[A.] Es más potente porque usa más participantes; supuesto: normalidad multivariada en cada condición
  \item[\textbf{\color{verde}B.}] \textbf{\color{verde} Elimina la variabilidad entre sujetos del error, aumentando la potencia; supuesto: esfericidad (homogeneidad de las varianzas de las diferencias entre condiciones)} \hfill $\checkmark$
  \item[C.] Permite comparar medias sin supuestos distribucionales; supuesto: homogeneidad de varianzas entre grupos
  \item[D.] No requiere análisis post hoc porque la F de medidas repetidas ya identifica entre qué condiciones hay diferencias
\end{enumerate}

\noindent\textbf{\color{azul}Explicación:} El ANOVA de medidas repetidas extrae del término error la variabilidad debida a las diferencias individuales entre sujetos (que en el diseño intrasujeto es constante a través de las condiciones), resultando en un error residual menor y mayor potencia estadística. El supuesto adicional es la esfericidad: que las varianzas de todas las diferencias entre pares de condiciones sean iguales.

\noindent\textbf{\color{gris}Por qué son incorrectas las otras opciones:}
\begin{itemize}
  \item {\small A — Incorrecto: el diseño de medidas repetidas usa los mismos participantes en todas las condiciones, no más participantes; la ventaja es la reducción del error, no el tamaño muestral.}
  \item {\small C — Incorrecto: el ANOVA de medidas repetidas sí requiere supuestos distribucionales; además, la homogeneidad de varianzas es un supuesto del ANOVA entre grupos, no del de medidas repetidas.}
  \item {\small D — Incorrecto: como el ANOVA entre grupos, el de medidas repetidas solo indica que hay diferencias significativas en algún punto; requiere contrastes o pruebas post hoc para identificar entre qué condiciones.}
\end{itemize}
\vspace{8pt}\hrule\vspace{10pt}

\begin{mdframed}[backgroundcolor=fondogris,linecolor=rojo,linewidth=1pt,innerleftmargin=10pt,innerrightmargin=10pt,innertopmargin=7pt,innerbottommargin=7pt]
\textbf{Ítem 11 / m5} \hfill {\small\bfseries\color{rojo} Nivel 3 · Avanzado}
\end{mdframed}

Un investigador diseña un experimento intrasujeto sobre el efecto del estrés en la memoria. Para inducir estrés usa el Trier Social Stress Test. Tras completar la condición de estrés, espera una semana (washout period) antes de aplicar la condición control. ¿Cuál es el propósito metodológico del período de lavado (washout)?

\begin{enumerate}
  \item[A.] Aumentar el número de medidas de la variable dependiente para mejorar la fiabilidad
  \item[\textbf{\color{verde}B.}] \textbf{\color{verde} Reducir el efecto de arrastre permitiendo que los efectos fisiológicos y psicológicos del estrés se disipen antes de la siguiente condición} \hfill $\checkmark$
  \item[C.] Garantizar que los participantes no recuerden el procedimiento de la primera condición
  \item[D.] Aleatorizar el orden de las condiciones entre los participantes de forma automática
\end{enumerate}

\noindent\textbf{\color{azul}Explicación:} El período de lavado (washout) es una técnica para reducir el efecto de arrastre en diseños intrasujeto: se introduce una pausa suficiente entre condiciones para que los efectos de la primera condición (fisiológicos, emocionales, de aprendizaje) se disipen antes de aplicar la siguiente. No elimina el efecto de arrastre irreversible, pero minimiza el arrastre temporal.

\noindent\textbf{\color{gris}Por qué son incorrectas las otras opciones:}
\begin{itemize}
  \item {\small A — Incorrecto: el período de lavado no añade medidas; es un intervalo sin medición entre condiciones.}
  \item {\small C — Incorrecto: el período de lavado no impide que los participantes recuerden el procedimiento; para eso se necesitarían materiales distintos o diseños entre grupos.}
  \item {\small D — Incorrecto: la aleatorización del orden se decide en el diseño; el período de lavado es un intervalo temporal que se implementa durante la recogida de datos.}
\end{itemize}
\vspace{8pt}\hrule\vspace{10pt}

\begin{mdframed}[backgroundcolor=fondogris,linecolor=rojo,linewidth=1pt,innerleftmargin=10pt,innerrightmargin=10pt,innertopmargin=7pt,innerbottommargin=7pt]
\textbf{Ítem 12 / mg10} \hfill {\small\bfseries\color{rojo} Nivel 3 · Avanzado}
\end{mdframed}

En un experimento de memoria intrasujeto con tres listas de palabras (condición A: palabras emocionales, condición B: palabras neutras, condición C: palabras técnicas), el investigador encuentra que los participantes recuerdan más palabras en la condición que realizan en segundo lugar, independientemente de cuál sea. Quiere confirmar estadísticamente si existe un efecto de posición. ¿Qué análisis añadiría?

\begin{enumerate}
  \item[A.] Una correlación de Pearson entre la puntuación en la tarea y el orden de presentación asignado a cada participante
  \item[\textbf{\color{verde}B.}] \textbf{\color{verde} Un ANOVA de medidas repetidas con la posición en la secuencia (1.ª, 2.ª, 3.ª) como factor intrasujeto adicional, para estimar el efecto de posición con independencia del tipo de condición} \hfill $\checkmark$
  \item[C.] Una prueba t de muestras independientes entre los participantes que empezaron por A y los que empezaron por B
  \item[D.] Una regresión logística con la condición como variable dependiente y la posición como predictor
\end{enumerate}

\noindent\textbf{\color{azul}Explicación:} Para cuantificar el efecto de posición, se puede modelar la posición en la secuencia como un factor adicional en el ANOVA de medidas repetidas. Si este factor resulta significativo, confirma que la posición per se afecta al recuerdo con independencia del tipo de material. Esto justifica la necesidad del contrabalanceo y permite estimar su magnitud.

\noindent\textbf{\color{gris}Por qué son incorrectas las otras opciones:}
\begin{itemize}
  \item {\small A — Incorrecto: la correlación de Pearson mide la relación lineal entre dos variables continuas; aquí la posición es categórica (1.ª, 2.ª, 3.ª) y la comparación implica varios grupos dentro del mismo sujeto.}
  \item {\small C — Incorrecto: la prueba t compara dos grupos independientes; no sirve para analizar el efecto de posición dentro de un diseño intrasujeto con múltiples condiciones.}
  \item {\small D — Incorrecto: la regresión logística se usa cuando la VD es categórica binaria; aquí la VD es continua (número de palabras recordadas) y la condición es una variable independiente, no dependiente.}
\end{itemize}
\vspace{8pt}\hrule\vspace{10pt}

\section{Factorial (FACT) · T.8}
{\small\color{gris} 13 ítems · Niveles: L1=5 · L2=5 · L3=3}

\begin{mdframed}[backgroundcolor=fondogris,linecolor=azul,linewidth=1pt,innerleftmargin=10pt,innerrightmargin=10pt,innertopmargin=7pt,innerbottommargin=7pt]
\textbf{Ítem 13 / f1} \hfill {\small\bfseries\color{azul} Nivel 1 · Básico}
\end{mdframed}

Una investigadora estudia simultáneamente el efecto del método de enseñanza (tradicional vs. activo) y del tamaño del grupo (pequeño vs. grande) sobre el rendimiento académico. Asigna aleatoriamente a los participantes a una de las cuatro combinaciones posibles. ¿Cómo se describe este diseño?

\begin{enumerate}
  \item[A.] Diseño entre grupos con cuatro niveles de una sola variable independiente
  \item[\textbf{\color{verde}B.}] \textbf{\color{verde} Diseño factorial 2×2 con dos variables independientes y cuatro condiciones experimentales} \hfill $\checkmark$
  \item[C.] Diseño cuasiexperimental porque el tamaño del grupo no puede manipularse éticamente
  \item[D.] Diseño intrasujeto porque cada participante pertenece a una sola condición
\end{enumerate}

\noindent\textbf{\color{azul}Explicación:} El diseño factorial incluye simultáneamente dos o más variables independientes (factores). Un diseño 2×2 tiene dos factores con dos niveles cada uno, generando 4 condiciones (celdas). Lo más valioso del diseño factorial es que permite estudiar no solo el efecto de cada factor por separado (efectos principales) sino también si el efecto de un factor depende del nivel del otro (interacción).

\noindent\textbf{\color{gris}Por qué son incorrectas las otras opciones:}
\begin{itemize}
  \item {\small A — Incorrecto: cuatro niveles de una sola VI sería un diseño unifactorial de 4 grupos, no un factorial.}
  \item {\small C — Incorrecto: si la investigadora asigna aleatoriamente a los participantes a los tamaños de grupo, el tamaño del grupo es una VI manipulada, no una variable natural.}
  \item {\small D — Incorrecto: en este diseño cada participante recibe UNA combinación de condiciones (entre sujetos), no todas.}
\end{itemize}
\vspace{8pt}\hrule\vspace{10pt}

\begin{mdframed}[backgroundcolor=fondogris,linecolor=azul,linewidth=1pt,innerleftmargin=10pt,innerrightmargin=10pt,innertopmargin=7pt,innerbottommargin=7pt]
\textbf{Ítem 14 / f2} \hfill {\small\bfseries\color{azul} Nivel 1 · Básico}
\end{mdframed}

En un diseño factorial 3×2, ¿cuántas condiciones experimentales (celdas) hay en total?

\begin{enumerate}
  \item[A.] 3
  \item[B.] 2
  \item[C.] 5
  \item[\textbf{\color{verde}D.}] \textbf{\color{verde} 6} \hfill $\checkmark$
\end{enumerate}

\noindent\textbf{\color{azul}Explicación:} En un diseño factorial, el número de condiciones (celdas) es el producto de los niveles de cada factor. Un diseño 3×2 tiene 3 niveles en el primer factor y 2 en el segundo: 3 × 2 = 6 condiciones. Un diseño 2×3×2 tendría 2×3×2 = 12 condiciones.

\noindent\textbf{\color{gris}Por qué son incorrectas las otras opciones:}
\begin{itemize}
  \item {\small A — Incorrecto: 3 es el número de niveles del primer factor, no el número total de celdas.}
  \item {\small B — Incorrecto: 2 es el número de niveles del segundo factor, no el número total de celdas.}
  \item {\small C — Incorrecto: 5 sería la suma de los niveles (3+2), no el producto, que es la operación correcta para calcular el número de celdas.}
\end{itemize}
\vspace{8pt}\hrule\vspace{10pt}

\begin{mdframed}[backgroundcolor=fondogris,linecolor=azul,linewidth=1pt,innerleftmargin=10pt,innerrightmargin=10pt,innertopmargin=7pt,innerbottommargin=7pt]
\textbf{Ítem 15 / n08} \hfill {\small\bfseries\color{azul} Nivel 1 · Básico}
\end{mdframed}

Un investigador estudia si el tipo de recompensa (verbal vs. material) y el tipo de tarea (creativa vs. rutinaria) afectan la motivación. Crea cuatro grupos: recompensa verbal + tarea creativa, recompensa verbal + tarea rutinaria, recompensa material + tarea creativa, y recompensa material + tarea rutinaria. ¿Qué diseño es este?

\begin{enumerate}
  \item[A.] Diseño experimental de cuatro grupos independientes con una sola variable independiente
  \item[\textbf{\color{verde}B.}] \textbf{\color{verde} Diseño factorial 2×2, porque cruza dos variables independientes con dos niveles cada una} \hfill $\checkmark$
  \item[C.] Diseño cuasiexperimental, porque los participantes eligieron el tipo de tarea que preferían
  \item[D.] Diseño intrasujeto, porque todos los participantes realizan los cuatro tipos de tarea
\end{enumerate}

\noindent\textbf{\color{azul}Explicación:} Un diseño factorial 2×2 tiene dos factores (variables independientes) con dos niveles cada uno, lo que genera 4 condiciones (celdas). Permite estudiar tanto el efecto independiente de cada factor como la posible interacción entre ellos.

\noindent\textbf{\color{gris}Por qué son incorrectas las otras opciones:}
\begin{itemize}
  \item {\small A — Incorrecto: si hubiera una sola VI con cuatro niveles sería un diseño unifactorial; aquí hay dos VIs cruzadas, lo que define al diseño factorial.}
  \item {\small C — Incorrecto: si los participantes eligen su condición no hay asignación controlada; en este diseño el investigador asigna las condiciones.}
  \item {\small D — Incorrecto: en el diseño intrasujeto cada participante pasa por todas las condiciones; aquí cada participante está en solo una de las cuatro celdas.}
\end{itemize}
\vspace{8pt}\hrule\vspace{10pt}

\begin{mdframed}[backgroundcolor=fondogris,linecolor=azul,linewidth=1pt,innerleftmargin=10pt,innerrightmargin=10pt,innertopmargin=7pt,innerbottommargin=7pt]
\textbf{Ítem 16 / fa07} \hfill {\small\bfseries\color{azul} Nivel 1 · Básico}
\end{mdframed}

Un investigador diseña un estudio en el que algunos participantes hacen ejercicio aeróbico y otros no, y además algunos reciben suplementación con omega-3 y otros no. Asigna aleatoriamente a los participantes a cada combinación. ¿Cuántas condiciones tiene este diseño y cómo se llama?

\begin{enumerate}
  \item[A.] Dos condiciones; diseño experimental de un factor con dos niveles
  \item[\textbf{\color{verde}B.}] \textbf{\color{verde} Cuatro condiciones; diseño factorial 2×2 con dos variables independientes} \hfill $\checkmark$
  \item[C.] Tres condiciones; diseño intrasujeto con tres niveles
  \item[D.] Seis condiciones; diseño factorial 3×2 con niveles desiguales
\end{enumerate}

\noindent\textbf{\color{azul}Explicación:} Un diseño factorial 2×2 tiene dos factores (ejercicio: sí/no; omega-3: sí/no) con dos niveles cada uno, lo que genera 2×2=4 condiciones: ejercicio+omega-3, ejercicio sin omega-3, sin ejercicio+omega-3, y sin ejercicio ni omega-3. El factorial permite estudiar si el efecto del ejercicio es diferente según si se toma omega-3 o no (interacción).

\noindent\textbf{\color{gris}Por qué son incorrectas las otras opciones:}
\begin{itemize}
  \item {\small A — Incorrecto: si solo hubiera un factor (por ejemplo, ejercicio sí/no) habría dos condiciones; pero al cruzar dos factores hay cuatro.}
  \item {\small C — Incorrecto: en el diseño intrasujeto los mismos participantes pasan por todas las condiciones; aquí cada participante está en solo una de las cuatro celdas.}
  \item {\small D — Incorrecto: el diseño 3×2 implica que uno de los factores tiene tres niveles; aquí ambos factores tienen exactamente dos niveles.}
\end{itemize}
\vspace{8pt}\hrule\vspace{10pt}

\begin{mdframed}[backgroundcolor=fondogris,linecolor=azul,linewidth=1pt,innerleftmargin=10pt,innerrightmargin=10pt,innertopmargin=7pt,innerbottommargin=7pt]
\textbf{Ítem 17 / fa08} \hfill {\small\bfseries\color{azul} Nivel 1 · Básico}
\end{mdframed}

Una investigadora estudia si el tipo de retroalimentación (positiva vs. negativa) afecta la motivación. Tras analizar los datos de su diseño entre grupos con dos condiciones, su director le sugiere añadir una segunda variable independiente: el tipo de tarea (creativa vs. rutinaria). ¿Qué ventaja principal añade incorporar esta segunda variable al diseño?

\begin{enumerate}
  \item[A.] Reduce el número total de participantes necesarios a la mitad
  \item[\textbf{\color{verde}B.}] \textbf{\color{verde} Permite detectar si el efecto de la retroalimentación sobre la motivación es el mismo para ambos tipos de tarea, o si depende del tipo de tarea (interacción)} \hfill $\checkmark$
  \item[C.] Convierte el diseño en intrasujeto, eliminando la variabilidad individual
  \item[D.] Garantiza que los grupos sean equivalentes antes del tratamiento
\end{enumerate}

\noindent\textbf{\color{azul}Explicación:} La ventaja exclusiva del diseño factorial es la posibilidad de detectar interacciones: saber si el efecto de una VI cambia según el nivel de otra. Si la retroalimentación positiva solo aumenta la motivación en tareas creativas pero no en las rutinarias, eso es una interacción que un diseño de un solo factor nunca podría revelar.

\noindent\textbf{\color{gris}Por qué son incorrectas las otras opciones:}
\begin{itemize}
  \item {\small A — Incorrecto: añadir un segundo factor no reduce los participantes; de hecho añade celdas y normalmente requiere más participantes para mantener el mismo tamaño muestral por celda.}
  \item {\small C — Incorrecto: añadir una segunda VI no cambia la estructura entre sujetos / intrasujeto del diseño; eso depende de cómo se asignen los participantes.}
  \item {\small D — Incorrecto: la equivalencia de los grupos se garantiza mediante la aleatorización, no por añadir factores al diseño.}
\end{itemize}
\vspace{8pt}\hrule\vspace{10pt}

\begin{mdframed}[backgroundcolor=fondogris,linecolor=ambar,linewidth=1pt,innerleftmargin=10pt,innerrightmargin=10pt,innertopmargin=7pt,innerbottommargin=7pt]
\textbf{Ítem 18 / f3} \hfill {\small\bfseries\color{ambar} Nivel 2 · Intermedio}
\end{mdframed}

En un diseño factorial 2×2, una investigadora obtiene estos resultados: con retroalimentación positiva, los estudiantes de primaria (M=7.8) y de secundaria (M=7.9) rinden casi igual. Con retroalimentación negativa, los de primaria rinden mucho peor (M=5.2) que los de secundaria (M=7.5). ¿Qué patrón describen estos resultados?

\begin{enumerate}
  \item[A.] Un efecto principal de la etapa educativa: los estudiantes de secundaria rinden mejor en general
  \item[B.] Un efecto principal del tipo de retroalimentación: la retroalimentación positiva mejora el rendimiento de todos
  \item[\textbf{\color{verde}C.}] \textbf{\color{verde} Una interacción: el efecto del tipo de retroalimentación sobre el rendimiento depende de la etapa educativa} \hfill $\checkmark$
  \item[D.] Dos efectos principales independientes sin interacción entre las variables
\end{enumerate}

\noindent\textbf{\color{azul}Explicación:} Hay interacción cuando el efecto de una VI sobre la VD cambia en función del nivel de otra VI. En este caso, la retroalimentación negativa perjudica mucho más a los estudiantes de primaria que a los de secundaria, mientras que la retroalimentación positiva produce efectos similares en ambos grupos. Este patrón diferencial es la definición de interacción.

\noindent\textbf{\color{gris}Por qué son incorrectas las otras opciones:}
\begin{itemize}
  \item {\small A — Incorrecto: si el efecto principal de la etapa fuera el único patrón, las diferencias entre primaria y secundaria serían constantes en ambos tipos de retroalimentación.}
  \item {\small B — Incorrecto: si solo hubiera efecto principal de la retroalimentación, ambas etapas educativas responderían de forma paralela a ambos tipos de feedback.}
  \item {\small D — Incorrecto: cuando hay interacción, los efectos principales no pueden interpretarse de forma aislada; los datos muestran claramente que el efecto del feedback depende de la etapa.}
\end{itemize}
\vspace{8pt}\hrule\vspace{10pt}

\begin{mdframed}[backgroundcolor=fondogris,linecolor=ambar,linewidth=1pt,innerleftmargin=10pt,innerrightmargin=10pt,innertopmargin=7pt,innerbottommargin=7pt]
\textbf{Ítem 19 / f4} \hfill {\small\bfseries\color{ambar} Nivel 2 · Intermedio}
\end{mdframed}

Una investigadora detecta una interacción significativa en su diseño factorial 2×2. Su director de tesis le sugiere interpretar únicamente los efectos principales e ignorar la interacción. ¿Cuál es el error metodológico de este consejo según León y Montero?

\begin{enumerate}
  \item[A.] Los efectos principales son siempre más relevantes que la interacción en un diseño factorial
  \item[\textbf{\color{verde}B.}] \textbf{\color{verde} Cuando la interacción es significativa, interpretar solo los efectos principales puede resultar engañoso porque el efecto de una VI depende del nivel de la otra} \hfill $\checkmark$
  \item[C.] Ignorar la interacción solo es problemático si el tamaño del efecto es mayor que .50
  \item[D.] Los efectos principales solo son válidos cuando no hay interacción y el diseño es equilibrado
\end{enumerate}

\noindent\textbf{\color{azul}Explicación:} Cuando existe una interacción significativa, el efecto de cada factor no es constante: varía según el nivel del otro factor. Por tanto, el efecto principal (promedio marginal) puede ocultar o distorsionar la realidad. La norma en diseños factoriales es interpretar primero la interacción; si es significativa, los efectos principales deben interpretarse condicionalmente (efectos simples).

\noindent\textbf{\color{gris}Por qué son incorrectas las otras opciones:}
\begin{itemize}
  \item {\small A — Incorrecto: en un diseño factorial, la interacción es con frecuencia el hallazgo más valioso e informativo.}
  \item {\small C — Incorrecto: el criterio no es el tamaño del efecto de la interacción sino si es estadísticamente significativa.}
  \item {\small D — Incorrecto: los efectos principales son válidos en diseños equilibrados tanto si hay interacción como si no; el problema es su interpretación aislada cuando existe interacción.}
\end{itemize}
\vspace{8pt}\hrule\vspace{10pt}

\begin{mdframed}[backgroundcolor=fondogris,linecolor=ambar,linewidth=1pt,innerleftmargin=10pt,innerrightmargin=10pt,innertopmargin=7pt,innerbottommargin=7pt]
\textbf{Ítem 20 / n13} \hfill {\small\bfseries\color{ambar} Nivel 2 · Intermedio}
\end{mdframed}

En un diseño factorial 2×2 sobre el efecto del sueño (bien descansado vs. privado de sueño) y el tipo de tarea (fácil vs. difícil) sobre el rendimiento, los resultados muestran: bien descansado + tarea fácil: 9/10; bien descansado + tarea difícil: 8/10; privado de sueño + tarea fácil: 8/10; privado de sueño + tarea difícil: 4/10. ¿Existe interacción?

\begin{enumerate}
  \item[A.] No, porque en todas las condiciones el grupo bien descansado rinde mejor
  \item[\textbf{\color{verde}B.}] \textbf{\color{verde} Sí, porque el efecto de la privación del sueño sobre el rendimiento es mucho mayor en la tarea difícil que en la fácil} \hfill $\checkmark$
  \item[C.] No, porque las diferencias entre condiciones son pequeñas y no significativas
  \item[D.] Sí, porque hay cuatro condiciones y en los diseños 2×2 siempre hay interacción
\end{enumerate}

\noindent\textbf{\color{azul}Explicación:} Hay interacción cuando el efecto de un factor depende del nivel del otro. Aquí la privación de sueño reduce el rendimiento 1 punto en la tarea fácil (9→8) pero 4 puntos en la difícil (8→4). Este patrón diferencial es la interacción: la privación de sueño perjudica mucho más cuando la tarea es difícil.

\noindent\textbf{\color{gris}Por qué son incorrectas las otras opciones:}
\begin{itemize}
  \item {\small A — Incorrecto: aunque el grupo bien descansado rinde mejor en general, lo importante es si ese efecto es constante en ambos tipos de tarea; no lo es, lo que indica interacción.}
  \item {\small C — Incorrecto: la pregunta es conceptual, no estadística; el patrón de los datos (diferencias asimétricas) indica interacción con independencia de la significación.}
  \item {\small D — Incorrecto: los diseños 2×2 no siempre tienen interacción; la interacción depende de los datos, no de la estructura del diseño.}
\end{itemize}
\vspace{8pt}\hrule\vspace{10pt}

\begin{mdframed}[backgroundcolor=fondogris,linecolor=ambar,linewidth=1pt,innerleftmargin=10pt,innerrightmargin=10pt,innertopmargin=7pt,innerbottommargin=7pt]
\textbf{Ítem 21 / fa09} \hfill {\small\bfseries\color{ambar} Nivel 2 · Intermedio}
\end{mdframed}

En un diseño factorial 2×2 completamente aleatorizado (Factor A: medicación sí/no; Factor B: psicoterapia sí/no), los resultados muestran un efecto principal significativo de A (F=12.3, p=.001) y de B (F=9.8, p=.002), pero la interacción A×B no es significativa (F=0.4, p=.53). ¿Cómo se interpreta correctamente este patrón?

\begin{enumerate}
  \item[A.] Los efectos de A y B son espurios porque la interacción no es significativa
  \item[\textbf{\color{verde}B.}] \textbf{\color{verde} Tanto la medicación como la psicoterapia mejoran el resultado de forma independiente y aditiva; su efecto combinado es la suma de sus efectos por separado} \hfill $\checkmark$
  \item[C.] La medicación y la psicoterapia solo son eficaces cuando se administran juntas
  \item[D.] No pueden interpretarse los efectos principales sin una interacción significativa
\end{enumerate}

\noindent\textbf{\color{azul}Explicación:} Cuando la interacción no es significativa, los efectos principales pueden interpretarse de forma directa e independiente. Cada factor tiene un efecto que no depende del nivel del otro factor: la medicación mejora el resultado tanto con psicoterapia como sin ella, y la psicoterapia mejora el resultado tanto con medicación como sin ella. Los efectos son aditivos.

\noindent\textbf{\color{gris}Por qué son incorrectas las otras opciones:}
\begin{itemize}
  \item {\small A — Incorrecto: la ausencia de interacción no invalida los efectos principales; al contrario, cuando no hay interacción los efectos principales son más limpios e interpretables.}
  \item {\small C — Incorrecto: si la combinación fuera la única condición eficaz, los efectos principales serían no significativos y habría interacción; aquí los efectos principales son independientes.}
  \item {\small D — Incorrecto: la norma es la contraria: cuando hay interacción significativa, los efectos principales deben interpretarse con cautela; cuando no la hay, se interpretan directamente.}
\end{itemize}
\vspace{8pt}\hrule\vspace{10pt}

\begin{mdframed}[backgroundcolor=fondogris,linecolor=ambar,linewidth=1pt,innerleftmargin=10pt,innerrightmargin=10pt,innertopmargin=7pt,innerbottommargin=7pt]
\textbf{Ítem 22 / fa10} \hfill {\small\bfseries\color{ambar} Nivel 2 · Intermedio}
\end{mdframed}

Una investigadora diseña un estudio factorial 2×2 mixto: el Factor A (tipo de entrenamiento: intensivo vs. moderado) es entre sujetos y el Factor B (momento de medida: antes vs. después del entrenamiento) es intrasujeto. ¿Qué información específica proporciona la interacción A×B en este diseño?

\begin{enumerate}
  \item[A.] Si los participantes que reciben entrenamiento intensivo son más inteligentes que los de entrenamiento moderado
  \item[\textbf{\color{verde}B.}] \textbf{\color{verde} Si el cambio producido por el entrenamiento (diferencia antes-después) es diferente según el tipo de entrenamiento recibido} \hfill $\checkmark$
  \item[C.] Si la medida de antes es fiable comparada con la medida de después
  \item[D.] Si los dos tipos de entrenamiento producen los mismos resultados absolutos en el postest
\end{enumerate}

\noindent\textbf{\color{azul}Explicación:} En un diseño mixto con un factor temporal intrasujeto (antes vs. después), la interacción entre el tipo de entrenamiento y el momento de medida indica si el cambio a lo largo del tiempo (la mejora producida) es diferente según el grupo. Una interacción significativa significa que el entrenamiento intensivo produce una ganancia mayor (o menor) que el moderado. Esta es la pregunta central en la mayoría de los estudios de intervención.

\noindent\textbf{\color{gris}Por qué son incorrectas las otras opciones:}
\begin{itemize}
  \item {\small A — Incorrecto: las diferencias previas en inteligencia entre grupos serían una amenaza al diseño, no lo que mide la interacción; además se controlan con la medida pretest.}
  \item {\small C — Incorrecto: la fiabilidad del instrumento de medida es una cuestión psicométrica, no lo que revela la interacción en un diseño factorial.}
  \item {\small D — Incorrecto: si los postest fueran iguales pero los pretest diferentes, habría interacción pero ningún efecto general; la interacción no compara puntuaciones absolutas sino patrones de cambio.}
\end{itemize}
\vspace{8pt}\hrule\vspace{10pt}

\begin{mdframed}[backgroundcolor=fondogris,linecolor=rojo,linewidth=1pt,innerleftmargin=10pt,innerrightmargin=10pt,innertopmargin=7pt,innerbottommargin=7pt]
\textbf{Ítem 23 / f5} \hfill {\small\bfseries\color{rojo} Nivel 3 · Avanzado}
\end{mdframed}

Una investigadora diseña un estudio factorial 2×2 en el que el Factor A (entrenamiento: sí vs. no) es entre sujetos y el Factor B (momento de medida: pretest vs. postest) es intrasujeto. Este diseño se denomina:

\begin{enumerate}
  \item[A.] Diseño factorial completamente aleatorizado, porque hay asignación aleatoria al Factor A
  \item[B.] Diseño de medidas repetidas completo, porque hay una medida temporal repetida
  \item[\textbf{\color{verde}C.}] \textbf{\color{verde} Diseño mixto o split-plot, porque combina un factor entre sujetos y un factor intrasujeto} \hfill $\checkmark$
  \item[D.] Diseño cuasiexperimental factorial, porque el pretest es anterior al postest
\end{enumerate}

\noindent\textbf{\color{azul}Explicación:} El diseño mixto (también llamado split-plot) combina al menos un factor entre sujetos (diferentes participantes en cada nivel) con al menos un factor intrasujeto (los mismos participantes en todos los niveles). Es muy común en investigación aplicada porque el tiempo suele ser un factor intrasujeto y el grupo experimental/control un factor entre sujetos.

\noindent\textbf{\color{gris}Por qué son incorrectas las otras opciones:}
\begin{itemize}
  \item {\small A — Incorrecto: el diseño completamente aleatorizado implica que todos los factores son entre sujetos; aquí el Factor B (tiempo) es intrasujeto.}
  \item {\small B — Incorrecto: el diseño de medidas repetidas completo implica que todos los factores son intrasujeto; aquí el Factor A (entrenamiento) es entre sujetos.}
  \item {\small D — Incorrecto: el diseño mixto puede incluir asignación aleatoria (como aquí) y no es cuasiexperimental; el cuasiexperimental se define por la ausencia de aleatorización.}
\end{itemize}
\vspace{8pt}\hrule\vspace{10pt}

\begin{mdframed}[backgroundcolor=fondogris,linecolor=rojo,linewidth=1pt,innerleftmargin=10pt,innerrightmargin=10pt,innertopmargin=7pt,innerbottommargin=7pt]
\textbf{Ítem 24 / f6} \hfill {\small\bfseries\color{rojo} Nivel 3 · Avanzado}
\end{mdframed}

En un diseño factorial 2×2 completamente aleatorizado se detecta una interacción significativa (F(1,76)=12.4, p=.001, η²=.14). ¿Qué análisis debe realizarse a continuación para interpretar correctamente la interacción?

\begin{enumerate}
  \item[A.] Calcular el tamaño del efecto de cada factor por separado usando d de Cohen
  \item[\textbf{\color{verde}B.}] \textbf{\color{verde} Análisis de efectos simples: comparar los niveles de un factor dentro de cada nivel del otro factor} \hfill $\checkmark$
  \item[C.] Repetir el ANOVA excluyendo los participantes con puntuaciones extremas para limpiar la interacción
  \item[D.] Calcular la correlación entre los dos factores para cuantificar su interdependencia
\end{enumerate}

\noindent\textbf{\color{azul}Explicación:} Los efectos simples descomponen la interacción examinando el efecto de un factor para cada nivel del otro factor. Por ejemplo, si los factores son A y B, se calcula el efecto de A solo para B=1 y luego solo para B=2. Esto revela en qué condiciones específicas se produce la diferencia y permite una interpretación sustantiva de la interacción.

\noindent\textbf{\color{gris}Por qué son incorrectas las otras opciones:}
\begin{itemize}
  \item {\small A — Incorrecto: el eta-cuadrado cuantifica el tamaño del efecto global de cada factor, pero no descompone la interacción en sus componentes interpretables.}
  \item {\small C — Incorrecto: excluir participantes con puntuaciones extremas es una manipulación de datos que sesgaría los resultados; la interacción no se "limpia" de ese modo.}
  \item {\small D — Incorrecto: la correlación entre los factores no tiene sentido cuando ambos son variables independientes categóricas manipuladas por el investigador.}
\end{itemize}
\vspace{8pt}\hrule\vspace{10pt}

\begin{mdframed}[backgroundcolor=fondogris,linecolor=rojo,linewidth=1pt,innerleftmargin=10pt,innerrightmargin=10pt,innertopmargin=7pt,innerbottommargin=7pt]
\textbf{Ítem 25 / fa11} \hfill {\small\bfseries\color{rojo} Nivel 3 · Avanzado}
\end{mdframed}

Una investigadora realiza un diseño factorial 2×2 y obtiene los siguientes resultados: condición A1B1: M=6.0; A1B2: M=6.2; A2B1: M=6.1; A2B2: M=6.0. Las cuatro medias son prácticamente iguales. ¿Qué se puede concluir sobre los efectos en este diseño?

\begin{enumerate}
  \item[A.] Hay una interacción disordinal (cruzada) porque las medias cambian de orden entre condiciones
  \item[\textbf{\color{verde}B.}] \textbf{\color{verde} No hay efectos principales ni interacción apreciables; el factor A, el factor B y su combinación no parecen afectar a la variable dependiente} \hfill $\checkmark$
  \item[C.] Hay un efecto principal claro de B porque las medias de B1 y B2 son diferentes según el factor A
  \item[D.] Hay interacción porque la diferencia A1B1-A1B2 no es igual que la diferencia A2B1-A2B2
\end{enumerate}

\noindent\textbf{\color{azul}Explicación:} Cuando las cuatro medias son prácticamente iguales, las diferencias marginales de A (promedio de A1 vs. promedio de A2) y de B son cercanas a cero, y las diferencias entre celdas también son mínimas. Esto indica la ausencia de efectos principales y de interacción: ni A ni B ni su combinación parecen influir sobre la VD. Esta conclusión siempre debe respaldarse con las pruebas estadísticas formales.

\noindent\textbf{\color{gris}Por qué son incorrectas las otras opciones:}
\begin{itemize}
  \item {\small A — Incorrecto: la interacción disordinal implica que las líneas de perfil se cruzan; con cuatro medias tan similares no hay cruce real.}
  \item {\small C — Incorrecto: para que haya efecto principal de B, las medias marginales de B1 y B2 deben diferir; con medias casi idénticas eso no ocurre.}
  \item {\small D — Incorrecto: aunque las diferencias A1B1-A1B2 y A2B1-A2B2 no son exactamente iguales, cuando las diferencias son mínimas (< 0.2 puntos) no hay base para hablar de interacción sustantiva.}
\end{itemize}
\vspace{8pt}\hrule\vspace{10pt}

\section{Cuasiexperimental (CUASI) · T.10}
{\small\color{gris} 13 ítems · Niveles: L1=6 · L2=4 · L3=3}

\begin{mdframed}[backgroundcolor=fondogris,linecolor=azul,linewidth=1pt,innerleftmargin=10pt,innerrightmargin=10pt,innertopmargin=7pt,innerbottommargin=7pt]
\textbf{Ítem 26 / c1} \hfill {\small\bfseries\color{azul} Nivel 1 · Básico}
\end{mdframed}

Un psicólogo evalúa un programa de convivencia en un instituto. Aplica el programa al grupo A (1.º de ESO) y mantiene al grupo B (2.º de ESO) sin intervención. Mide el clima de aula antes y después en ambos grupos. No puede reasignar a los estudiantes porque los grupos ya están formados por el centro. Según León y Montero, ¿qué define este diseño como cuasiexperimental y no como un experimento verdadero?

\begin{enumerate}
  \item[A.] Que mide el clima de aula antes y después de la intervención
  \item[B.] Que el investigador no puede manipular ninguna variable independiente
  \item[\textbf{\color{verde}C.}] \textbf{\color{verde} La ausencia de asignación aleatoria de participantes a las condiciones, usando grupos preexistentes} \hfill $\checkmark$
  \item[D.] Que el grupo B no recibe ningún tipo de intervención activa
\end{enumerate}

\noindent\textbf{\color{azul}Explicación:} El criterio definitorio del diseño cuasiexperimental según León y Montero es la ausencia de asignación aleatoria. Hay manipulación de la VI (el programa se aplica deliberadamente a un grupo), pero los participantes no se asignan al azar a las condiciones porque los grupos ya existen. Esto limita la validez interna porque los grupos pueden diferir en variables extrañas.

\noindent\textbf{\color{gris}Por qué son incorrectas las otras opciones:}
\begin{itemize}
  \item {\small A — Incorrecto: medir antes y después (pretest-postest) puede hacerse tanto en experimentos verdaderos como en cuasiexperimentos; no es lo que define al diseño como cuasiexperimental.}
  \item {\small B — Incorrecto: en el cuasiexperimental sí hay manipulación de la VI; es precisamente lo que lo distingue del ex post facto.}
  \item {\small D — Incorrecto: tener un grupo de comparación que no recibe tratamiento es conveniente pero no define al diseño como cuasiexperimental.}
\end{itemize}
\vspace{8pt}\hrule\vspace{10pt}

\begin{mdframed}[backgroundcolor=fondogris,linecolor=azul,linewidth=1pt,innerleftmargin=10pt,innerrightmargin=10pt,innertopmargin=7pt,innerbottommargin=7pt]
\textbf{Ítem 27 / n04} \hfill {\small\bfseries\color{azul} Nivel 1 · Básico}
\end{mdframed}

Un director de colegio aplica un programa de lectura rápida a todos los alumnos de 5.º A. Los alumnos de 5.º B no reciben el programa y sirven como comparación. El director mide la velocidad lectora antes y después en ambas clases. No puede mezclar a los alumnos entre grupos. ¿Qué tipo de diseño es este?

\begin{enumerate}
  \item[A.] Diseño experimental entre grupos, porque hay un grupo que recibe el programa y otro que no
  \item[\textbf{\color{verde}B.}] \textbf{\color{verde} Diseño cuasiexperimental de grupos no equivalentes, porque hay intervención pero sin asignación aleatoria} \hfill $\checkmark$
  \item[C.] Diseño ex post facto, porque el programa ya existía en el colegio antes del estudio
  \item[D.] Diseño descriptivo longitudinal, porque mide en dos momentos distintos
\end{enumerate}

\noindent\textbf{\color{azul}Explicación:} El diseño cuasiexperimental se caracteriza por la presencia de una intervención (el programa) pero sin asignación aleatoria de los participantes a las condiciones: los grupos ya estaban formados (las clases). Esto limita la validez interna respecto al experimento verdadero.

\noindent\textbf{\color{gris}Por qué son incorrectas las otras opciones:}
\begin{itemize}
  \item {\small A — Incorrecto: en el experimento entre grupos los participantes se asignan al azar a las condiciones; aquí las clases son grupos preexistentes.}
  \item {\small C — Incorrecto: en el ex post facto no hay intervención activa del investigador; aquí el director aplica deliberadamente el programa.}
  \item {\small D — Incorrecto: el diseño longitudinal descriptivo solo describe cambios en el tiempo sin aplicar ninguna intervención.}
\end{itemize}
\vspace{8pt}\hrule\vspace{10pt}

\begin{mdframed}[backgroundcolor=fondogris,linecolor=azul,linewidth=1pt,innerleftmargin=10pt,innerrightmargin=10pt,innertopmargin=7pt,innerbottommargin=7pt]
\textbf{Ítem 28 / v02} \hfill {\small\bfseries\color{azul} Nivel 1 · Básico}
\end{mdframed}

Un investigador realiza un experimento muy controlado en el laboratorio con condiciones artificiales muy distintas al entorno real. Obtiene resultados muy claros. Sin embargo, cuando intenta aplicar sus conclusiones a situaciones del mundo real, los resultados no se replican. ¿Qué tipo de validez falla principalmente?

\begin{enumerate}
  \item[A.] Validez interna, porque los grupos no eran equivalentes al inicio del estudio
  \item[\textbf{\color{verde}B.}] \textbf{\color{verde} Validez externa o ecológica, porque las condiciones artificiales del laboratorio no se trasladan a la vida real} \hfill $\checkmark$
  \item[C.] Validez de constructo, porque las variables no estaban bien definidas conceptualmente
  \item[D.] Fiabilidad, porque los resultados no son consistentes entre distintos estudios
\end{enumerate}

\noindent\textbf{\color{azul}Explicación:} La validez ecológica es un subtipo de validez externa que indica en qué medida los resultados de un estudio se pueden generalizar a situaciones reales. Los experimentos de laboratorio muy artificiales pueden tener alta validez interna pero baja validez ecológica: lo que ocurre en el laboratorio puede no reflejar lo que ocurre en la vida cotidiana.

\noindent\textbf{\color{gris}Por qué son incorrectas las otras opciones:}
\begin{itemize}
  \item {\small A — Incorrecto: la validez interna se refiere al control de variables extrañas dentro del estudio; aquí el problema es la generalización al mundo real, no el control interno.}
  \item {\small C — Incorrecto: la validez de constructo se refiere a si los instrumentos miden realmente el constructo que dicen medir; ese no es el problema descrito.}
  \item {\small D — Incorrecto: la falta de replicación puede deberse a validez externa baja, no necesariamente a falta de fiabilidad del instrumento.}
\end{itemize}
\vspace{8pt}\hrule\vspace{10pt}

\begin{mdframed}[backgroundcolor=fondogris,linecolor=azul,linewidth=1pt,innerleftmargin=10pt,innerrightmargin=10pt,innertopmargin=7pt,innerbottommargin=7pt]
\textbf{Ítem 29 / n11} \hfill {\small\bfseries\color{azul} Nivel 1 · Básico}
\end{mdframed}

Un investigador evalúa si un nuevo método de enseñanza del inglés mejora la pronunciación. Lo aplica en dos colegios que ya eligieron usar este método, y compara con dos colegios que siguen con el método tradicional. Mide la pronunciación antes y después en todos. Los colegios no fueron asignados aleatoriamente. ¿Qué diseño es este?

\begin{enumerate}
  \item[A.] Diseño experimental con cuatro grupos
  \item[\textbf{\color{verde}B.}] \textbf{\color{verde} Diseño cuasiexperimental de grupos no equivalentes con pretest y postest} \hfill $\checkmark$
  \item[C.] Diseño ex post facto prospectivo, porque sigue a los colegios en el tiempo
  \item[D.] Diseño descriptivo longitudinal con dos tipos de colegios
\end{enumerate}

\noindent\textbf{\color{azul}Explicación:} El diseño cuasiexperimental de grupos no equivalentes con pretest-postest es el más común en evaluación educativa: hay intervención activa (el nuevo método), medida antes y después, y grupo de comparación, pero los grupos (colegios) ya existían antes y no se asignaron al azar.

\noindent\textbf{\color{gris}Por qué son incorrectas las otras opciones:}
\begin{itemize}
  \item {\small A — Incorrecto: para ser experimental los colegios debería haber sido asignados aleatoriamente al nuevo método; aquí los colegios eligieron por sí mismos.}
  \item {\small C — Incorrecto: el ex post facto prospectivo parte de una causa preexistente en los participantes y no implica que el investigador aplique una intervención.}
  \item {\small D — Incorrecto: el diseño longitudinal descriptivo no incluye intervención activa; aquí el investigador aplica (o permite) deliberadamente el nuevo método.}
\end{itemize}
\vspace{8pt}\hrule\vspace{10pt}

\begin{mdframed}[backgroundcolor=fondogris,linecolor=azul,linewidth=1pt,innerleftmargin=10pt,innerrightmargin=10pt,innertopmargin=7pt,innerbottommargin=7pt]
\textbf{Ítem 30 / cu06} \hfill {\small\bfseries\color{azul} Nivel 1 · Básico}
\end{mdframed}

Una investigadora aplica un programa de mindfulness a los empleados de una empresa (grupo A) y compara sus niveles de estrés con los empleados de otra empresa similar que no recibe el programa (grupo B). Ambos grupos se miden antes y después. Los empleados no fueron asignados aleatoriamente entre empresas. ¿Qué tipo de diseño es este?

\begin{enumerate}
  \item[A.] Diseño experimental entre grupos, porque hay una empresa que recibe el programa y otra que no
  \item[\textbf{\color{verde}B.}] \textbf{\color{verde} Diseño cuasiexperimental de grupos no equivalentes con pretest-postest, porque hay intervención pero sin aleatorización de los participantes} \hfill $\checkmark$
  \item[C.] Diseño ex post facto prospectivo, porque se sigue a los empleados a lo largo del tiempo
  \item[D.] Diseño descriptivo longitudinal, porque mide el estrés en dos momentos distintos
\end{enumerate}

\noindent\textbf{\color{azul}Explicación:} El diseño cuasiexperimental de grupos no equivalentes con pretest-postest es característico de la evaluación de intervenciones en entornos organizacionales o educativos: hay una intervención activa (el programa de mindfulness), hay medida antes y después, hay grupo de comparación, pero los participantes pertenecen a grupos naturales preexistentes (las empresas) y no se asignaron al azar.

\noindent\textbf{\color{gris}Por qué son incorrectas las otras opciones:}
\begin{itemize}
  \item {\small A — Incorrecto: el experimento verdadero requiere asignación aleatoria de los participantes a las condiciones; aquí los grupos son las empresas completas, que no se asignaron al azar.}
  \item {\small C — Incorrecto: el ex post facto prospectivo parte de una condición preexistente en los participantes y no implica que el investigador aplique ninguna intervención activa.}
  \item {\small D — Incorrecto: el diseño longitudinal descriptivo solo mide y describe cambios sin aplicar ninguna intervención; aquí hay una intervención deliberada.}
\end{itemize}
\vspace{8pt}\hrule\vspace{10pt}

\begin{mdframed}[backgroundcolor=fondogris,linecolor=azul,linewidth=1pt,innerleftmargin=10pt,innerrightmargin=10pt,innertopmargin=7pt,innerbottommargin=7pt]
\textbf{Ítem 31 / cu07} \hfill {\small\bfseries\color{azul} Nivel 1 · Básico}
\end{mdframed}

Un investigador quiere evaluar si una campaña publicitaria sobre nutrición saludable redujo el consumo de alimentos ultraprocesados. No puede asignar personas al azar. Recoge datos de ventas mensuales de ultraprocesados en la misma región durante los 18 meses anteriores y los 18 meses posteriores al lanzamiento de la campaña. ¿Qué diseño cuasiexperimental está usando?

\begin{enumerate}
  \item[A.] Diseño de grupos no equivalentes, porque compara dos regiones distintas
  \item[\textbf{\color{verde}B.}] \textbf{\color{verde} Diseño de series temporales interrumpidas, porque tiene muchas medidas antes y después de una intervención en un mismo grupo} \hfill $\checkmark$
  \item[C.] Diseño de caso único A-B-A, porque alterna fases con y sin la campaña
  \item[D.] Diseño ex post facto retrospectivo, porque los datos de ventas ya existían antes del estudio
\end{enumerate}

\noindent\textbf{\color{azul}Explicación:} El diseño de series temporales interrumpidas usa múltiples medidas en el mismo grupo antes y después de una intervención (la campaña). La serie de medidas previas permite establecer la tendencia de base, y la serie posterior muestra si hay un cambio en el nivel o la pendiente de esa tendencia atribuible a la campaña. Es el diseño cuasiexperimental más usado en políticas públicas.

\noindent\textbf{\color{gris}Por qué son incorrectas las otras opciones:}
\begin{itemize}
  \item {\small A — Incorrecto: el diseño de grupos no equivalentes compara dos grupos distintos; aquí solo hay un grupo (la misma región) medido repetidamente.}
  \item {\small C — Incorrecto: el diseño A-B-A de caso único requiere retirar la intervención para restablecer la línea de base; las campañas publicitarias no pueden "retirarse" del conocimiento colectivo.}
  \item {\small D — Incorrecto: en el ex post facto no hay intervención activa del investigador; aquí la campaña es una intervención deliberada cuyo efecto se evalúa.}
\end{itemize}
\vspace{8pt}\hrule\vspace{10pt}

\begin{mdframed}[backgroundcolor=fondogris,linecolor=ambar,linewidth=1pt,innerleftmargin=10pt,innerrightmargin=10pt,innertopmargin=7pt,innerbottommargin=7pt]
\textbf{Ítem 32 / c2} \hfill {\small\bfseries\color{ambar} Nivel 2 · Intermedio}
\end{mdframed}

Un equipo evalúa una intervención para reducir el absentismo escolar en tres colegios. Dos colegios reciben la intervención y uno actúa como grupo de comparación. Se miden las tasas de absentismo en los 12 meses anteriores y en los 12 meses posteriores a la intervención. ¿Cuál es la principal amenaza a la validez interna de este diseño cuasiexperimental?

\begin{enumerate}
  \item[A.] La mortalidad experimental, porque algunos alumnos pueden cambiar de colegio durante el estudio
  \item[\textbf{\color{verde}B.}] \textbf{\color{verde} La no equivalencia inicial de los grupos, ya que los colegios no se asignaron aleatoriamente y pueden diferir en factores socioeconómicos, culturales o de organización} \hfill $\checkmark$
  \item[C.] La instrumentación, porque el registro de absentismo puede variar entre medidas
  \item[D.] La regresión a la media, porque los colegios con absentismo alto tienden a mejorar espontáneamente
\end{enumerate}

\noindent\textbf{\color{azul}Explicación:} La principal amenaza en los diseños cuasiexperimentales con grupos no equivalentes es precisamente la falta de equivalencia inicial: los grupos pueden diferir en variables extrañas antes de la intervención (nivel socioeconómico, motivación institucional, calidad docente). Sin aleatorización, no puede descartarse que las diferencias postest se deban a esas variables previas y no a la intervención.

\noindent\textbf{\color{gris}Por qué son incorrectas las otras opciones:}
\begin{itemize}
  \item {\small A — Incorrecto: la mortalidad experimental es una amenaza real pero secundaria; puede controlarse con análisis de sensibilidad.}
  \item {\small C — Incorrecto: la instrumentación es una amenaza si el instrumento de medida cambia entre momentos; los registros de absentismo suelen ser consistentes.}
  \item {\small D — Incorrecto: la regresión a la media afecta cuando se seleccionan grupos por puntuaciones extremas; aquí los colegios no se seleccionaron por tener el absentismo más alto.}
\end{itemize}
\vspace{8pt}\hrule\vspace{10pt}

\begin{mdframed}[backgroundcolor=fondogris,linecolor=ambar,linewidth=1pt,innerleftmargin=10pt,innerrightmargin=10pt,innertopmargin=7pt,innerbottommargin=7pt]
\textbf{Ítem 33 / c3} \hfill {\small\bfseries\color{ambar} Nivel 2 · Intermedio}
\end{mdframed}

Un investigador usa un diseño de series temporales interrumpidas para evaluar el impacto de una nueva ley de tráfico sobre los accidentes. Recoge datos de accidentalidad mensualmente durante 24 meses antes y 24 meses después de la entrada en vigor de la ley. ¿Cuál es la ventaja principal de este diseño respecto al simple pretest-postest sin grupo control?

\begin{enumerate}
  \item[A.] Permite aleatorizar la aplicación de la ley entre regiones para controlar variables extrañas
  \item[\textbf{\color{verde}B.}] \textbf{\color{verde} El patrón de múltiples medidas antes y después permite distinguir el efecto de la ley de tendencias previas y fluctuaciones estacionales} \hfill $\checkmark$
  \item[C.] Elimina completamente la amenaza de historia porque hay muchas medidas antes de la intervención
  \item[D.] Convierte el diseño en experimental porque hay una interrupción deliberada de la serie temporal
\end{enumerate}

\noindent\textbf{\color{azul}Explicación:} El diseño de series temporales interrumpidas usa múltiples medidas antes y después del punto de intervención. Esto permite estimar la tendencia previa y compararla con la tendencia posterior, distinguiendo el efecto de la intervención de cambios graduales preexistentes, efectos estacionales o fluctuaciones aleatorias. Sigue sin controlar completamente la amenaza de historia.

\noindent\textbf{\color{gris}Por qué son incorrectas las otras opciones:}
\begin{itemize}
  \item {\small A — Incorrecto: la ley se aplica a toda la población simultáneamente; no es posible aleatorizar su aplicación geográfica en este caso.}
  \item {\small C — Incorrecto: la amenaza de historia (eventos externos coincidentes con la ley) sigue siendo una amenaza aunque haya muchas medidas previas.}
  \item {\small D — Incorrecto: la interrupción de la serie es la intervención, pero no hay asignación aleatoria ni control experimental; el diseño sigue siendo cuasiexperimental.}
\end{itemize}
\vspace{8pt}\hrule\vspace{10pt}

\begin{mdframed}[backgroundcolor=fondogris,linecolor=ambar,linewidth=1pt,innerleftmargin=10pt,innerrightmargin=10pt,innertopmargin=7pt,innerbottommargin=7pt]
\textbf{Ítem 34 / cu08} \hfill {\small\bfseries\color{ambar} Nivel 2 · Intermedio}
\end{mdframed}

En un diseño cuasiexperimental de grupos no equivalentes, el pretest muestra que el grupo experimental parte con una media significativamente más alta que el grupo control (M=72 vs. M=60). El investigador quiere comparar los postest pero su director le dice que primero debe comprobar la "equivalencia inicial". ¿Por qué es tan importante esta comprobación en los diseños cuasiexperimentales?

\begin{enumerate}
  \item[A.] Porque sin equivalencia inicial el investigador no puede calcular la diferencia entre pretest y postest
  \item[\textbf{\color{verde}B.}] \textbf{\color{verde} Porque si los grupos parten de niveles distintos, las diferencias en el postest pueden deberse a esas diferencias previas y no a la intervención, lo que amenaza la validez interna} \hfill $\checkmark$
  \item[C.] Porque la equivalencia inicial es un requisito legal en los estudios con seres humanos
  \item[D.] Porque solo si los grupos son equivalentes al inicio pueden asignarse aleatoriamente al postest
\end{enumerate}

\noindent\textbf{\color{azul}Explicación:} La equivalencia inicial es la principal amenaza a la validez interna en los diseños cuasiexperimentales. Si los grupos parten de niveles diferentes, cualquier diferencia postest puede reflejar las diferencias previas y no el efecto de la intervención. Por eso es esencial documentar las diferencias iniciales y, si existen, usar técnicas de ajuste estadístico como el ANCOVA.

\noindent\textbf{\color{gris}Por qué son incorrectas las otras opciones:}
\begin{itemize}
  \item {\small A — Incorrecto: el cálculo de la diferencia pretest-postest es aritméticamente posible independientemente de la equivalencia; el problema es su interpretación causal.}
  \item {\small C — Incorrecto: la equivalencia inicial no es un requisito legal sino una condición metodológica necesaria para la validez interna de la inferencia.}
  \item {\small D — Incorrecto: la asignación aleatoria ocurre antes del estudio, no tras el pretest; y precisamente la falta de aleatorización previa es lo que genera el problema de no equivalencia.}
\end{itemize}
\vspace{8pt}\hrule\vspace{10pt}

\begin{mdframed}[backgroundcolor=fondogris,linecolor=ambar,linewidth=1pt,innerleftmargin=10pt,innerrightmargin=10pt,innertopmargin=7pt,innerbottommargin=7pt]
\textbf{Ítem 35 / cu09} \hfill {\small\bfseries\color{ambar} Nivel 2 · Intermedio}
\end{mdframed}

Un equipo evalúa un programa de inserción laboral para personas desempleadas de larga duración. Aplican el programa a todos los que se inscriben voluntariamente y comparan su tasa de empleo a los 6 meses con la de personas desempleadas que no se inscribieron. ¿Cuál es el principal problema de selección en este diseño y cómo afecta a la validez interna?

\begin{enumerate}
  \item[A.] La autoselección: las personas que se inscriben voluntariamente pueden tener más motivación, recursos o habilidades que las que no se inscriben, lo que amenaza gravemente la validez interna
  \item[\textbf{\color{verde}B.}] \textbf{\color{verde} La mortalidad experimental: algunos participantes abandonarán el programa antes de los 6 meses} \hfill $\checkmark$
  \item[C.] El tamaño muestral: si pocos se inscriben, el estudio tendrá poca potencia estadística
  \item[D.] El efecto de novedad: los participantes pueden esforzarse más solo por el hecho de participar en un programa nuevo
\end{enumerate}

\noindent\textbf{\color{azul}Explicación:} La autoselección es una forma grave de sesgo de selección en los estudios con participación voluntaria: quienes eligen participar pueden diferir sistemáticamente de quienes no lo hacen en variables relevantes (motivación, capital social, actitud proactiva). Esto hace que cualquier diferencia en el resultado pueda atribuirse a las características previas de los participantes y no al programa.

\noindent\textbf{\color{gris}Por qué son incorrectas las otras opciones:}
\begin{itemize}
  \item {\small B — Incorrecto: la mortalidad experimental es una amenaza real pero secundaria; el problema principal es que los grupos no eran comparables desde el inicio.}
  \item {\small C — Incorrecto: el tamaño muestral afecta a la potencia estadística, no a la validez interna; un estudio puede ser inválido internamente aunque tenga mucha potencia.}
  \item {\small D — Incorrecto: el efecto de novedad (efecto Hawthorne) es una amenaza real pero no tan central como la autoselección en este caso concreto.}
\end{itemize}
\vspace{8pt}\hrule\vspace{10pt}

\begin{mdframed}[backgroundcolor=fondogris,linecolor=rojo,linewidth=1pt,innerleftmargin=10pt,innerrightmargin=10pt,innertopmargin=7pt,innerbottommargin=7pt]
\textbf{Ítem 36 / c4} \hfill {\small\bfseries\color{rojo} Nivel 3 · Avanzado}
\end{mdframed}

Un investigador evalúa un programa de reducción del estrés en empleados de dos empresas similares (una recibe el programa, otra no). Mide el estrés antes y después. Los datos pretest muestran que los grupos difieren significativamente (empresa intervenida: M=72; empresa control: M=61). ¿Qué estrategia analítica permite controlar mejor esta diferencia inicial para estimar el efecto del programa?

\begin{enumerate}
  \item[A.] Ignorar el pretest y comparar solo las puntuaciones postest entre grupos con una prueba t
  \item[\textbf{\color{verde}B.}] \textbf{\color{verde} Análisis de covarianza (ANCOVA) usando el pretest como covariable para ajustar las diferencias iniciales} \hfill $\checkmark$
  \item[C.] Recodificar la variable de estrés en rangos percentiles para equiparar ambos grupos
  \item[D.] Eliminar los participantes de la empresa intervenida con puntuaciones pretest superiores a 70
\end{enumerate}

\noindent\textbf{\color{azul}Explicación:} El ANCOVA usa el pretest como covariable: ajusta estadísticamente las puntuaciones postest eliminando la variabilidad atribuible a las diferencias preexistentes entre grupos. Esto no resuelve el problema de la no equivalencia (que requeriría aleatorización), pero ofrece una estimación del efecto de la intervención más ajustada que comparar los postest sin control.

\noindent\textbf{\color{gris}Por qué son incorrectas las otras opciones:}
\begin{itemize}
  \item {\small A — Incorrecto: comparar solo los postest sin ajustar las diferencias iniciales sesga la estimación del efecto: si el grupo intervenido partía de niveles más altos, su cambio puede parecer mayor o menor de lo que realmente es.}
  \item {\small C — Incorrecto: recodificar en percentiles no equipara los grupos; solo cambia la escala de medida sin resolver la no equivalencia inicial.}
  \item {\small D — Incorrecto: eliminar participantes en función del pretest distorsionaría la muestra y sesgaría aún más los resultados.}
\end{itemize}
\vspace{8pt}\hrule\vspace{10pt}

\begin{mdframed}[backgroundcolor=fondogris,linecolor=rojo,linewidth=1pt,innerleftmargin=10pt,innerrightmargin=10pt,innertopmargin=7pt,innerbottommargin=7pt]
\textbf{Ítem 37 / c5} \hfill {\small\bfseries\color{rojo} Nivel 3 · Avanzado}
\end{mdframed}

Un investigador quiere evaluar el efecto de una nueva metodología docente en cuatro colegios. Dos colegios ya usan la metodología nueva (grupo experimental) y dos siguen con la metodología tradicional (grupo control). No hay asignación aleatoria. Recopila datos de rendimiento de los últimos cinco años. ¿Cuál es la denominación específica de este diseño cuasiexperimental según León y Montero?

\begin{enumerate}
  \item[A.] Diseño de caso único con línea de base múltiple
  \item[\textbf{\color{verde}B.}] \textbf{\color{verde} Diseño cuasiexperimental de grupos no equivalentes con serie histórica (retrospectiva)} \hfill $\checkmark$
  \item[C.] Diseño ex post facto prospectivo con grupo de comparación
  \item[D.] Diseño experimental de medidas repetidas con bloqueo por colegio
\end{enumerate}

\noindent\textbf{\color{azul}Explicación:} El diseño cuasiexperimental de grupos no equivalentes con serie histórica combina la comparación entre grupos preexistentes (no equivalentes por falta de aleatorización) con múltiples medidas retrospectivas que permiten examinar si las tendencias previas diferían entre grupos antes de la intervención, fortaleciendo así la inferencia sobre el efecto de la metodología.

\noindent\textbf{\color{gris}Por qué son incorrectas las otras opciones:}
\begin{itemize}
  \item {\small A — Incorrecto: el diseño de caso único se aplica a un solo participante con replicación intrasujeto de la intervención a lo largo del tiempo.}
  \item {\small C — Incorrecto: el ex post facto prospectivo parte de una causa presente y sigue a los participantes hacia el futuro; aquí los datos son retrospectivos y hay manipulación (los colegios con nueva metodología la reciben deliberadamente).}
  \item {\small D — Incorrecto: no hay asignación aleatoria de colegios ni de alumnos; el bloqueo es una estrategia experimental que requiere aleatorización dentro de cada bloque.}
\end{itemize}
\vspace{8pt}\hrule\vspace{10pt}

\begin{mdframed}[backgroundcolor=fondogris,linecolor=rojo,linewidth=1pt,innerleftmargin=10pt,innerrightmargin=10pt,innertopmargin=7pt,innerbottommargin=7pt]
\textbf{Ítem 38 / cu10} \hfill {\small\bfseries\color{rojo} Nivel 3 · Avanzado}
\end{mdframed}

Un investigador evalúa el impacto de una nueva ley que prohíbe el tabaco en bares y restaurantes. Recoge datos de hospitalizaciones por infarto de miocardio mes a mes durante los 3 años anteriores y los 3 años posteriores a la ley en la región donde se aprobó. Un colega le señala que, aunque hay un descenso en los infartos, podría deberse a una campaña de salud cardiovascular nacional que comenzó al mismo tiempo que la ley. ¿Qué amenaza a la validez interna identifica el colega?

\begin{enumerate}
  \item[A.] Maduración, porque la población envejece y tiene menos infartos con el tiempo
  \item[\textbf{\color{verde}B.}] \textbf{\color{verde} Historia, porque un evento externo coincidente con la intervención puede explicar el cambio observado} \hfill $\checkmark$
  \item[C.] Sesgo de selección, porque los hospitales que reportan infartos pueden ser distintos antes y después de la ley
  \item[D.] Regresión a la media, porque los años con más infartos tienden a ser seguidos por años con menos
\end{enumerate}

\noindent\textbf{\color{azul}Explicación:} La amenaza de historia ocurre cuando un evento externo ajeno a la intervención estudiada ocurre al mismo tiempo que ella y puede ser el responsable del cambio observado. En el diseño de series temporales sin grupo de comparación, es imposible descartar que la mejora se deba a la campaña de salud nacional en lugar de a la ley antitabaco. La solución es añadir una región de comparación que no tuviera la ley.

\noindent\textbf{\color{gris}Por qué son incorrectas las otras opciones:}
\begin{itemize}
  \item {\small A — Incorrecto: la maduración son cambios graduales en los participantes o la población con el tiempo, no eventos externos discretos; además un efecto de maduración produciría un descenso progresivo, no un cambio brusco al inicio de la ley.}
  \item {\small C — Incorrecto: el sesgo de selección afecta a la comparabilidad inicial de los grupos; en un diseño de series temporales en una sola región no hay grupos distintos de personas.}
  \item {\small D — Incorrecto: la regresión a la media ocurre cuando se seleccionan grupos en momentos de puntuaciones extremas; aquí no se seleccionaron períodos extremos sino que se tomaron todos los meses consecutivos.}
\end{itemize}
\vspace{8pt}\hrule\vspace{10pt}

\section{Descriptivos (DES) · T.4--5}
{\small\color{gris} 14 ítems · Niveles: L1=8 · L2=4 · L3=2}

\begin{mdframed}[backgroundcolor=fondogris,linecolor=azul,linewidth=1pt,innerleftmargin=10pt,innerrightmargin=10pt,innertopmargin=7pt,innerbottommargin=7pt]
\textbf{Ítem 39 / d1} \hfill {\small\bfseries\color{azul} Nivel 1 · Básico}
\end{mdframed}

Una investigadora quiere conocer la prevalencia del síndrome de burnout en enfermeras de hospitales públicos de Madrid. Diseña un cuestionario validado y lo aplica a una muestra representativa de 400 enfermeras, sin manipular ninguna variable ni establecer grupos de comparación. ¿Qué tipo de diseño está empleando según León y Montero?

\begin{enumerate}
  \item[A.] Diseño cuasiexperimental de grupos no equivalentes
  \item[\textbf{\color{verde}B.}] \textbf{\color{verde} Diseño descriptivo por encuestas} \hfill $\checkmark$
  \item[C.] Diseño ex post facto prospectivo
  \item[D.] Diseño experimental entre grupos con grupo control
\end{enumerate}

\noindent\textbf{\color{azul}Explicación:} El diseño descriptivo por encuestas se caracteriza porque el investigador recoge información sobre una o más variables en una muestra sin manipular ninguna condición y sin establecer grupos con fines causales. El objetivo es describir la situación tal como existe.

\noindent\textbf{\color{gris}Por qué son incorrectas las otras opciones:}
\begin{itemize}
  \item {\small A — Incorrecto: el cuasiexperimental implica manipulación de la VI aunque sin aleatorización, lo que no ocurre aquí.}
  \item {\small C — Incorrecto: el ex post facto prospectivo parte de una causa ya presente y sigue a los participantes hacia el futuro para observar efectos.}
  \item {\small D — Incorrecto: el diseño experimental requiere manipulación deliberada de una VI y asignación aleatoria a condiciones.}
\end{itemize}
\vspace{8pt}\hrule\vspace{10pt}

\begin{mdframed}[backgroundcolor=fondogris,linecolor=azul,linewidth=1pt,innerleftmargin=10pt,innerrightmargin=10pt,innertopmargin=7pt,innerbottommargin=7pt]
\textbf{Ítem 40 / d2} \hfill {\small\bfseries\color{azul} Nivel 1 · Básico}
\end{mdframed}

Un psicólogo registra durante tres semanas las conductas de ayuda espontánea entre niños de una guardería durante el recreo. Utiliza un sistema de categorías predefinidas y no interactúa con los niños. Al finalizar calcula la frecuencia de cada tipo de conducta. ¿Qué diseño descriptivo describe mejor este estudio?

\begin{enumerate}
  \item[A.] Encuesta longitudinal con seguimiento
  \item[\textbf{\color{verde}B.}] \textbf{\color{verde} Observación sistemática no participante} \hfill $\checkmark$
  \item[C.] Diseño de caso único con línea de base
  \item[D.] Investigación-acción participativa
\end{enumerate}

\noindent\textbf{\color{azul}Explicación:} La observación sistemática se define por el registro planificado y objetivo de conductas en su contexto natural, usando categorías predefinidas y sin intervenir en la situación. El investigador es un observador externo que no altera el entorno.

\noindent\textbf{\color{gris}Por qué son incorrectas las otras opciones:}
\begin{itemize}
  \item {\small A — Incorrecto: la encuesta implica preguntar a los participantes; aquí no hay preguntas sino registro directo de conductas.}
  \item {\small C — Incorrecto: el diseño de caso único aplica y retira intervenciones para establecer causalidad; aquí solo se observa.}
  \item {\small D — Incorrecto: la investigación-acción implica co-participación del investigador para transformar la realidad.}
\end{itemize}
\vspace{8pt}\hrule\vspace{10pt}

\begin{mdframed}[backgroundcolor=fondogris,linecolor=azul,linewidth=1pt,innerleftmargin=10pt,innerrightmargin=10pt,innertopmargin=7pt,innerbottommargin=7pt]
\textbf{Ítem 41 / d3} \hfill {\small\bfseries\color{azul} Nivel 1 · Básico}
\end{mdframed}

Un equipo aplica el mismo cuestionario de satisfacción con la vida a estudiantes universitarios de tres ciudades distintas. El objetivo es comparar los niveles de satisfacción entre ciudades, sin manipular ninguna variable. Según León y Montero, ¿cómo se denomina la inclusión de las tres ciudades en este diseño descriptivo?

\begin{enumerate}
  \item[A.] Asignación aleatoria de participantes a condiciones
  \item[\textbf{\color{verde}B.}] \textbf{\color{verde} Comparación entre grupos naturales dentro del diseño descriptivo} \hfill $\checkmark$
  \item[C.] Diseño factorial de tres factores
  \item[D.] Bloqueo por localización geográfica
\end{enumerate}

\noindent\textbf{\color{azul}Explicación:} Dentro del diseño descriptivo por encuestas, es posible comparar grupos que ya existen en la realidad —grupos naturales como ciudades, sexos o edades— sin que ello convierta el diseño en experimental o cuasiexperimental. La comparación tiene carácter descriptivo y no permite inferencias causales.

\noindent\textbf{\color{gris}Por qué son incorrectas las otras opciones:}
\begin{itemize}
  \item {\small A — Incorrecto: la asignación aleatoria define al experimento verdadero; aquí los grupos son ciudades ya existentes, no asignadas.}
  \item {\small C — Incorrecto: un diseño factorial requiere cruzar dos o más variables independientes manipuladas; aquí no hay manipulación.}
  \item {\small D — Incorrecto: el bloqueo es una estrategia experimental para controlar una variable extraña antes de aleatorizar.}
\end{itemize}
\vspace{8pt}\hrule\vspace{10pt}

\begin{mdframed}[backgroundcolor=fondogris,linecolor=azul,linewidth=1pt,innerleftmargin=10pt,innerrightmargin=10pt,innertopmargin=7pt,innerbottommargin=7pt]
\textbf{Ítem 42 / n02} \hfill {\small\bfseries\color{azul} Nivel 1 · Básico}
\end{mdframed}

Un investigador quiere saber qué porcentaje de adolescentes de España consume videojuegos más de tres horas al día. Envía un cuestionario a 2.000 adolescentes seleccionados aleatoriamente de distintas comunidades autónomas. No compara grupos ni aplica ninguna intervención. ¿Qué diseño está usando?

\begin{enumerate}
  \item[A.] Diseño experimental de un factor, porque mide una sola variable
  \item[\textbf{\color{verde}B.}] \textbf{\color{verde} Diseño descriptivo por encuestas, porque recoge datos sobre una variable en una muestra sin manipular nada} \hfill $\checkmark$
  \item[C.] Diseño ex post facto prospectivo, porque hará seguimiento de los adolescentes en el futuro
  \item[D.] Diseño cuasiexperimental con pretest, porque obtiene datos antes de una posible intervención
\end{enumerate}

\noindent\textbf{\color{azul}Explicación:} El diseño descriptivo por encuestas se usa para describir cómo es una variable en una población, sin manipular condiciones. El objetivo aquí es puramente descriptivo: estimar la prevalencia del consumo excesivo de videojuegos.

\noindent\textbf{\color{gris}Por qué son incorrectas las otras opciones:}
\begin{itemize}
  \item {\small A — Incorrecto: medir una sola variable no define al diseño experimental; lo que lo define es la manipulación deliberada de la VI y la aleatorización.}
  \item {\small C — Incorrecto: el ex post facto prospectivo parte de una causa ya presente en los participantes y los sigue en el tiempo; aquí no hay seguimiento ni causa previa identificada.}
  \item {\small D — Incorrecto: el diseño cuasiexperimental requiere una intervención que el investigador aplica; aquí no hay ninguna.}
\end{itemize}
\vspace{8pt}\hrule\vspace{10pt}

\begin{mdframed}[backgroundcolor=fondogris,linecolor=azul,linewidth=1pt,innerleftmargin=10pt,innerrightmargin=10pt,innertopmargin=7pt,innerbottommargin=7pt]
\textbf{Ítem 43 / cov01} \hfill {\small\bfseries\color{azul} Nivel 1 · Básico}
\end{mdframed}

Un investigador observa que los países con mayor consumo de chocolate per cápita tienen más premios Nobel. Calcula el coeficiente de correlación y obtiene r=.79. ¿Qué demuestran estos datos?

\begin{enumerate}
  \item[A.] Que comer chocolate mejora la capacidad intelectual y por tanto aumenta los premios Nobel
  \item[\textbf{\color{verde}B.}] \textbf{\color{verde} Que existe una covariación estadística entre el consumo de chocolate y los premios Nobel, pero esto no implica que uno cause al otro} \hfill $\checkmark$
  \item[C.] Que el chocolate es una variable independiente que puede manipularse para incrementar los premios Nobel
  \item[D.] Que el diseño es experimental porque hay una variable que cambia con la otra
\end{enumerate}

\noindent\textbf{\color{azul}Explicación:} La covariación significa que dos variables cambian juntas de forma sistemática: cuando una sube, la otra también tiende a subir (o a bajar). Pero covariación no es causalidad: pueden existir terceras variables que expliquen la correlación (por ejemplo, el nivel de riqueza del país). Este es el principio de "correlación no implica causalidad".

\noindent\textbf{\color{gris}Por qué son incorrectas las otras opciones:}
\begin{itemize}
  \item {\small A — Incorrecto: inferir que el chocolate mejora la inteligencia a partir de una correlación es un error clásico de confundir correlación con causalidad.}
  \item {\small C — Incorrecto: en un diseño correlacional descriptivo ninguna variable se "manipula"; ambas se miden tal como existen.}
  \item {\small D — Incorrecto: el diseño experimental requiere manipulación deliberada de una VI por parte del investigador; aquí solo se miden datos existentes.}
\end{itemize}
\vspace{8pt}\hrule\vspace{10pt}

\begin{mdframed}[backgroundcolor=fondogris,linecolor=azul,linewidth=1pt,innerleftmargin=10pt,innerrightmargin=10pt,innertopmargin=7pt,innerbottommargin=7pt]
\textbf{Ítem 44 / n09} \hfill {\small\bfseries\color{azul} Nivel 1 · Básico}
\end{mdframed}

Un investigador quiere saber si los profesores de primaria y secundaria difieren en su nivel de satisfacción laboral. Aplica la misma escala de satisfacción a 200 profesores de cada nivel. No interviene ni manipula nada. ¿Qué diseño es este?

\begin{enumerate}
  \item[A.] Experimento entre grupos, porque compara dos grupos diferentes
  \item[\textbf{\color{verde}B.}] \textbf{\color{verde} Diseño descriptivo con comparación de grupos naturales, porque mide una variable en grupos ya existentes sin manipulación} \hfill $\checkmark$
  \item[C.] Diseño cuasiexperimental, porque los profesores no pueden asignarse aleatoriamente a los niveles educativos
  \item[D.] Diseño ex post facto, porque el nivel educativo en que trabajan ya es una variable preexistente
\end{enumerate}

\noindent\textbf{\color{azul}Explicación:} Cuando un investigador compara grupos naturales (profesores de primaria vs. secundaria) midiendo una variable sin manipular nada, está haciendo un diseño descriptivo con comparación de grupos. El nivel educativo no es una VI manipulada: es una característica preexistente de los participantes.

\noindent\textbf{\color{gris}Por qué son incorrectas las otras opciones:}
\begin{itemize}
  \item {\small A — Incorrecto: el experimento entre grupos requiere manipulación deliberada y asignación aleatoria; aquí no hay ninguna.}
  \item {\small C — Incorrecto: el cuasiexperimental requiere que el investigador aplique una intervención aunque sin aleatorizar; aquí solo se mide.}
  \item {\small D — Incorrecto: aunque el nivel educativo es preexistente, el objetivo no es estudiar sus efectos sobre un resultado futuro sino describir y comparar la satisfacción actual; es un diseño descriptivo.}
\end{itemize}
\vspace{8pt}\hrule\vspace{10pt}

\begin{mdframed}[backgroundcolor=fondogris,linecolor=azul,linewidth=1pt,innerleftmargin=10pt,innerrightmargin=10pt,innertopmargin=7pt,innerbottommargin=7pt]
\textbf{Ítem 45 / n14} \hfill {\small\bfseries\color{azul} Nivel 1 · Básico}
\end{mdframed}

Un investigador lleva a cabo un estudio en el que registra el número de infracciones de tráfico cometidas por conductores de distintas franjas de edad (18-25, 26-40, 41-60, más de 60 años). No interviene ni manipula nada. Solo describe la distribución de infracciones por grupo de edad. ¿Qué diseño es este?

\begin{enumerate}
  \item[A.] Experimento unifactorial entre grupos con cuatro niveles
  \item[\textbf{\color{verde}B.}] \textbf{\color{verde} Diseño descriptivo con comparación de grupos naturales por categorías de edad} \hfill $\checkmark$
  \item[C.] Diseño ex post facto retrospectivo, porque las infracciones pasadas ya han ocurrido
  \item[D.] Diseño cuasiexperimental, porque compara cuatro grupos sin aleatorización
\end{enumerate}

\noindent\textbf{\color{azul}Explicación:} Cuando el objetivo es simplemente describir cómo se distribuye una variable en distintos grupos naturales (aquí, la edad) sin manipular nada, el diseño es descriptivo. La comparación entre grupos de edad no implica manipulación ni inferencia causal.

\noindent\textbf{\color{gris}Por qué son incorrectas las otras opciones:}
\begin{itemize}
  \item {\small A — Incorrecto: el experimento unifactorial requiere que el investigador manipule los niveles de la VI; aquí la edad no se manipula.}
  \item {\small C — Incorrecto: el ex post facto retrospectivo parte de un efecto y busca causas pasadas; aquí el objetivo es simplemente describir la distribución actual de infracciones.}
  \item {\small D — Incorrecto: el cuasiexperimental requiere una intervención activa del investigador; aquí solo se mide y describe.}
\end{itemize}
\vspace{8pt}\hrule\vspace{10pt}

\begin{mdframed}[backgroundcolor=fondogris,linecolor=azul,linewidth=1pt,innerleftmargin=10pt,innerrightmargin=10pt,innertopmargin=7pt,innerbottommargin=7pt]
\textbf{Ítem 46 / de07} \hfill {\small\bfseries\color{azul} Nivel 1 · Básico}
\end{mdframed}

Una investigadora quiere describir los hábitos de ejercicio físico de los mayores de 65 años en España. Para ello aplica una encuesta telefónica a 1.500 personas seleccionadas mediante muestreo aleatorio simple. No compara grupos ni aplica ninguna intervención. ¿Qué tipo de muestreo garantiza que todos los mayores de 65 años tengan la misma probabilidad de ser seleccionados?

\begin{enumerate}
  \item[A.] Muestreo por conveniencia, porque se selecciona a quienes están disponibles
  \item[\textbf{\color{verde}B.}] \textbf{\color{verde} Muestreo aleatorio simple, porque todos los miembros de la población tienen la misma probabilidad de ser seleccionados} \hfill $\checkmark$
  \item[C.] Muestreo intencional, porque la investigadora elige a quienes mejor representan la experiencia de ejercicio
  \item[D.] Muestreo de bola de nieve, porque cada seleccionado propone a otros participantes
\end{enumerate}

\noindent\textbf{\color{azul}Explicación:} El muestreo aleatorio simple es el tipo de muestreo probabilístico más básico: cada miembro de la población tiene exactamente la misma probabilidad de ser incluido en la muestra. Esto garantiza la representatividad estadística y permite generalizar los resultados a la población con un error conocido.

\noindent\textbf{\color{gris}Por qué son incorrectas las otras opciones:}
\begin{itemize}
  \item {\small A — Incorrecto: el muestreo por conveniencia selecciona a quienes están más accesibles; no garantiza representatividad y es un muestreo no probabilístico.}
  \item {\small C — Incorrecto: el muestreo intencional o purposive selecciona deliberadamente a participantes con características específicas; es un muestreo no probabilístico propio de la investigación cualitativa.}
  \item {\small D — Incorrecto: el muestreo de bola de nieve se usa cuando la población es de difícil acceso y cada participante recluta a otros; no garantiza representatividad ni igual probabilidad de selección.}
\end{itemize}
\vspace{8pt}\hrule\vspace{10pt}

\begin{mdframed}[backgroundcolor=fondogris,linecolor=ambar,linewidth=1pt,innerleftmargin=10pt,innerrightmargin=10pt,innertopmargin=7pt,innerbottommargin=7pt]
\textbf{Ítem 47 / d4} \hfill {\small\bfseries\color{ambar} Nivel 2 · Intermedio}
\end{mdframed}

Una investigadora aplica una encuesta sobre hábitos de lectura a 800 estudiantes mediante muestreo estratificado proporcional por tipo de centro (público vs. privado) y curso escolar. ¿Qué ventaja aporta este tipo de muestreo respecto al muestreo aleatorio simple?

\begin{enumerate}
  \item[A.] Permite establecer relaciones causales entre tipo de centro y hábitos de lectura
  \item[\textbf{\color{verde}B.}] \textbf{\color{verde} Garantiza representación proporcional de todos los estratos, aumentando la precisión de las estimaciones} \hfill $\checkmark$
  \item[C.] Elimina el sesgo de autoinforme propio de los cuestionarios de opinión
  \item[D.] Convierte el diseño en cuasiexperimental al incluir grupos de comparación
\end{enumerate}

\noindent\textbf{\color{azul}Explicación:} El muestreo estratificado proporcional asegura que cada subgrupo relevante de la población (estrato) quede representado en la muestra en la misma proporción que en la población. Esto reduce el error muestral y permite comparaciones entre estratos con mayor precisión estadística.

\noindent\textbf{\color{gris}Por qué son incorrectas las otras opciones:}
\begin{itemize}
  \item {\small A — Incorrecto: ningún tipo de muestreo, por sofisticado que sea, permite inferir causalidad sin manipulación experimental.}
  \item {\small C — Incorrecto: el sesgo de autoinforme es una limitación del instrumento (cuestionario), no del tipo de muestreo.}
  \item {\small D — Incorrecto: incluir estratos no convierte el diseño en cuasiexperimental; este requiere manipulación de una VI.}
\end{itemize}
\vspace{8pt}\hrule\vspace{10pt}

\begin{mdframed}[backgroundcolor=fondogris,linecolor=ambar,linewidth=1pt,innerleftmargin=10pt,innerrightmargin=10pt,innertopmargin=7pt,innerbottommargin=7pt]
\textbf{Ítem 48 / d5} \hfill {\small\bfseries\color{ambar} Nivel 2 · Intermedio}
\end{mdframed}

Dos investigadores realizan un estudio sobre el estado emocional de los trabajadores de una empresa. Uno aplica un cuestionario en un único momento (enero). El otro aplica el mismo cuestionario tres veces (enero, junio y diciembre) a los mismos trabajadores. ¿Cuál es la ventaja principal del segundo diseño según León y Montero?

\begin{enumerate}
  \item[A.] Permite aleatorizar la asignación de trabajadores a condiciones
  \item[\textbf{\color{verde}B.}] \textbf{\color{verde} Permite analizar cambios y tendencias a lo largo del tiempo en la misma muestra} \hfill $\checkmark$
  \item[C.] Elimina las amenazas a la validez interna propias de los estudios transversales
  \item[D.] Permite inferir que el tiempo transcurrido causa los cambios observados
\end{enumerate}

\noindent\textbf{\color{azul}Explicación:} El diseño longitudinal mide las mismas variables en los mismos participantes a lo largo del tiempo, lo que permite detectar cambios intraindividuales y tendencias temporales. Es una modalidad del diseño descriptivo por encuestas y no permite inferencias causales, pero sí seguir la evolución de los fenómenos.

\noindent\textbf{\color{gris}Por qué son incorrectas las otras opciones:}
\begin{itemize}
  \item {\small A — Incorrecto: la aleatorización no es relevante en diseños descriptivos y no es posible en este contexto.}
  \item {\small C — Incorrecto: los diseños longitudinales tienen sus propias amenazas, como la mortalidad experimental y los efectos de práctica.}
  \item {\small D — Incorrecto: medir en varios momentos no equivale a manipular el tiempo; la correlación temporal no implica causalidad.}
\end{itemize}
\vspace{8pt}\hrule\vspace{10pt}

\begin{mdframed}[backgroundcolor=fondogris,linecolor=ambar,linewidth=1pt,innerleftmargin=10pt,innerrightmargin=10pt,innertopmargin=7pt,innerbottommargin=7pt]
\textbf{Ítem 49 / de08} \hfill {\small\bfseries\color{ambar} Nivel 2 · Intermedio}
\end{mdframed}

Un investigador aplica una batería de tests psicológicos a 300 adolescentes y calcula la correlación entre el tiempo frente a pantallas y el rendimiento académico, obteniendo r=-.48. Su colega le pregunta si puede concluir que reducir el tiempo de pantallas mejoraría el rendimiento. ¿Qué le respondería correctamente el investigador?

\begin{enumerate}
  \item[A.] Sí, porque la correlación negativa indica que reducir las pantallas causa mejoras en el rendimiento
  \item[\textbf{\color{verde}B.}] \textbf{\color{verde} No, porque para inferir causalidad hacen falta las tres condiciones: covariación (presente), antecesión temporal (no establecida) y descarte de terceras variables (no realizado)} \hfill $\checkmark$
  \item[C.] Sí, siempre que la correlación sea estadísticamente significativa con p<.05
  \item[D.] No, porque r=-.48 es una correlación demasiado débil para tener implicaciones prácticas
\end{enumerate}

\noindent\textbf{\color{azul}Explicación:} La covariación (r=-.48) es solo la primera de las tres condiciones para inferir causalidad. Sin datos longitudinales no puede saberse si el tiempo de pantallas precede al bajo rendimiento o al revés (antecesión temporal). Además, terceras variables como el nivel socioeconómico o el tiempo de estudio podrían explicar ambas variables simultáneamente.

\noindent\textbf{\color{gris}Por qué son incorrectas las otras opciones:}
\begin{itemize}
  \item {\small A — Incorrecto: una correlación negativa no implica causalidad; solo indica que las dos variables varían en sentido contrario, pero no que una cause la otra.}
  \item {\small C — Incorrecto: la significación estadística solo indica que la correlación no es debida al azar en la muestra; no establece causalidad.}
  \item {\small D — Incorrecto: r=-.48 es una correlación de magnitud moderada con relevancia práctica; el problema no es la magnitud de la correlación sino la imposibilidad de inferir causalidad a partir de ella.}
\end{itemize}
\vspace{8pt}\hrule\vspace{10pt}

\begin{mdframed}[backgroundcolor=fondogris,linecolor=ambar,linewidth=1pt,innerleftmargin=10pt,innerrightmargin=10pt,innertopmargin=7pt,innerbottommargin=7pt]
\textbf{Ítem 50 / de09} \hfill {\small\bfseries\color{ambar} Nivel 2 · Intermedio}
\end{mdframed}

Una investigadora quiere estudiar si la satisfacción laboral de los empleados cambia a lo largo de un año. Para ello mide la satisfacción en los mismos 200 empleados en enero, abril, julio y octubre. ¿Cuál es la denominación correcta de este tipo de diseño descriptivo y cuál es su limitación principal?

\begin{enumerate}
  \item[A.] Diseño transversal; limitación: no permite comparar distintos grupos de empleados
  \item[\textbf{\color{verde}B.}] \textbf{\color{verde} Diseño longitudinal o de panel; limitación: la mortalidad experimental (pérdida de participantes) puede sesgar los resultados si los que abandonan difieren de los que permanecen} \hfill $\checkmark$
  \item[C.] Diseño cuasiexperimental de series temporales; limitación: no hay grupo de comparación sin intervención
  \item[D.] Diseño experimental de medidas repetidas; limitación: no hay manipulación de ninguna variable independiente
\end{enumerate}

\noindent\textbf{\color{azul}Explicación:} El diseño longitudinal o de panel mide las mismas variables en los mismos participantes en varios momentos temporales. Su principal amenaza es la mortalidad experimental o atrición: si los participantes que abandonan el estudio son sistemáticamente distintos de los que permanecen (por ejemplo, si los más insatisfechos dejan la empresa), los resultados finales pueden sobrestimar la satisfacción.

\noindent\textbf{\color{gris}Por qué son incorrectas las otras opciones:}
\begin{itemize}
  \item {\small A — Incorrecto: el diseño transversal mide en un único momento temporal; aquí hay cuatro medidas a lo largo del año, lo que lo convierte en longitudinal.}
  \item {\small C — Incorrecto: el diseño de series temporales cuasiexperimental requiere una intervención que se evalúa; aquí no hay ninguna intervención, solo medición descriptiva.}
  \item {\small D — Incorrecto: el diseño experimental de medidas repetidas requiere manipulación de al menos una variable independiente; aquí solo se mide la satisfacción sin manipular nada.}
\end{itemize}
\vspace{8pt}\hrule\vspace{10pt}

\begin{mdframed}[backgroundcolor=fondogris,linecolor=rojo,linewidth=1pt,innerleftmargin=10pt,innerrightmargin=10pt,innertopmargin=7pt,innerbottommargin=7pt]
\textbf{Ítem 51 / d6} \hfill {\small\bfseries\color{rojo} Nivel 3 · Avanzado}
\end{mdframed}

Una investigadora concluye que, dado que en su encuesta encontró una correlación de r=.62 entre el número de horas de ejercicio y el bienestar psicológico, el ejercicio físico mejora el bienestar. Su colega le señala que esta conclusión no se desprende del diseño. ¿Cuál es la limitación fundamental que señala el colega, según León y Montero?

\begin{enumerate}
  \item[A.] La muestra es insuficiente para detectar una correlación de esa magnitud con potencia estadística adecuada
  \item[\textbf{\color{verde}B.}] \textbf{\color{verde} La correlación solo mide la fuerza de asociación lineal, pero no establece dirección causal ni descarta terceras variables explicativas} \hfill $\checkmark$
  \item[C.] El coeficiente r de Pearson no es adecuado para datos de encuesta
  \item[D.] El diseño longitudinal habría sido necesario para calcular correlaciones válidas
\end{enumerate}

\noindent\textbf{\color{azul}Explicación:} La limitación fundamental del diseño descriptivo correlacional es que no puede establecer causalidad. Una correlación alta implica covariación, pero no permite saber si A causa B, si B causa A, o si ambas son causadas por una tercera variable (variable de confusión). Solo el experimento con manipulación y aleatorización permite inferir causalidad.

\noindent\textbf{\color{gris}Por qué son incorrectas las otras opciones:}
\begin{itemize}
  \item {\small A — Incorrecto: una correlación de .62 es moderada-alta; el tamaño muestral no es el problema central señalado.}
  \item {\small C — Incorrecto: el coeficiente r de Pearson es perfectamente aplicable a datos de encuesta con escalas continuas.}
  \item {\small D — Incorrecto: la correlación puede calcularse tanto en diseños transversales como longitudinales; el problema es de inferencia causal.}
\end{itemize}
\vspace{8pt}\hrule\vspace{10pt}

\begin{mdframed}[backgroundcolor=fondogris,linecolor=rojo,linewidth=1pt,innerleftmargin=10pt,innerrightmargin=10pt,innertopmargin=7pt,innerbottommargin=7pt]
\textbf{Ítem 52 / de10} \hfill {\small\bfseries\color{rojo} Nivel 3 · Avanzado}
\end{mdframed}

Un investigador quiere estimar la prevalencia de la depresión en adultos mayores de una ciudad. Selecciona una muestra usando muestreo por conglomerados: primero elige aleatoriamente 10 barrios, luego elige aleatoriamente 5 centros de salud en cada barrio, y finalmente selecciona aleatoriamente a 20 adultos de cada centro. ¿Cuál es la principal desventaja de este tipo de muestreo respecto al muestreo aleatorio simple?

\begin{enumerate}
  \item[A.] Que no puede calcularse ningún estadístico inferencial con muestras por conglomerados
  \item[B.] Que es más costoso económicamente que el muestreo aleatorio simple para el mismo tamaño muestral
  \item[\textbf{\color{verde}C.}] \textbf{\color{verde} Que la homogeneidad dentro de los conglomerados (barrios, centros de salud) puede reducir la precisión de las estimaciones, requiriendo muestras más grandes para el mismo nivel de error que el muestreo aleatorio simple} \hfill $\checkmark$
  \item[D.] Que los adultos dentro del mismo centro de salud tienen más probabilidad de ser seleccionados que los de centros con más pacientes
\end{enumerate}

\noindent\textbf{\color{azul}Explicación:} El muestreo por conglomerados es logísticamente eficiente (no hay que contactar a personas dispersas) pero estadísticamente menos preciso que el muestreo aleatorio simple. Esto se debe a que las personas dentro del mismo conglomerado (barrio, centro de salud) tienden a parecerse entre sí (homogeneidad intracluster), por lo que aportan menos información independiente que personas seleccionadas al azar de toda la población.

\noindent\textbf{\color{gris}Por qué son incorrectas las otras opciones:}
\begin{itemize}
  \item {\small A — Incorrecto: con el muestreo por conglomerados sí pueden calcularse estadísticos inferenciales; simplemente hay que tener en cuenta el efecto de diseño en los errores estándar.}
  \item {\small B — Incorrecto: una de las ventajas del muestreo por conglomerados es precisamente que suele ser más barato que el muestreo aleatorio simple, ya que concentra la recogida de datos en lugares específicos.}
  \item {\small D — Incorrecto: en el muestreo estratificado proporcional pueden surgir diferencias de probabilidad; en el muestreo por conglomerados bien diseñado se igualan las probabilidades de selección.}
\end{itemize}
\vspace{8pt}\hrule\vspace{10pt}

\section{Cualitativo (CUAL) · T.12--15}
{\small\color{gris} 14 ítems · Niveles: L1=6 · L2=5 · L3=3}

\begin{mdframed}[backgroundcolor=fondogris,linecolor=azul,linewidth=1pt,innerleftmargin=10pt,innerrightmargin=10pt,innertopmargin=7pt,innerbottommargin=7pt]
\textbf{Ítem 53 / q1} \hfill {\small\bfseries\color{azul} Nivel 1 · Básico}
\end{mdframed}

Una investigadora convive durante ocho meses en una comunidad rural de montaña para entender cómo sus habitantes perciben los cambios en el clima y cómo estos afectan a sus prácticas agrícolas tradicionales. Participa en las actividades cotidianas, realiza entrevistas informales y toma notas de campo. ¿Qué diseño cualitativo describe mejor este estudio?

\begin{enumerate}
  \item[A.] Fenomenología, porque estudia la experiencia vivida de los cambios climáticos
  \item[B.] Investigación-acción, porque busca transformar las prácticas agrícolas de la comunidad
  \item[\textbf{\color{verde}C.}] \textbf{\color{verde} Etnografía, porque implica inmersión prolongada para comprender la cultura del grupo desde dentro} \hfill $\checkmark$
  \item[D.] Estudio de casos, porque analiza en profundidad una comunidad concreta
\end{enumerate}

\noindent\textbf{\color{azul}Explicación:} La etnografía se define por la inmersión prolongada del investigador en el contexto de estudio mediante observación participante. El objetivo es comprender la cultura, las prácticas, los significados y las normas del grupo desde la perspectiva de sus propios miembros (perspectiva émica). Es la tradición metodológica más ligada a la antropología cultural.

\noindent\textbf{\color{gris}Por qué son incorrectas las otras opciones:}
\begin{itemize}
  \item {\small A — Incorrecto: la fenomenología se centra en la estructura de la experiencia vivida de un fenómeno concreto, normalmente mediante entrevistas en profundidad; no requiere inmersión prolongada en la comunidad.}
  \item {\small B — Incorrecto: la investigación-acción implica el objetivo explícito de transformar la realidad junto a los participantes, en ciclos de acción y reflexión; aquí solo hay observación y comprensión.}
  \item {\small D — Incorrecto: el estudio de casos analiza una unidad delimitada usando múltiples fuentes de evidencia; la etnografía tiene un marco epistemológico y metodológico propio centrado en la cultura.}
\end{itemize}
\vspace{8pt}\hrule\vspace{10pt}

\begin{mdframed}[backgroundcolor=fondogris,linecolor=azul,linewidth=1pt,innerleftmargin=10pt,innerrightmargin=10pt,innertopmargin=7pt,innerbottommargin=7pt]
\textbf{Ítem 54 / q2} \hfill {\small\bfseries\color{azul} Nivel 1 · Básico}
\end{mdframed}

Un investigador quiere entender cómo viven la experiencia de la infertilidad las parejas que llevan más de tres años en tratamiento de reproducción asistida. Realiza entrevistas en profundidad y analiza sus narrativas buscando temas y estructuras de significado comunes. ¿Qué diseño cualitativo es el más adecuado?

\begin{enumerate}
  \item[A.] Etnografía, porque requiere observación participante prolongada en la clínica
  \item[B.] Investigación-acción, porque busca transformar las condiciones de atención a los pacientes
  \item[\textbf{\color{verde}C.}] \textbf{\color{verde} Estudio fenomenológico, porque busca describir la estructura de la experiencia vivida de la infertilidad} \hfill $\checkmark$
  \item[D.] Estudio de casos múltiples, porque incluye a más de una pareja en el análisis
\end{enumerate}

\noindent\textbf{\color{azul}Explicación:} La fenomenología estudia cómo las personas viven y dan sentido a sus experiencias en primera persona. El Análisis Fenomenológico Interpretativo (AFI) es el método más utilizado en psicología clínica y de la salud. Se aplica cuando la pregunta de investigación es "¿cómo viven/experimentan X las personas?" y el interés está en la subjetividad y el significado.

\noindent\textbf{\color{gris}Por qué son incorrectas las otras opciones:}
\begin{itemize}
  \item {\small A — Incorrecto: la etnografía requiere inmersión prolongada en el contexto natural; el estudio de la infertilidad se centra en la experiencia subjetiva, no en la cultura de una clínica.}
  \item {\small B — Incorrecto: la investigación-acción tiene un objetivo transformador explícito; aquí el objetivo es comprensivo, no de cambio.}
  \item {\small D — Incorrecto: incluir múltiples participantes no convierte un estudio en estudio de casos múltiples; el estudio de casos se define por el análisis en profundidad de unidades delimitadas, no por el número de entrevistados.}
\end{itemize}
\vspace{8pt}\hrule\vspace{10pt}

\begin{mdframed}[backgroundcolor=fondogris,linecolor=azul,linewidth=1pt,innerleftmargin=10pt,innerrightmargin=10pt,innertopmargin=7pt,innerbottommargin=7pt]
\textbf{Ítem 55 / n06} \hfill {\small\bfseries\color{azul} Nivel 1 · Básico}
\end{mdframed}

Una investigadora pasa tres meses en un barrio de inmigrantes, participa en sus fiestas, acude a sus reuniones de vecinos y toma notas detalladas sobre sus costumbres y formas de relacionarse. Su objetivo es entender la cultura del barrio desde dentro. ¿Qué diseño cualitativo usa?

\begin{enumerate}
  \item[A.] Fenomenología, porque estudia la experiencia vivida de los inmigrantes
  \item[B.] Investigación-acción, porque quiere ayudar al barrio a resolver sus problemas
  \item[\textbf{\color{verde}C.}] \textbf{\color{verde} Etnografía, porque implica inmersión prolongada para comprender la cultura de un grupo} \hfill $\checkmark$
  \item[D.] Estudio de casos, porque analiza en profundidad un barrio concreto
\end{enumerate}

\noindent\textbf{\color{azul}Explicación:} La etnografía se define por la inmersión prolongada del investigador mediante observación participante, con el objetivo de describir y comprender la cultura de un grupo desde la perspectiva de sus propios miembros. Es el método clásico de la antropología cultural.

\noindent\textbf{\color{gris}Por qué son incorrectas las otras opciones:}
\begin{itemize}
  \item {\small A — Incorrecto: la fenomenología estudia experiencias individuales vividas mediante entrevistas; no requiere inmersión en la comunidad.}
  \item {\small B — Incorrecto: la investigación-acción busca transformar la realidad junto a los participantes; aquí solo se observa y comprende.}
  \item {\small D — Incorrecto: el estudio de casos analiza una unidad delimitada usando múltiples fuentes; la etnografía tiene un marco metodológico propio centrado en la cultura.}
\end{itemize}
\vspace{8pt}\hrule\vspace{10pt}

\begin{mdframed}[backgroundcolor=fondogris,linecolor=azul,linewidth=1pt,innerleftmargin=10pt,innerrightmargin=10pt,innertopmargin=7pt,innerbottommargin=7pt]
\textbf{Ítem 56 / n15} \hfill {\small\bfseries\color{azul} Nivel 1 · Básico}
\end{mdframed}

Un investigador realiza entrevistas en profundidad a diez personas que han sufrido un ictus para entender cómo viven el proceso de rehabilitación y qué significado le dan a los pequeños avances. ¿Qué diseño cualitativo describe mejor este estudio?

\begin{enumerate}
  \item[A.] Etnografía, porque requiere convivir prolongadamente con los pacientes en el hospital
  \item[B.] Investigación-acción, porque busca mejorar el protocolo de rehabilitación con los pacientes
  \item[\textbf{\color{verde}C.}] \textbf{\color{verde} Fenomenología, porque busca describir la estructura de la experiencia vivida del proceso de rehabilitación} \hfill $\checkmark$
  \item[D.] Estudio de casos, porque analiza en profundidad a diez personas individuales
\end{enumerate}

\noindent\textbf{\color{azul}Explicación:} La fenomenología se aplica cuando la pregunta es "¿cómo viven y qué significado dan las personas a una experiencia concreta?". Las entrevistas en profundidad son el instrumento habitual y el análisis busca la estructura esencial de la experiencia vivida (en este caso, la rehabilitación tras un ictus).

\noindent\textbf{\color{gris}Por qué son incorrectas las otras opciones:}
\begin{itemize}
  \item {\small A — Incorrecto: la etnografía implica inmersión prolongada para comprender la cultura de un grupo; aquí el objetivo es entender la experiencia subjetiva individual, no la cultura hospitalaria.}
  \item {\small B — Incorrecto: la investigación-acción tiene un objetivo transformador explícito (cambiar el protocolo); aquí el objetivo es comprender la experiencia, no transformar la práctica.}
  \item {\small D — Incorrecto: el estudio de casos analiza una unidad delimitada (una persona, un programa, una institución) en su contexto; aquí el interés no es el caso individual sino la estructura de una experiencia compartida.}
\end{itemize}
\vspace{8pt}\hrule\vspace{10pt}

\begin{mdframed}[backgroundcolor=fondogris,linecolor=azul,linewidth=1pt,innerleftmargin=10pt,innerrightmargin=10pt,innertopmargin=7pt,innerbottommargin=7pt]
\textbf{Ítem 57 / ql08} \hfill {\small\bfseries\color{azul} Nivel 1 · Básico}
\end{mdframed}

Un investigador quiere entender por qué muchos estudiantes universitarios abandonan los estudios en el primer año. Realiza entrevistas en profundidad a 15 estudiantes que abandonaron y analiza sus narrativas para identificar patrones de significado compartidos. No parte de hipótesis previas. ¿Qué diseño cualitativo se ajusta mejor?

\begin{enumerate}
  \item[A.] Observación participante etnográfica, porque hay que conocer la cultura universitaria desde dentro
  \item[B.] Estudio de casos múltiples, porque hay 15 casos individuales que analizar
  \item[\textbf{\color{verde}C.}] \textbf{\color{verde} Diseño cualitativo con análisis temático inductivo, porque busca patrones de significado en las narrativas sin hipótesis previas} \hfill $\checkmark$
  \item[D.] Investigación-acción, porque quiere cambiar la política de abandono universitario
\end{enumerate}

\noindent\textbf{\color{azul}Explicación:} El análisis temático inductivo es el procedimiento de codificación y categorización que parte de los datos sin teoría previa para identificar temas recurrentes en las narrativas. Es el método de análisis cualitativo más utilizado en psicología y ciencias sociales cuando no se parte de una teoría específica sino de las experiencias de los participantes.

\noindent\textbf{\color{gris}Por qué son incorrectas las otras opciones:}
\begin{itemize}
  \item {\small A — Incorrecto: la observación participante requiere que el investigador comparta el contexto de los participantes de forma prolongada; aquí el investigador solo realiza entrevistas.}
  \item {\small B — Incorrecto: el estudio de casos analiza una unidad en profundidad en su contexto; aquí el interés no es cada estudiante como caso sino los patrones de significado compartidos entre ellos.}
  \item {\small D — Incorrecto: la investigación-acción tiene un objetivo transformador explícito (cambiar algo); aquí el objetivo es puramente comprensivo (entender por qué abandonan).}
\end{itemize}
\vspace{8pt}\hrule\vspace{10pt}

\begin{mdframed}[backgroundcolor=fondogris,linecolor=azul,linewidth=1pt,innerleftmargin=10pt,innerrightmargin=10pt,innertopmargin=7pt,innerbottommargin=7pt]
\textbf{Ítem 58 / ql09} \hfill {\small\bfseries\color{azul} Nivel 1 · Básico}
\end{mdframed}

Una investigadora estudia cómo los médicos residentes aprenden a gestionar la incertidumbre clínica durante su formación. Recoge datos mediante observaciones en la planta, entrevistas semiestructuradas y análisis de registros clínicos. Su objetivo es construir una teoría sobre ese proceso de aprendizaje a partir de los datos. ¿Qué diseño cualitativo está usando?

\begin{enumerate}
  \item[A.] Fenomenología, porque estudia la experiencia vivida de la incertidumbre
  \item[B.] Estudio de casos, porque analiza un hospital concreto
  \item[\textbf{\color{verde}C.}] \textbf{\color{verde} Teoría fundamentada, porque construye teoría inductivamente a partir de múltiples fuentes de datos} \hfill $\checkmark$
  \item[D.] Etnografía, porque observa in situ las prácticas del grupo
\end{enumerate}

\noindent\textbf{\color{azul}Explicación:} La teoría fundamentada (Grounded Theory) se distingue de los demás diseños cualitativos por su objetivo explícito de construir teoría a partir de los datos mediante codificación sistemática y comparación constante. El uso de múltiples fuentes (observación, entrevistas, documentos) y el muestreo teórico son rasgos adicionales de este enfoque.

\noindent\textbf{\color{gris}Por qué son incorrectas las otras opciones:}
\begin{itemize}
  \item {\small A — Incorrecto: la fenomenología se centra en la estructura de la experiencia vivida de un fenómeno específico; no tiene como objetivo construir teoría sino describir la esencia de una experiencia.}
  \item {\small B — Incorrecto: el estudio de casos analiza una unidad delimitada (el hospital) como objeto de estudio en sí mismo; la teoría fundamentada usa el contexto como fuente de datos para construir teoría, no como fin en sí mismo.}
  \item {\small D — Incorrecto: la etnografía se centra en la descripción cultural mediante inmersión prolongada; aunque puede haber observación, el objetivo de construir teoría es específico de la teoría fundamentada.}
\end{itemize}
\vspace{8pt}\hrule\vspace{10pt}

\begin{mdframed}[backgroundcolor=fondogris,linecolor=ambar,linewidth=1pt,innerleftmargin=10pt,innerrightmargin=10pt,innertopmargin=7pt,innerbottommargin=7pt]
\textbf{Ítem 59 / q3} \hfill {\small\bfseries\color{ambar} Nivel 2 · Intermedio}
\end{mdframed}

En una investigación cualitativa, la investigadora lleva a cabo 12 entrevistas y concluye que las últimas tres no aportaron ninguna categoría ni tema nuevo respecto a las nueve anteriores. ¿Cómo se denomina este criterio de cierre y qué justifica su uso?

\begin{enumerate}
  \item[A.] Validez convergente, porque los datos de distintos participantes convergen en los mismos temas
  \item[\textbf{\color{verde}B.}] \textbf{\color{verde} Saturación teórica, porque añadir nuevos datos no modifica ni enriquece las categorías analíticas ya desarrolladas} \hfill $\checkmark$
  \item[C.] Fiabilidad interjueces, porque dos investigadores independientes obtienen los mismos temas
  \item[D.] Transferibilidad, porque los resultados se pueden aplicar a contextos similares
\end{enumerate}

\noindent\textbf{\color{azul}Explicación:} La saturación teórica es el criterio cualitativo para decidir cuándo dejar de recoger datos: se alcanza cuando la incorporación de nuevos participantes o casos no genera categorías, temas o relaciones nuevas. Justifica la suficiencia de la muestra en términos de riqueza de información, no de representatividad estadística.

\noindent\textbf{\color{gris}Por qué son incorrectas las otras opciones:}
\begin{itemize}
  \item {\small A — Incorrecto: la validez convergente es un concepto de validez de constructo en investigación cuantitativa; no es el criterio de cierre muestral en estudios cualitativos.}
  \item {\small C — Incorrecto: la fiabilidad interjueces mide el acuerdo entre codificadores independientes; es un criterio de rigor del proceso analítico, no de la cantidad de datos.}
  \item {\small D — Incorrecto: la transferibilidad es el equivalente cualitativo de la validez externa y se refiere a la posibilidad de aplicar los hallazgos a otros contextos, no al criterio de cierre muestral.}
\end{itemize}
\vspace{8pt}\hrule\vspace{10pt}

\begin{mdframed}[backgroundcolor=fondogris,linecolor=ambar,linewidth=1pt,innerleftmargin=10pt,innerrightmargin=10pt,innertopmargin=7pt,innerbottommargin=7pt]
\textbf{Ítem 60 / q4} \hfill {\small\bfseries\color{ambar} Nivel 2 · Intermedio}
\end{mdframed}

Una investigadora usa tres fuentes de datos para estudiar el liderazgo en una organización: entrevistas a directivos, observación de reuniones y análisis de documentos internos. ¿Qué estrategia de rigor está aplicando y cuál es su propósito?

\begin{enumerate}
  \item[A.] Saturación teórica; el propósito es asegurarse de que los datos son suficientes para cerrar el estudio
  \item[\textbf{\color{verde}B.}] \textbf{\color{verde} Triangulación de fuentes; el propósito es aumentar la credibilidad de los hallazgos verificando que diferentes fuentes apuntan en la misma dirección} \hfill $\checkmark$
  \item[C.] Reducción fenomenológica; el propósito es eliminar los supuestos previos del investigador sobre el liderazgo
  \item[D.] Dependencia o auditabilidad; el propósito es que un revisor externo pueda reconstruir el proceso analítico
\end{enumerate}

\noindent\textbf{\color{azul}Explicación:} La triangulación es la estrategia de combinar múltiples fuentes de datos, métodos, investigadores o teorías para aumentar la credibilidad de los resultados cualitativos. Cuando distintas fuentes convergen en los mismos hallazgos, se incrementa la confianza en su validez. La triangulación de fuentes compara datos recogidos a través de diferentes métodos sobre el mismo fenómeno.

\noindent\textbf{\color{gris}Por qué son incorrectas las otras opciones:}
\begin{itemize}
  \item {\small A — Incorrecto: la saturación teórica se refiere al criterio de cierre muestral, no al uso de múltiples fuentes de datos.}
  \item {\small C — Incorrecto: la reducción fenomenológica es un movimiento epistemológico del investigador para poner entre paréntesis sus supuestos; no es el nombre de una estrategia de recogida de datos multifuente.}
  \item {\small D — Incorrecto: la dependencia o auditabilidad se refiere a la posibilidad de que un revisor externo reconstruya el proceso analítico a partir de un rastro documental del estudio.}
\end{itemize}
\vspace{8pt}\hrule\vspace{10pt}

\begin{mdframed}[backgroundcolor=fondogris,linecolor=ambar,linewidth=1pt,innerleftmargin=10pt,innerrightmargin=10pt,innertopmargin=7pt,innerbottommargin=7pt]
\textbf{Ítem 61 / q5} \hfill {\small\bfseries\color{ambar} Nivel 2 · Intermedio}
\end{mdframed}

Un equipo de investigación usa la teoría fundamentada (grounded theory) para estudiar cómo los profesores noveles aprenden a gestionar el aula. ¿Cuál es la característica más distintiva de este enfoque respecto a otros métodos cualitativos?

\begin{enumerate}
  \item[A.] Que parte de un marco teórico establecido y lo aplica al análisis de los datos cualitativos recogidos
  \item[\textbf{\color{verde}B.}] \textbf{\color{verde} Que construye teoría inductivamente a partir de los datos, sin partir de hipótesis previas, mediante muestreo teórico y análisis comparativo constante} \hfill $\checkmark$
  \item[C.] Que requiere observación participante prolongada del investigador en el aula natural de los profesores
  \item[D.] Que analiza únicamente documentos y textos escritos sin entrevistas ni observaciones directas
\end{enumerate}

\noindent\textbf{\color{azul}Explicación:} La teoría fundamentada (Glaser y Strauss) genera teoría inductivamente a partir de los datos mediante un proceso de codificación sistemática y comparación constante. El muestreo teórico dirige la recogida de nuevos datos basándose en los conceptos emergentes del análisis previo. El resultado es una teoría sustantiva enraizada en los datos, no una mera descripción.

\noindent\textbf{\color{gris}Por qué son incorrectas las otras opciones:}
\begin{itemize}
  \item {\small A — Incorrecto: la teoría fundamentada parte de los datos para construir teoría, no aplica teorías preexistentes; partir de un marco teórico previo es característico del enfoque hipotético-deductivo cuantitativo.}
  \item {\small C — Incorrecto: la observación participante prolongada define a la etnografía; la teoría fundamentada puede usar entrevistas, observaciones y documentos, pero su sello no es la inmersión etnográfica.}
  \item {\small D — Incorrecto: el análisis exclusivo de documentos es característico del análisis documental o del análisis de contenido; la teoría fundamentada suele combinar múltiples fuentes.}
\end{itemize}
\vspace{8pt}\hrule\vspace{10pt}

\begin{mdframed}[backgroundcolor=fondogris,linecolor=ambar,linewidth=1pt,innerleftmargin=10pt,innerrightmargin=10pt,innertopmargin=7pt,innerbottommargin=7pt]
\textbf{Ítem 62 / ql10} \hfill {\small\bfseries\color{ambar} Nivel 2 · Intermedio}
\end{mdframed}

Una investigadora cualitativa quiere garantizar que sus interpretaciones de las entrevistas reflejan fielmente las experiencias de los participantes. Para ello, comparte las categorías y temas emergentes con los propios participantes para que confirmen o cuestionen las interpretaciones. ¿Cómo se llama esta estrategia y a qué criterio de rigor cualitativo contribuye?

\begin{enumerate}
  \item[A.] Triangulación metodológica; contribuye a la dependencia del estudio
  \item[B.] Revisión por pares o auditoría; contribuye a la confirmabilidad del estudio
  \item[\textbf{\color{verde}C.}] \textbf{\color{verde} Validación por los participantes o member checking; contribuye a la credibilidad del estudio} \hfill $\checkmark$
  \item[D.] Muestreo teórico; contribuye a la transferibilidad del estudio
\end{enumerate}

\noindent\textbf{\color{azul}Explicación:} El member checking (validación por los participantes) consiste en devolver a los participantes las interpretaciones del investigador para que las validen, corrijan o enriquezcan. Es una estrategia clave para aumentar la credibilidad (equivalente cualitativo de la validez interna): si los participantes reconocen sus experiencias en las interpretaciones del investigador, aumenta la confianza en que esas interpretaciones son fieles a la realidad estudiada.

\noindent\textbf{\color{gris}Por qué son incorrectas las otras opciones:}
\begin{itemize}
  \item {\small A — Incorrecto: la triangulación combina múltiples fuentes, métodos o investigadores; el member checking solo involucra a los participantes para validar interpretaciones específicas.}
  \item {\small B — Incorrecto: la auditoría o revisión por pares es una estrategia para evaluar el proceso analítico desde fuera; el member checking implica a los propios participantes del estudio.}
  \item {\small D — Incorrecto: el muestreo teórico es la estrategia de la teoría fundamentada por la que se seleccionan nuevos participantes o casos en función de los conceptos emergentes del análisis.}
\end{itemize}
\vspace{8pt}\hrule\vspace{10pt}

\begin{mdframed}[backgroundcolor=fondogris,linecolor=ambar,linewidth=1pt,innerleftmargin=10pt,innerrightmargin=10pt,innertopmargin=7pt,innerbottommargin=7pt]
\textbf{Ítem 63 / ql11} \hfill {\small\bfseries\color{ambar} Nivel 2 · Intermedio}
\end{mdframed}

Un investigador utiliza un enfoque de estudio de casos para analizar cómo tres escuelas rurales distintas han logrado mantener alta la motivación de sus alumnos pese a tener escasos recursos. Combina entrevistas a docentes, observaciones de aula y análisis de documentos institucionales. ¿Cómo se denomina el tipo de estudio de casos que usa múltiples casos para iluminar un fenómeno general?

\begin{enumerate}
  \item[A.] Estudio de caso intrínseco, porque cada escuela es interesante por sí misma
  \item[\textbf{\color{verde}B.}] \textbf{\color{verde} Estudio de caso colectivo o múltiple, porque analiza varios casos para comprender mejor el fenómeno} \hfill $\checkmark$
  \item[C.] Estudio de caso instrumental único, porque el instrumento de recogida de datos es múltiple
  \item[D.] Diseño de series temporales de casos, porque las escuelas se analizan en distintos momentos del tiempo
\end{enumerate}

\noindent\textbf{\color{azul}Explicación:} El estudio de casos múltiples o colectivo (Stake) analiza varios casos simultáneamente para construir una comprensión más completa y contrastada del fenómeno de interés. Los resultados de los casos individuales se comparan entre sí para identificar patrones comunes y diferencias, lo que fortalece la transferibilidad de las conclusiones.

\noindent\textbf{\color{gris}Por qué son incorrectas las otras opciones:}
\begin{itemize}
  \item {\small A — Incorrecto: el caso intrínseco es aquel en que el investigador estudia el caso porque es inherentemente interesante en sí mismo, no para iluminar un fenómeno general.}
  \item {\small C — Incorrecto: el término "instrumental único" se refiere a un solo caso que se estudia para entender algo más allá de él mismo; aquí hay tres casos, lo que lo convierte en múltiple o colectivo.}
  \item {\small D — Incorrecto: el diseño de series temporales es un diseño cuasiexperimental cuantitativo; no es una categoría dentro de los estudios de casos cualitativos.}
\end{itemize}
\vspace{8pt}\hrule\vspace{10pt}

\begin{mdframed}[backgroundcolor=fondogris,linecolor=rojo,linewidth=1pt,innerleftmargin=10pt,innerrightmargin=10pt,innertopmargin=7pt,innerbottommargin=7pt]
\textbf{Ítem 64 / q6} \hfill {\small\bfseries\color{rojo} Nivel 3 · Avanzado}
\end{mdframed}

Una investigadora cualitativa publica un estudio sobre la experiencia de los pacientes en urgencias. Un revisor cuantitativo cuestiona la "validez" y "generalización" del estudio. ¿Cuál es la respuesta metodológicamente más rigurosa de la investigadora?

\begin{enumerate}
  \item[A.] Que su estudio es igualmente generalizable porque usó una muestra de más de 30 participantes
  \item[\textbf{\color{verde}B.}] \textbf{\color{verde} Que en la tradición cualitativa los criterios de rigor son credibilidad, transferibilidad, dependencia y confirmabilidad, que responden a preguntas distintas de las que abordan la validez interna y externa cuantitativas} \hfill $\checkmark$
  \item[C.] Que la ausencia de grupo control en los estudios cualitativos es una limitación aceptable en psicología clínica
  \item[D.] Que los estudios cualitativos no necesitan criterios de rigor porque su objetivo es comprender, no medir
\end{enumerate}

\noindent\textbf{\color{azul}Explicación:} La tradición cualitativa tiene sus propios criterios de rigor (Lincoln y Guba): credibilidad (equivalente a la validez interna), transferibilidad (equivalente a la validez externa), dependencia (equivalente a la fiabilidad) y confirmabilidad (equivalente a la objetividad). Estos criterios reconocen la naturaleza interpretativa del conocimiento cualitativo y no son inferiores a los cuantitativos, sino apropiados para un paradigma diferente.

\noindent\textbf{\color{gris}Por qué son incorrectas las otras opciones:}
\begin{itemize}
  \item {\small A — Incorrecto: el tamaño muestral no garantiza la generalización estadística en estudios cualitativos; los criterios de rigor cualitativos no se basan en el número de participantes.}
  \item {\small C — Incorrecto: la ausencia de grupo control no es relevante en investigación cualitativa porque no se busca establecer causalidad; la crítica no procede.}
  \item {\small D — Incorrecto: afirmar que los estudios cualitativos no necesitan criterios de rigor es incorrecto; tienen criterios propios y bien desarrollados que garantizan la calidad del proceso y los resultados.}
\end{itemize}
\vspace{8pt}\hrule\vspace{10pt}

\begin{mdframed}[backgroundcolor=fondogris,linecolor=rojo,linewidth=1pt,innerleftmargin=10pt,innerrightmargin=10pt,innertopmargin=7pt,innerbottommargin=7pt]
\textbf{Ítem 65 / q7} \hfill {\small\bfseries\color{rojo} Nivel 3 · Avanzado}
\end{mdframed}

Un equipo quiere evaluar un programa de atención a personas mayores en una residencia. El objetivo es tanto comprender la experiencia de los residentes como proponer y probar mejoras con su participación. El proceso se articula en ciclos: diagnóstico → planificación → acción → reflexión. ¿Qué diseño se ajusta a este objetivo?

\begin{enumerate}
  \item[A.] Estudio de casos instrumental, porque el programa es la unidad de análisis
  \item[B.] Diseño cuasiexperimental de series temporales, porque hay medidas antes y después de cada ciclo
  \item[\textbf{\color{verde}C.}] \textbf{\color{verde} Investigación-acción participativa, porque busca simultáneamente generar conocimiento y transformar la realidad con los participantes como co-investigadores} \hfill $\checkmark$
  \item[D.] Diseño longitudinal descriptivo, porque se recogen datos en múltiples momentos temporales
\end{enumerate}

\noindent\textbf{\color{azul}Explicación:} La investigación-acción (Lewin, Kemmis) se define por su doble objetivo: producir conocimiento y transformar la realidad, con los participantes como co-investigadores activos. El proceso cíclico de diagnóstico-planificación-acción-reflexión es su sello metodológico. La variante participativa (IAP) enfatiza el papel central de la comunidad en todo el proceso.

\noindent\textbf{\color{gris}Por qué son incorrectas las otras opciones:}
\begin{itemize}
  \item {\small A — Incorrecto: el estudio de casos instrumental analiza un caso en profundidad para iluminar un fenómeno, pero no tiene el objetivo explícito de transformación ni la estructura cíclica de acción y reflexión.}
  \item {\small B — Incorrecto: el diseño cuasiexperimental de series temporales mide el impacto de una intervención definida externamente; no implica participación de los residentes ni ciclos de reflexión compartida.}
  \item {\small D — Incorrecto: el diseño longitudinal descriptivo mide variables a lo largo del tiempo sin objetivo transformador ni participación de los sujetos en el diseño de la investigación.}
\end{itemize}
\vspace{8pt}\hrule\vspace{10pt}

\begin{mdframed}[backgroundcolor=fondogris,linecolor=rojo,linewidth=1pt,innerleftmargin=10pt,innerrightmargin=10pt,innertopmargin=7pt,innerbottommargin=7pt]
\textbf{Ítem 66 / ql12} \hfill {\small\bfseries\color{rojo} Nivel 3 · Avanzado}
\end{mdframed}

Un equipo de investigación combina en un mismo estudio un cuestionario de bienestar aplicado a 500 docentes (fase cuantitativa) con grupos de discusión posteriores a una submuestra de 20 docentes para explorar en profundidad los factores que explican los resultados del cuestionario (fase cualitativa). ¿Cómo se denomina este diseño integrado y cuál es su lógica?

\begin{enumerate}
  \item[A.] Diseño experimental secuencial, porque el experimento cuantitativo precede al análisis cualitativo
  \item[\textbf{\color{verde}B.}] \textbf{\color{verde} Diseño mixto secuencial explicativo, porque la fase cuantitativa identifica patrones y la fase cualitativa los explica en profundidad} \hfill $\checkmark$
  \item[C.] Diseño cuasiexperimental con validación cualitativa, porque el cuestionario es la VI y los grupos de discusión la VD
  \item[D.] Diseño longitudinal con triangulación, porque las dos fases se repiten a lo largo del tiempo
\end{enumerate}

\noindent\textbf{\color{azul}Explicación:} El diseño mixto secuencial explicativo (Creswell) combina una primera fase cuantitativa para obtener resultados generales con una segunda fase cualitativa que profundiza en los casos, subgrupos o patrones que el cuestionario no puede explicar por sí solo. La integración ocurre en la fase de interpretación, donde los datos cualitativos "explican" los resultados cuantitativos.

\noindent\textbf{\color{gris}Por qué son incorrectas las otras opciones:}
\begin{itemize}
  \item {\small A — Incorrecto: el diseño experimental secuencial implica manipulación experimental en la primera fase; aquí la primera fase es un cuestionario descriptivo sin manipulación.}
  \item {\small C — Incorrecto: el cuestionario no es una VI (no hay manipulación); los grupos de discusión no miden el efecto de ninguna variable manipulada.}
  \item {\small D — Incorrecto: el diseño longitudinal implica medir las mismas variables en distintos momentos del tiempo; aquí hay dos métodos distintos aplicados en dos fases, no la misma variable medida repetidamente.}
\end{itemize}
\vspace{8pt}\hrule\vspace{10pt}

\section{Ex post facto (EPF) · T.11}
{\small\color{gris} 16 ítems · Niveles: L1=9 · L2=5 · L3=2}

\begin{mdframed}[backgroundcolor=fondogris,linecolor=azul,linewidth=1pt,innerleftmargin=10pt,innerrightmargin=10pt,innertopmargin=7pt,innerbottommargin=7pt]
\textbf{Ítem 67 / ep1} \hfill {\small\bfseries\color{azul} Nivel 1 · Básico}
\end{mdframed}

Una investigadora localiza a adultos de 40 años con diagnóstico de diabetes tipo 2 y a adultos sanos de la misma edad. A continuación les pregunta por sus hábitos alimentarios durante la adolescencia, su nivel de actividad física pasada y su historial familiar. El objetivo es identificar qué factores podrían haber contribuido al desarrollo de la diabetes. ¿Qué tipo de diseño describe mejor este estudio?

\begin{enumerate}
  \item[A.] Diseño experimental entre grupos, porque compara dos grupos con y sin diabetes
  \item[B.] Diseño descriptivo por encuestas, porque recoge información mediante cuestionarios
  \item[\textbf{\color{verde}C.}] \textbf{\color{verde} Diseño ex post facto retrospectivo, porque parte del efecto ya ocurrido y busca causas hacia atrás en el tiempo} \hfill $\checkmark$
  \item[D.] Diseño cuasiexperimental de grupos no equivalentes, porque los grupos no se asignaron aleatoriamente
\end{enumerate}

\noindent\textbf{\color{azul}Explicación:} El diseño ex post facto retrospectivo parte de un efecto ya ocurrido (la diabetes) y mira hacia atrás en el tiempo para identificar posibles causas anteriores. El investigador no ha manipulado ninguna variable: los grupos ya difieren en la variable de resultado (diabetes sí/no) y se reconstruyen las exposiciones pasadas.

\noindent\textbf{\color{gris}Por qué son incorrectas las otras opciones:}
\begin{itemize}
  \item {\small A — Incorrecto: el experimento requiere manipulación deliberada de la VI; aquí nadie asignó la diabetes a ningún grupo.}
  \item {\small B — Incorrecto: el diseño descriptivo por encuestas describe una situación actual sin comparar grupos en función de un efecto previo.}
  \item {\small D — Incorrecto: el cuasiexperimental implica que el investigador aplica una intervención aunque no aleatorice; aquí no hay ninguna intervención.}
\end{itemize}
\vspace{8pt}\hrule\vspace{10pt}

\begin{mdframed}[backgroundcolor=fondogris,linecolor=azul,linewidth=1pt,innerleftmargin=10pt,innerrightmargin=10pt,innertopmargin=7pt,innerbottommargin=7pt]
\textbf{Ítem 68 / ep2} \hfill {\small\bfseries\color{azul} Nivel 1 · Básico}
\end{mdframed}

Un equipo identifica hoy a 500 adolescentes de 14 años, la mitad con consumo habitual de alcohol y la mitad sin él. Les siguen durante 10 años registrando si desarrollan problemas de salud mental. No manipulan ninguna variable. ¿Qué tipo de diseño ex post facto es este?

\begin{enumerate}
  \item[A.] Retrospectivo, porque mira hacia atrás para identificar causas del consumo de alcohol
  \item[\textbf{\color{verde}B.}] \textbf{\color{verde} Prospectivo, porque parte de una causa presente (el consumo) y avanza en el tiempo para observar efectos futuros} \hfill $\checkmark$
  \item[C.] Correlacional transversal, porque mide las dos variables en un único momento temporal
  \item[D.] Cuasiexperimental, porque hay un grupo expuesto y un grupo no expuesto a la variable de estudio
\end{enumerate}

\noindent\textbf{\color{azul}Explicación:} El diseño ex post facto prospectivo parte de una condición o exposición ya presente en los participantes (el consumo de alcohol) y los sigue hacia el futuro para observar qué efectos aparecen. La dirección temporal es de la causa al efecto, al contrario que el retrospectivo. No hay manipulación en ninguno de los dos tipos.

\noindent\textbf{\color{gris}Por qué son incorrectas las otras opciones:}
\begin{itemize}
  \item {\small A — Incorrecto: el retrospectivo parte del efecto ya ocurrido y busca causas hacia atrás; aquí el efecto (problemas de salud mental) todavía no ha ocurrido.}
  \item {\small C — Incorrecto: el diseño es longitudinal, no transversal; y la comparación no es correlacional sino entre grupos con y sin exposición.}
  \item {\small D — Incorrecto: el cuasiexperimental implica que el investigador aplica deliberadamente una intervención; aquí el consumo de alcohol es una condición preexistente no manipulada.}
\end{itemize}
\vspace{8pt}\hrule\vspace{10pt}

\begin{mdframed}[backgroundcolor=fondogris,linecolor=azul,linewidth=1pt,innerleftmargin=10pt,innerrightmargin=10pt,innertopmargin=7pt,innerbottommargin=7pt]
\textbf{Ítem 69 / n03} \hfill {\small\bfseries\color{azul} Nivel 1 · Básico}
\end{mdframed}

Una psicóloga selecciona a 50 personas que sufrieron maltrato en la infancia y a 50 personas sin esa historia. Compara sus niveles actuales de autoestima. No ha manipulado nada: el maltrato ya ocurrió en el pasado. ¿Qué diseño es este?

\begin{enumerate}
  \item[A.] Experimento entre grupos, porque compara dos grupos con y sin una condición
  \item[\textbf{\color{verde}B.}] \textbf{\color{verde} Diseño ex post facto retrospectivo, porque parte de grupos que ya difieren en una variable pasada y mide un efecto presente} \hfill $\checkmark$
  \item[C.] Diseño cuasiexperimental, porque no pudo asignar aleatoriamente el maltrato
  \item[D.] Diseño de caso único, porque estudia en profundidad la historia de cada participante
\end{enumerate}

\noindent\textbf{\color{azul}Explicación:} El diseño ex post facto retrospectivo parte de grupos que ya difieren en una condición previa (maltrato sí/no) y mira sus efectos actuales. El investigador no ha manipulado nada; el "tratamiento" (el maltrato) ya ocurrió antes de que comenzara el estudio.

\noindent\textbf{\color{gris}Por qué son incorrectas las otras opciones:}
\begin{itemize}
  \item {\small A — Incorrecto: el experimento requiere que el investigador asigne deliberadamente la VI; nadie asignó el maltrato.}
  \item {\small C — Incorrecto: el cuasiexperimental implica que el investigador sí aplica una intervención aunque no aleatorice; aquí no hay intervención alguna.}
  \item {\small D — Incorrecto: el diseño de caso único estudia a un solo participante con replicación de intervenciones a lo largo del tiempo.}
\end{itemize}
\vspace{8pt}\hrule\vspace{10pt}

\begin{mdframed}[backgroundcolor=fondogris,linecolor=azul,linewidth=1pt,innerleftmargin=10pt,innerrightmargin=10pt,innertopmargin=7pt,innerbottommargin=7pt]
\textbf{Ítem 70 / n07} \hfill {\small\bfseries\color{azul} Nivel 1 · Básico}
\end{mdframed}

Un epidemiólogo identifica hoy a 300 fumadores de 30 años y a 300 no fumadores de la misma edad. Los seguirá durante 20 años para registrar si desarrollan enfermedades respiratorias. No manipula el tabaquismo. ¿Qué diseño es este?

\begin{enumerate}
  \item[A.] Diseño experimental entre grupos, porque compara fumadores y no fumadores
  \item[\textbf{\color{verde}B.}] \textbf{\color{verde} Diseño ex post facto prospectivo, porque parte de una causa presente y sigue a los participantes hacia el futuro} \hfill $\checkmark$
  \item[C.] Diseño descriptivo transversal, porque mide en un único momento temporal
  \item[D.] Diseño cuasiexperimental, porque aplica una intervención (el seguimiento) a los dos grupos
\end{enumerate}

\noindent\textbf{\color{azul}Explicación:} El ex post facto prospectivo parte de una condición ya presente en los participantes (fumar) y avanza en el tiempo para observar los efectos. A diferencia del retrospectivo, aquí el efecto (la enfermedad respiratoria) todavía no ha ocurrido cuando empieza el estudio.

\noindent\textbf{\color{gris}Por qué son incorrectas las otras opciones:}
\begin{itemize}
  \item {\small A — Incorrecto: el experimento requiere que el investigador asigne deliberadamente la VI; nadie asigna el tabaquismo.}
  \item {\small C — Incorrecto: el diseño transversal mide en un único momento; aquí hay seguimiento de 20 años.}
  \item {\small D — Incorrecto: el seguimiento es solo una técnica de recogida de datos, no una intervención experimental; el cuasiexperimental requiere aplicar deliberadamente una condición.}
\end{itemize}
\vspace{8pt}\hrule\vspace{10pt}

\begin{mdframed}[backgroundcolor=fondogris,linecolor=azul,linewidth=1pt,innerleftmargin=10pt,innerrightmargin=10pt,innertopmargin=7pt,innerbottommargin=7pt]
\textbf{Ítem 71 / int02} \hfill {\small\bfseries\color{azul} Nivel 1 · Básico}
\end{mdframed}

Una investigadora estudia si los niños que ven más de dos horas de televisión al día tienen peores hábitos de sueño. Encuesta a 500 familias y analiza la relación entre horas de televisión y calidad del sueño. No cambia los hábitos de ninguna familia. ¿Esta investigadora interviene en las variables que estudia?

\begin{enumerate}
  \item[A.] Sí, porque registra activamente los hábitos de televisión y sueño de las familias
  \item[\textbf{\color{verde}B.}] \textbf{\color{verde} No, porque solo mide variables que ya existen sin modificarlas ni aplicar ningún tratamiento} \hfill $\checkmark$
  \item[C.] Sí, porque selecciona a las familias con más de dos horas de televisión al día
  \item[D.] No, porque usa muestreo aleatorio para seleccionar a los participantes
\end{enumerate}

\noindent\textbf{\color{azul}Explicación:} La ausencia de intervención o manipulación es la característica definitoria de los diseños no experimentales (descriptivos y ex post facto). El investigador mide variables que ya existen en los participantes pero no altera su estado ni aplica ninguna condición experimental.

\noindent\textbf{\color{gris}Por qué son incorrectas las otras opciones:}
\begin{itemize}
  \item {\small A — Incorrecto: registrar o medir no es lo mismo que intervenir; la intervención implica cambiar deliberadamente el estado o la conducta de los participantes.}
  \item {\small C — Incorrecto: seleccionar participantes según una característica es un criterio de inclusión, no una intervención sobre las variables.}
  \item {\small D — Incorrecto: el muestreo aleatorio es una técnica de selección de la muestra, no una intervención sobre las variables estudiadas.}
\end{itemize}
\vspace{8pt}\hrule\vspace{10pt}

\begin{mdframed}[backgroundcolor=fondogris,linecolor=azul,linewidth=1pt,innerleftmargin=10pt,innerrightmargin=10pt,innertopmargin=7pt,innerbottommargin=7pt]
\textbf{Ítem 72 / cov02} \hfill {\small\bfseries\color{azul} Nivel 1 · Básico}
\end{mdframed}

Para poder inferir que la variable A causa el efecto B, según León y Montero se necesitan tres condiciones: covariación, antecesión temporal y descarte de variables de confusión. ¿Cuál de las siguientes situaciones demuestra covariación entre dos variables?

\begin{enumerate}
  \item[A.] El investigador aplica el tratamiento A antes de medir el resultado B
  \item[\textbf{\color{verde}B.}] \textbf{\color{verde} Cuando los valores de A aumentan, los valores de B también aumentan de forma sistemática en los datos} \hfill $\checkmark$
  \item[C.] El investigador descarta que una tercera variable C explique la relación entre A y B
  \item[D.] El investigador asigna aleatoriamente a los participantes a las condiciones A y B
\end{enumerate}

\noindent\textbf{\color{azul}Explicación:} La covariación es la primera condición para inferir causalidad: A y B deben variar juntas de forma sistemática. Si A sube cuando B sube (o baja), hay covariación. Pero la covariación sola no basta para establecer causalidad: también se necesita que A preceda a B en el tiempo y que se descarten explicaciones alternativas.

\noindent\textbf{\color{gris}Por qué son incorrectas las otras opciones:}
\begin{itemize}
  \item {\small A — Incorrecto: aplicar A antes de medir B es la condición de antecesión temporal, no de covariación.}
  \item {\small C — Incorrecto: descartar terceras variables es la condición de control de variables de confusión, no de covariación.}
  \item {\small D — Incorrecto: la asignación aleatoria es una estrategia experimental para controlar variables extrañas; no es la definición de covariación.}
\end{itemize}
\vspace{8pt}\hrule\vspace{10pt}

\begin{mdframed}[backgroundcolor=fondogris,linecolor=azul,linewidth=1pt,innerleftmargin=10pt,innerrightmargin=10pt,innertopmargin=7pt,innerbottommargin=7pt]
\textbf{Ítem 73 / ant01} \hfill {\small\bfseries\color{azul} Nivel 1 · Básico}
\end{mdframed}

Para demostrar que el consumo de alcohol en la adolescencia causa problemas de memoria en la adultez, es necesario demostrar que el consumo de alcohol ocurrió ANTES de los problemas de memoria. Esta condición se llama:

\begin{enumerate}
  \item[A.] Covariación, porque alcohol y memoria deben cambiar juntos de forma sistemática
  \item[\textbf{\color{verde}B.}] \textbf{\color{verde} Antecesión temporal, porque la causa debe preceder al efecto en el tiempo} \hfill $\checkmark$
  \item[C.] Descarte de terceras variables, porque hay que eliminar otras explicaciones
  \item[D.] Validez interna, porque hay que controlar las variables extrañas dentro del estudio
\end{enumerate}

\noindent\textbf{\color{azul}Explicación:} La antecesión temporal es la segunda condición necesaria para inferir causalidad: la causa (A) debe ocurrir antes que el efecto (B) en el tiempo. En los diseños retrospectivos esta condición es especialmente difícil de establecer porque se reconstruye el pasado desde el presente, con los consiguientes sesgos de memoria.

\noindent\textbf{\color{gris}Por qué son incorrectas las otras opciones:}
\begin{itemize}
  \item {\small A — Incorrecto: la covariación se refiere a que las variables varíen juntas de forma sistemática; no tiene que ver con el orden temporal.}
  \item {\small C — Incorrecto: el descarte de terceras variables es la tercera condición para inferir causalidad; la segunda (la que se pregunta aquí) es la antecesión temporal.}
  \item {\small D — Incorrecto: la validez interna es la propiedad global del diseño que permite atribuir los cambios en la VD a la VI; no es el nombre de la condición temporal.}
\end{itemize}
\vspace{8pt}\hrule\vspace{10pt}

\begin{mdframed}[backgroundcolor=fondogris,linecolor=azul,linewidth=1pt,innerleftmargin=10pt,innerrightmargin=10pt,innertopmargin=7pt,innerbottommargin=7pt]
\textbf{Ítem 74 / n10} \hfill {\small\bfseries\color{azul} Nivel 1 · Básico}
\end{mdframed}

Una investigadora encuentra a adultos mayores con deterioro cognitivo leve y a adultos mayores sin deterioro. Les pregunta sobre sus hábitos de lectura y actividad intelectual durante los años previos para identificar factores de protección. ¿Qué diseño es este?

\begin{enumerate}
  \item[A.] Diseño experimental, porque compara dos grupos con y sin deterioro
  \item[\textbf{\color{verde}B.}] \textbf{\color{verde} Diseño ex post facto retrospectivo, porque parte del efecto (el deterioro) y busca causas pasadas} \hfill $\checkmark$
  \item[C.] Diseño prospectivo, porque analiza variables que pueden predecir el deterioro en el futuro
  \item[D.] Diseño descriptivo por encuestas, porque usa cuestionarios para recoger datos
\end{enumerate}

\noindent\textbf{\color{azul}Explicación:} El diseño retrospectivo parte de un efecto ya ocurrido (el deterioro cognitivo) y mira hacia atrás para identificar posibles causas (los hábitos intelectuales previos). La clave es la dirección temporal: del efecto presente a las causas pasadas.

\noindent\textbf{\color{gris}Por qué son incorrectas las otras opciones:}
\begin{itemize}
  \item {\small A — Incorrecto: el experimento requiere que el investigador asigne deliberadamente el deterioro cognitivo; esto sería imposible e inmoral.}
  \item {\small C — Incorrecto: el diseño prospectivo parte de una causa presente y mira hacia el futuro; aquí se mira hacia el pasado, lo que lo hace retrospectivo.}
  \item {\small D — Incorrecto: aunque se usan cuestionarios, el propósito y la estructura del diseño son los del ex post facto retrospectivo, no los del meramente descriptivo.}
\end{itemize}
\vspace{8pt}\hrule\vspace{10pt}

\begin{mdframed}[backgroundcolor=fondogris,linecolor=azul,linewidth=1pt,innerleftmargin=10pt,innerrightmargin=10pt,innertopmargin=7pt,innerbottommargin=7pt]
\textbf{Ítem 75 / ep06} \hfill {\small\bfseries\color{azul} Nivel 1 · Básico}
\end{mdframed}

Una investigadora quiere estudiar si los niños que recibieron lactancia materna tienen un cociente intelectual más alto a los 10 años que los que no la recibieron. Identifica en registros hospitalarios a niños de 10 años con y sin lactancia registrada al nacer. No ha manipulado nada. ¿Qué tipo de diseño es este?

\begin{enumerate}
  \item[A.] Diseño experimental entre grupos, porque compara dos grupos con y sin una condición
  \item[\textbf{\color{verde}B.}] \textbf{\color{verde} Diseño ex post facto retrospectivo, porque parte de grupos que ya difieren en una variable pasada (lactancia) y mide un efecto presente (CI)} \hfill $\checkmark$
  \item[C.] Diseño cuasiexperimental, porque aplica la lactancia materna como intervención a un grupo
  \item[D.] Diseño descriptivo transversal, porque mide el CI en un único momento temporal
\end{enumerate}

\noindent\textbf{\color{azul}Explicación:} El diseño es retrospectivo porque la variable de exposición (lactancia materna) ya ocurrió antes de que comenzara el estudio; el investigador la identifica en registros pasados. El efecto (CI actual) se mide en el presente. La investigadora no ha manipulado ni asignado la lactancia: se trata de un dato preexistente.

\noindent\textbf{\color{gris}Por qué son incorrectas las otras opciones:}
\begin{itemize}
  \item {\small A — Incorrecto: en el experimento el investigador asigna deliberadamente la lactancia; asignar lactancia materna sería éticamente inviable.}
  \item {\small C — Incorrecto: el cuasiexperimental implica que el investigador aplica una intervención; la lactancia ya ocurrió y nadie la aplicó como parte del estudio.}
  \item {\small D — Incorrecto: aunque el CI se mide en un único momento, el diseño es más complejo que uno meramente descriptivo porque compara grupos definidos por una exposición pasada.}
\end{itemize}
\vspace{8pt}\hrule\vspace{10pt}

\begin{mdframed}[backgroundcolor=fondogris,linecolor=ambar,linewidth=1pt,innerleftmargin=10pt,innerrightmargin=10pt,innertopmargin=7pt,innerbottommargin=7pt]
\textbf{Ítem 76 / ep3} \hfill {\small\bfseries\color{ambar} Nivel 2 · Intermedio}
\end{mdframed}

Una investigadora compara el rendimiento académico de 200 estudiantes con TDAH y 200 sin TDAH, emparejados por edad y sexo. No ha manipulado el diagnóstico. Sus resultados muestran que el grupo con TDAH rinde significativamente peor (p<.001). ¿Qué conclusión está justificada y cuál no, según León y Montero?

\begin{enumerate}
  \item[A.] Está justificado afirmar que el TDAH causa el bajo rendimiento; no está justificado afirmar que el tratamiento mejora el rendimiento
  \item[\textbf{\color{verde}B.}] \textbf{\color{verde} Está justificado afirmar que existe una asociación entre TDAH y rendimiento académico; no está justificado inferir causalidad porque no hubo manipulación ni aleatorización} \hfill $\checkmark$
  \item[C.] Está justificado generalizar el resultado a toda la población clínica sin restricciones muestrales
  \item[D.] Está justificado afirmar causalidad porque el TDAH es anterior en el tiempo al bajo rendimiento observado en el estudio
\end{enumerate}

\noindent\textbf{\color{azul}Explicación:} En el diseño ex post facto, la ausencia de manipulación y de aleatorización impide inferir causalidad, aunque la asociación sea estadísticamente significativa. El hecho de que el TDAH sea temporalmente anterior no basta para establecer causalidad: podrían existir terceras variables (nivel socioeconómico, calidad del entorno familiar) que expliquen tanto el diagnóstico como el rendimiento.

\noindent\textbf{\color{gris}Por qué son incorrectas las otras opciones:}
\begin{itemize}
  \item {\small A — Incorrecto: afirmar causalidad directa del TDAH al rendimiento excede lo que el diseño ex post facto permite; además, no se ha evaluado ningún tratamiento.}
  \item {\small C — Incorrecto: la generalización depende de que la muestra sea representativa de la población clínica objetivo, lo que no se garantiza automáticamente.}
  \item {\small D — Incorrecto: la precedencia temporal es condición necesaria pero no suficiente para establecer causalidad; también se requiere descartar variables de confusión y demostrar covariación controlada.}
\end{itemize}
\vspace{8pt}\hrule\vspace{10pt}

\begin{mdframed}[backgroundcolor=fondogris,linecolor=ambar,linewidth=1pt,innerleftmargin=10pt,innerrightmargin=10pt,innertopmargin=7pt,innerbottommargin=7pt]
\textbf{Ítem 77 / ep4} \hfill {\small\bfseries\color{ambar} Nivel 2 · Intermedio}
\end{mdframed}

Dos investigadores estudian la relación entre el tabaquismo en la adolescencia y el desarrollo de enfermedades cardiovasculares. El primero localiza a pacientes adultos con enfermedad cardiovascular y les pregunta si fumaron de jóvenes. El segundo identifica hoy a adolescentes fumadores y no fumadores y los seguirá 25 años. ¿Cuál de los dos diseños tiene mayor validez interna y por qué?

\begin{enumerate}
  \item[A.] El primero, porque recoge datos en pacientes ya diagnosticados y puede controlar mejor las variables extrañas
  \item[\textbf{\color{verde}B.}] \textbf{\color{verde} El segundo, porque mide el tabaquismo antes de que ocurra el efecto, eliminando el sesgo de recuerdo retrospectivo} \hfill $\checkmark$
  \item[C.] Ambos tienen la misma validez interna porque ninguno manipula la VI ni aleatoriza a los participantes
  \item[D.] El primero, porque el tamaño muestral de pacientes con enfermedad cardiovascular es mayor
\end{enumerate}

\noindent\textbf{\color{azul}Explicación:} El diseño prospectivo tiene mayor validez interna que el retrospectivo porque mide la causa antes de que ocurra el efecto, evitando el principal sesgo del retrospectivo: la distorsión del recuerdo (recordar haber fumado más o menos de lo que realmente se fumó, influido por el diagnóstico actual). No obstante, ninguno de los dos permite inferir causalidad estricta.

\noindent\textbf{\color{gris}Por qué son incorrectas las otras opciones:}
\begin{itemize}
  \item {\small A — Incorrecto: el diseño retrospectivo está especialmente expuesto al sesgo de memoria; los pacientes pueden recordar su tabaquismo adolescente de forma distorsionada por su estado actual de salud.}
  \item {\small C — Incorrecto: aunque ambos comparten la limitación de no poder establecer causalidad, difieren en la calidad de la medición de la exposición (tabaquismo), siendo mejor el prospectivo.}
  \item {\small D — Incorrecto: el tamaño muestral afecta a la potencia estadística, pero no a la validez interna que es un asunto de diseño y control de sesgos.}
\end{itemize}
\vspace{8pt}\hrule\vspace{10pt}

\begin{mdframed}[backgroundcolor=fondogris,linecolor=ambar,linewidth=1pt,innerleftmargin=10pt,innerrightmargin=10pt,innertopmargin=7pt,innerbottommargin=7pt]
\textbf{Ítem 78 / ant02} \hfill {\small\bfseries\color{ambar} Nivel 2 · Intermedio}
\end{mdframed}

Una investigadora usa un diseño transversal para estudiar la relación entre el nivel de estrés laboral y las enfermedades cardiovasculares: mide ambas variables en el mismo momento. Encuentra una correlación significativa. Su colega señala que hay un problema con la inferencia causal. ¿Cuál es ese problema?

\begin{enumerate}
  \item[A.] La muestra es demasiado pequeña para encontrar correlaciones significativas
  \item[\textbf{\color{verde}B.}] \textbf{\color{verde} Al medir ambas variables en el mismo momento no puede establecerse cuál precedió a la otra en el tiempo, lo que impide cumplir la condición de antecesión temporal} \hfill $\checkmark$
  \item[C.] La correlación solo puede calcularse si las variables son experimentalmente manipuladas
  \item[D.] El diseño transversal no permite calcular correlaciones entre variables continuas
\end{enumerate}

\noindent\textbf{\color{azul}Explicación:} El diseño transversal mide todas las variables en un único momento temporal. Esto impide saber si el estrés precedió a las enfermedades cardiovasculares o si, al contrario, tener enfermedad cardiovascular genera estrés. Sin antecesión temporal clara, no puede inferirse causalidad aunque la correlación sea alta.

\noindent\textbf{\color{gris}Por qué son incorrectas las otras opciones:}
\begin{itemize}
  \item {\small A — Incorrecto: la pregunta habla de una correlación significativa, lo que implica que el tamaño muestral fue suficiente.}
  \item {\small C — Incorrecto: la correlación puede calcularse en cualquier diseño con dos variables cuantitativas; la manipulación experimental no es un requisito para calcularla.}
  \item {\small D — Incorrecto: el diseño transversal permite perfectamente calcular correlaciones entre variables continuas; el problema es de inferencia causal, no de posibilidad estadística.}
\end{itemize}
\vspace{8pt}\hrule\vspace{10pt}

\begin{mdframed}[backgroundcolor=fondogris,linecolor=ambar,linewidth=1pt,innerleftmargin=10pt,innerrightmargin=10pt,innertopmargin=7pt,innerbottommargin=7pt]
\textbf{Ítem 79 / ep07} \hfill {\small\bfseries\color{ambar} Nivel 2 · Intermedio}
\end{mdframed}

Un equipo de investigadores compara el rendimiento cognitivo de un grupo de adultos que practicaron deporte de forma regular durante su adolescencia y un grupo que no lo practicó, emparejando ambos grupos por edad, sexo y nivel educativo actual. Encuentran que los deportistas rinden mejor. Un colega argumenta que el emparejamiento no resuelve todos los problemas de validez. ¿Cuál es el argumento más sólido?

\begin{enumerate}
  \item[A.] El emparejamiento hace que la muestra sea demasiado pequeña para tener potencia estadística suficiente
  \item[\textbf{\color{verde}B.}] \textbf{\color{verde} El emparejamiento solo controla las variables en las que se empareja; pueden existir otras variables no controladas (p. ej., motivación, hábitos alimentarios) que expliquen tanto la práctica deportiva como el rendimiento cognitivo} \hfill $\checkmark$
  \item[C.] El emparejamiento no es válido porque las variables de emparejamiento (edad, sexo, nivel educativo) no están relacionadas con el rendimiento cognitivo
  \item[D.] El emparejamiento convierte el diseño en experimental, por lo que sí puede inferirse causalidad
\end{enumerate}

\noindent\textbf{\color{azul}Explicación:} El emparejamiento es una técnica útil para controlar variables conocidas, pero no puede garantizar la equivalencia en todas las variables relevantes. Las variables no medidas ni controladas (confusores residuales) pueden seguir explicando la relación observada. Esta limitación es estructural en todos los diseños ex post facto y solo la aleatorización puede resolverla.

\noindent\textbf{\color{gris}Por qué son incorrectas las otras opciones:}
\begin{itemize}
  \item {\small A — Incorrecto: el emparejamiento no implica necesariamente una muestra pequeña; puede realizarse con muestras grandes.}
  \item {\small C — Incorrecto: la edad, el sexo y el nivel educativo sí están relacionados con el rendimiento cognitivo; precisamente por eso se usan como variables de emparejamiento.}
  \item {\small D — Incorrecto: el emparejamiento es una técnica de control estadístico, no de asignación aleatoria; no convierte el diseño en experimental ni permite inferir causalidad.}
\end{itemize}
\vspace{8pt}\hrule\vspace{10pt}

\begin{mdframed}[backgroundcolor=fondogris,linecolor=ambar,linewidth=1pt,innerleftmargin=10pt,innerrightmargin=10pt,innertopmargin=7pt,innerbottommargin=7pt]
\textbf{Ítem 80 / ep08} \hfill {\small\bfseries\color{ambar} Nivel 2 · Intermedio}
\end{mdframed}

Una investigadora realiza un estudio prospectivo: identifica hoy a 400 trabajadores de oficina con elevado sedentarismo y a 400 con actividad física regular, y los sigue durante 5 años midiendo su salud cardiovascular anualmente. ¿Cuál es la ventaja de este diseño frente a un estudio retrospectivo sobre el mismo tema?

\begin{enumerate}
  \item[A.] Que permite asignar aleatoriamente a los trabajadores a la condición de sedentarismo o actividad física
  \item[\textbf{\color{verde}B.}] \textbf{\color{verde} Que la medición de la causa (el sedentarismo actual) ocurre antes del efecto (los problemas cardiovasculares futuros), evitando el sesgo de recuerdo retrospectivo y estableciendo la antecesión temporal con más solidez} \hfill $\checkmark$
  \item[C.] Que al tener más años de seguimiento la muestra se vuelve más representativa de la población general
  \item[D.] Que elimina completamente la posibilidad de que terceras variables expliquen la relación entre sedentarismo y salud cardiovascular
\end{enumerate}

\noindent\textbf{\color{azul}Explicación:} El diseño prospectivo mide la exposición (sedentarismo) antes de que ocurra el efecto (problemas cardiovasculares), lo que establece la antecesión temporal con más claridad y evita el sesgo de recuerdo (que en el retrospectivo contamina la reconstrucción de hábitos pasados). Sin embargo, no elimina el problema de las terceras variables ni convierte el diseño en experimental.

\noindent\textbf{\color{gris}Por qué son incorrectas las otras opciones:}
\begin{itemize}
  \item {\small A — Incorrecto: la asignación aleatoria es característica del experimento, no del estudio prospectivo; asignar aleatoriamente a las personas a ser sedentarias sería imposible e inmoral.}
  \item {\small C — Incorrecto: el tiempo de seguimiento no afecta directamente a la representatividad muestral; la representatividad depende del proceso de selección de la muestra, no de la duración.}
  \item {\small D — Incorrecto: ni el diseño prospectivo ni ningún otro diseño observacional elimina completamente el problema de las terceras variables; solo la aleatorización puede hacerlo.}
\end{itemize}
\vspace{8pt}\hrule\vspace{10pt}

\begin{mdframed}[backgroundcolor=fondogris,linecolor=rojo,linewidth=1pt,innerleftmargin=10pt,innerrightmargin=10pt,innertopmargin=7pt,innerbottommargin=7pt]
\textbf{Ítem 81 / ep5} \hfill {\small\bfseries\color{rojo} Nivel 3 · Avanzado}
\end{mdframed}

Un investigador usa un diseño ex post facto para comparar el bienestar psicológico de personas que migraron en los últimos cinco años y personas que no migraron, emparejadas por edad, sexo y nivel educativo. Concluye que la migración reduce el bienestar. Un colega argumenta que esta conclusión no es válida. ¿Cuál es el argumento metodológico más sólido del colega?

\begin{enumerate}
  \item[A.] La muestra de migrantes es difícil de reclutar, lo que reduce la validez externa del estudio
  \item[\textbf{\color{verde}B.}] \textbf{\color{verde} El emparejamiento por edad, sexo y nivel educativo no controla todas las variables de confusión posibles (p. ej. personalidad, red de apoyo social previa, razón de la migración) que podrían explicar las diferencias en bienestar} \hfill $\checkmark$
  \item[C.] El bienestar psicológico no puede medirse con fiabilidad suficiente mediante cuestionarios de autoinforme
  \item[D.] El diseño ex post facto es adecuado solo para estudios retrospectivos, no para variables con efectos a largo plazo como la migración
\end{enumerate}

\noindent\textbf{\color{azul}Explicación:} En el diseño ex post facto, el emparejamiento controla solo las variables en las que se empareja, pero no puede garantizar la equivalencia en todas las variables relevantes. Variables como la personalidad resiliente, las razones de la migración (voluntaria vs. forzada) o la calidad de la red social previa podrían diferir entre grupos y explicar las diferencias en bienestar mejor que la migración en sí.

\noindent\textbf{\color{gris}Por qué son incorrectas las otras opciones:}
\begin{itemize}
  \item {\small A — Incorrecto: la dificultad de reclutamiento afecta a la validez externa (generalización), no a la validez interna que es lo que se cuestiona aquí.}
  \item {\small C — Incorrecto: aunque los autoinformes tienen limitaciones, existen escalas de bienestar con fiabilidad y validez bien establecidas; esta no es la crítica central al diseño.}
  \item {\small D — Incorrecto: el diseño ex post facto puede ser retrospectivo o prospectivo; la duración de los efectos no es un criterio para su aplicabilidad.}
\end{itemize}
\vspace{8pt}\hrule\vspace{10pt}

\begin{mdframed}[backgroundcolor=fondogris,linecolor=rojo,linewidth=1pt,innerleftmargin=10pt,innerrightmargin=10pt,innertopmargin=7pt,innerbottommargin=7pt]
\textbf{Ítem 82 / ep09} \hfill {\small\bfseries\color{rojo} Nivel 3 · Avanzado}
\end{mdframed}

Un investigador quiere estudiar si la pobreza en la infancia predice problemas de salud mental en la adultez. Diseña un estudio en el que identifica a adultos de 35 años con diagnóstico de trastorno de ansiedad o depresión y a adultos sin diagnóstico, y les pregunta retrospectivamente por la situación económica de su familia durante su infancia. Un metodólogo le señala tres problemas fundamentales de este diseño. ¿Cuál de los siguientes NO es un problema metodológico real de este diseño?

\begin{enumerate}
  \item[A.] El sesgo de memoria: los adultos con trastornos pueden recordar su infancia de forma más negativa que la realidad
  \item[B.] La ausencia de antecesión temporal establecida objetivamente: la pobreza se reconstruye retrospectivamente y puede estar contaminada por el estado mental actual
  \item[\textbf{\color{verde}C.}] \textbf{\color{verde} La ausencia de grupo control: el diseño compara adultos con y sin diagnóstico, por lo que sí existe un grupo de referencia} \hfill $\checkmark$
  \item[D.] La imposibilidad de descartar terceras variables: otros factores de la infancia (trauma, falta de apoyo emocional) pueden covariar con la pobreza y ser los verdaderos causantes de los trastornos
\end{enumerate}

\noindent\textbf{\color{azul}Explicación:} El diseño sí incluye un grupo de referencia (adultos sin diagnóstico), así que la ausencia de grupo de comparación NO es un problema real de este diseño. Los problemas genuinos son: el sesgo de recuerdo retrospectivo, la dificultad para establecer la antecesión temporal de forma objetiva y la imposibilidad de controlar todas las variables de confusión propias del entorno de pobreza.

\noindent\textbf{\color{gris}Por qué son incorrectas las otras opciones:}
\begin{itemize}
  \item {\small A — Incorrecto (sí es un problema real): el sesgo de memoria o recuerdo es la amenaza más característica de los diseños retrospectivos; las personas con trastornos ansiosos o depresivos pueden recordar su pasado de forma más negativa (memoria congruente con el estado de ánimo).}
  \item {\small B — Incorrecto (sí es un problema real): aunque la pobreza es anterior en el tiempo al diagnóstico adulto, la medición retrospectiva es subjetiva y puede estar distorsionada, lo que debilita la condición de antecesión temporal objetivamente establecida.}
  \item {\small D — Incorrecto (sí es un problema real): la pobreza covaría con múltiples factores adversos (traumas, negligencia, escaso apoyo educativo) que también predicen trastornos mentales; separarlos estadísticamente sin aleatorización es muy difícil.}
\end{itemize}
\vspace{8pt}\hrule\vspace{10pt}

\section{Experimental entre grupos (EXP-EG) · T.6}
{\small\color{gris} 18 ítems · Niveles: L1=8 · L2=7 · L3=3}

\begin{mdframed}[backgroundcolor=fondogris,linecolor=azul,linewidth=1pt,innerleftmargin=10pt,innerrightmargin=10pt,innertopmargin=7pt,innerbottommargin=7pt]
\textbf{Ítem 83 / e1} \hfill {\small\bfseries\color{azul} Nivel 1 · Básico}
\end{mdframed}

Una investigadora asigna al azar a 90 participantes en tres grupos: el primero recibe terapia cognitivo-conductual, el segundo recibe mindfulness y el tercero es un grupo control en lista de espera. Tras 12 semanas mide los síntomas de ansiedad en los tres grupos. ¿Cuál de las siguientes afirmaciones describe mejor este diseño?

\begin{enumerate}
  \item[A.] Es un diseño ex post facto porque los grupos ya existían antes del estudio
  \item[\textbf{\color{verde}B.}] \textbf{\color{verde} Es un experimento entre grupos porque hay asignación aleatoria, manipulación de la VI y medición de la VD} \hfill $\checkmark$
  \item[C.] Es un diseño cuasiexperimental porque compara más de dos grupos
  \item[D.] Es un diseño factorial porque estudia el efecto de dos tipos de terapia
\end{enumerate}

\noindent\textbf{\color{azul}Explicación:} El experimento entre grupos (between-subjects) se define por tres características: manipulación deliberada de la variable independiente, asignación aleatoria a las condiciones y medición de la variable dependiente para comparar grupos. La asignación aleatoria garantiza la equivalencia inicial y permite atribuir causalmente las diferencias a la VI.

\noindent\textbf{\color{gris}Por qué son incorrectas las otras opciones:}
\begin{itemize}
  \item {\small A — Incorrecto: en el ex post facto los grupos difieren en una variable que ya existía; aquí los grupos se crean aleatoriamente asignando tratamientos.}
  \item {\small C — Incorrecto: el criterio del cuasiexperimental es la falta de aleatorización, no el número de grupos.}
  \item {\small D — Incorrecto: el diseño factorial requiere cruzar dos o más VI; aquí hay una sola VI (tipo de tratamiento) con tres niveles.}
\end{itemize}
\vspace{8pt}\hrule\vspace{10pt}

\begin{mdframed}[backgroundcolor=fondogris,linecolor=azul,linewidth=1pt,innerleftmargin=10pt,innerrightmargin=10pt,innertopmargin=7pt,innerbottommargin=7pt]
\textbf{Ítem 84 / e2} \hfill {\small\bfseries\color{azul} Nivel 1 · Básico}
\end{mdframed}

Un investigador quiere estudiar si el color del fondo de una pantalla (blanco, azul o verde) afecta la velocidad de lectura. Asigna aleatoriamente a 60 participantes (20 por condición) y mide el número de palabras por minuto. ¿Cuál es la variable dependiente en este diseño?

\begin{enumerate}
  \item[A.] El color del fondo de pantalla
  \item[B.] La asignación aleatoria de los participantes
  \item[\textbf{\color{verde}C.}] \textbf{\color{verde} La velocidad de lectura en palabras por minuto} \hfill $\checkmark$
  \item[D.] El número de participantes por condición
\end{enumerate}

\noindent\textbf{\color{azul}Explicación:} La variable dependiente (VD) es la que se mide para observar el efecto de la manipulación experimental. En este caso, el investigador manipula el color del fondo (VI) y mide la velocidad de lectura (VD). La asignación aleatoria es la técnica de control, no una variable del diseño.

\noindent\textbf{\color{gris}Por qué son incorrectas las otras opciones:}
\begin{itemize}
  \item {\small A — Incorrecto: el color del fondo es la variable independiente: es lo que el investigador manipula deliberadamente.}
  \item {\small B — Incorrecto: la asignación aleatoria es el procedimiento de control experimental, no una variable del diseño.}
  \item {\small D — Incorrecto: el número de participantes por condición es un parámetro del diseño, no lo que se mide.}
\end{itemize}
\vspace{8pt}\hrule\vspace{10pt}

\begin{mdframed}[backgroundcolor=fondogris,linecolor=azul,linewidth=1pt,innerleftmargin=10pt,innerrightmargin=10pt,innertopmargin=7pt,innerbottommargin=7pt]
\textbf{Ítem 85 / n01} \hfill {\small\bfseries\color{azul} Nivel 1 · Básico}
\end{mdframed}

Una investigadora divide al azar a 40 personas con fobia a los perros en dos grupos. El grupo A recibe terapia de exposición gradual durante cuatro semanas; el grupo B no recibe tratamiento. Al finalizar, mide el nivel de miedo en ambos grupos. ¿Qué diseño es este?

\begin{enumerate}
  \item[A.] Diseño ex post facto, porque la fobia ya existía antes del estudio
  \item[\textbf{\color{verde}B.}] \textbf{\color{verde} Diseño experimental entre grupos, porque hay aleatorización, manipulación de la terapia y comparación de grupos} \hfill $\checkmark$
  \item[C.] Diseño cuasiexperimental, porque los grupos no podían formarse completamente al azar
  \item[D.] Diseño descriptivo por encuestas, porque mide el miedo mediante cuestionarios
\end{enumerate}

\noindent\textbf{\color{azul}Explicación:} Los tres ingredientes del experimento verdadero están presentes: manipulación de la VI (recibir o no la terapia), asignación aleatoria a los grupos y medición de la VD (nivel de miedo). Esto define al diseño experimental entre grupos o between-subjects.

\noindent\textbf{\color{gris}Por qué son incorrectas las otras opciones:}
\begin{itemize}
  \item {\small A — Incorrecto: en el ex post facto el investigador no aplica ningún tratamiento; aquí la investigadora sí asigna deliberadamente la terapia.}
  \item {\small C — Incorrecto: el cuasiexperimental se define por la ausencia de aleatorización; aquí se aleatoriza explícitamente.}
  \item {\small D — Incorrecto: el diseño descriptivo no manipula variables; aquí hay una intervención activa (la terapia de exposición).}
\end{itemize}
\vspace{8pt}\hrule\vspace{10pt}

\begin{mdframed}[backgroundcolor=fondogris,linecolor=azul,linewidth=1pt,innerleftmargin=10pt,innerrightmargin=10pt,innertopmargin=7pt,innerbottommargin=7pt]
\textbf{Ítem 86 / v01} \hfill {\small\bfseries\color{azul} Nivel 1 · Básico}
\end{mdframed}

Un experimento sobre el efecto del estrés en la memoria se realiza en un laboratorio con estudiantes universitarios de psicología que se ofrecen voluntariamente. Los resultados muestran que el estrés perjudica significativamente el recuerdo. El autor afirma que estos resultados se aplican a toda la población adulta. ¿Qué problema metodológico tiene esta afirmación?

\begin{enumerate}
  \item[A.] Validez interna débil, porque el estrés no estaba bien controlado en el laboratorio
  \item[\textbf{\color{verde}B.}] \textbf{\color{verde} Validez externa limitada, porque la muestra de estudiantes voluntarios no es representativa de toda la población adulta} \hfill $\checkmark$
  \item[C.] Ausencia de asignación aleatoria, lo que convierte el diseño en cuasiexperimental
  \item[D.] El experimento no tiene grupo control, por lo que no puede establecerse el efecto del estrés
\end{enumerate}

\noindent\textbf{\color{azul}Explicación:} La validez externa se refiere a la capacidad de generalizar los resultados a otras personas, contextos y momentos. Usar estudiantes universitarios voluntarios (muestra de conveniencia) limita la generalización porque pueden diferir de la población general en características relevantes para el rendimiento de memoria bajo estrés.

\noindent\textbf{\color{gris}Por qué son incorrectas las otras opciones:}
\begin{itemize}
  \item {\small A — Incorrecto: si el estrés estaba bien manipulado en el laboratorio, la validez interna puede ser alta; el problema es la generalización, no el control interno.}
  \item {\small C — Incorrecto: en un experimento de laboratorio sí puede haber asignación aleatoria aunque la muestra no sea representativa de la población.}
  \item {\small D — Incorrecto: el texto no dice que falta grupo control; el problema señalado es la representatividad de la muestra para generalizar.}
\end{itemize}
\vspace{8pt}\hrule\vspace{10pt}

\begin{mdframed}[backgroundcolor=fondogris,linecolor=azul,linewidth=1pt,innerleftmargin=10pt,innerrightmargin=10pt,innertopmargin=7pt,innerbottommargin=7pt]
\textbf{Ítem 87 / int01} \hfill {\small\bfseries\color{azul} Nivel 1 · Básico}
\end{mdframed}

¿Cuál de los siguientes estudios incluye una intervención activa por parte del investigador?

\begin{enumerate}
  \item[A.] Un investigador observa cuántas veces al día los empleados de una oficina consultan su móvil
  \item[\textbf{\color{verde}B.}] \textbf{\color{verde} Un investigador aplica un programa de mindfulness a un grupo y compara su estrés con un grupo que no lo recibe} \hfill $\checkmark$
  \item[C.] Un investigador pregunta a estudiantes cuántas horas estudian y qué nota han sacado
  \item[D.] Un investigador revisa los expedientes académicos de los últimos cinco años para ver la evolución del rendimiento
\end{enumerate}

\noindent\textbf{\color{azul}Explicación:} Una intervención activa significa que el investigador aplica deliberadamente una condición experimental (el programa de mindfulness) a los participantes. Esto es lo que distingue a los diseños experimentales y cuasiexperimentales de los descriptivos y ex post facto, donde el investigador solo observa o mide variables existentes.

\noindent\textbf{\color{gris}Por qué son incorrectas las otras opciones:}
\begin{itemize}
  \item {\small A — Incorrecto: observar cuántas veces se consulta el móvil es una observación sistemática sin intervención alguna.}
  \item {\small C — Incorrecto: preguntar sobre horas de estudio y notas es un diseño descriptivo correlacional, sin ninguna intervención.}
  \item {\small D — Incorrecto: revisar expedientes históricos es un diseño ex post facto que no implica ninguna intervención activa.}
\end{itemize}
\vspace{8pt}\hrule\vspace{10pt}

\begin{mdframed}[backgroundcolor=fondogris,linecolor=azul,linewidth=1pt,innerleftmargin=10pt,innerrightmargin=10pt,innertopmargin=7pt,innerbottommargin=7pt]
\textbf{Ítem 88 / cu01} \hfill {\small\bfseries\color{azul} Nivel 1 · Básico}
\end{mdframed}

Una psicóloga trabaja con un único niño con autismo que tiene conductas autolesivas frecuentes. Primero observa la conducta durante dos semanas sin hacer nada (fase A). Luego aplica una intervención conductual durante dos semanas (fase B). Después retira la intervención otras dos semanas (fase A). Finalmente la reaplicasun dos semanas más (fase B). ¿Qué tipo de diseño es este?

\begin{enumerate}
  \item[A.] Diseño descriptivo por observación sistemática, porque solo observa la conducta sin manipular
  \item[\textbf{\color{verde}B.}] \textbf{\color{verde} Diseño de caso único A-B-A-B, porque alterna sistemáticamente la presencia y ausencia de la intervención en un solo participante} \hfill $\checkmark$
  \item[C.] Diseño cuasiexperimental de series temporales, porque hay múltiples medidas antes y después
  \item[D.] Diseño ex post facto retrospectivo, porque analiza conductas pasadas del niño
\end{enumerate}

\noindent\textbf{\color{azul}Explicación:} El diseño de caso único A-B-A-B (también llamado de reversión o de retirada) alterna fases de línea de base (A) con fases de intervención (B) en un solo participante. Si la conducta mejora solo en las fases B y empeora al volver a A, se puede inferir que la intervención causa el cambio. Es el diseño más usado en modificación de conducta aplicada.

\noindent\textbf{\color{gris}Por qué son incorrectas las otras opciones:}
\begin{itemize}
  \item {\small A — Incorrecto: la psicóloga no solo observa, también aplica y retira activamente una intervención; hay manipulación deliberada de la VI.}
  \item {\small C — Incorrecto: el diseño de series temporales se aplica a grupos y no incluye la retirada sistemática de la intervención para demostrar causalidad.}
  \item {\small D — Incorrecto: en el ex post facto no hay manipulación activa; aquí la psicóloga aplica y retira la intervención de forma controlada.}
\end{itemize}
\vspace{8pt}\hrule\vspace{10pt}

\begin{mdframed}[backgroundcolor=fondogris,linecolor=azul,linewidth=1pt,innerleftmargin=10pt,innerrightmargin=10pt,innertopmargin=7pt,innerbottommargin=7pt]
\textbf{Ítem 89 / cu02} \hfill {\small\bfseries\color{azul} Nivel 1 · Básico}
\end{mdframed}

En un diseño de caso único, un terapeuta observa que la conducta agresiva de su paciente disminuye durante la fase de intervención (B), vuelve a aumentar al retirar la intervención (A) y vuelve a disminuir al reaplicarla (B). ¿Qué conclusión permite este patrón de resultados?

\begin{enumerate}
  \item[A.] Que la conducta agresiva disminuye naturalmente con el tiempo gracias a la maduración
  \item[\textbf{\color{verde}B.}] \textbf{\color{verde} Que la intervención es la causa de los cambios en la conducta, porque el cambio ocurre sistemáticamente solo cuando se aplica} \hfill $\checkmark$
  \item[C.] Que el diseño no es válido porque solo hay un participante y no se puede generalizar
  \item[D.] Que hay una correlación entre la intervención y la conducta, pero no puede establecerse causalidad
\end{enumerate}

\noindent\textbf{\color{azul}Explicación:} La lógica del diseño de caso único de reversión permite inferir causalidad: si la conducta cambia cada vez que la intervención se aplica o se retira, la explicación más parsimoniosa es que la intervención causa el cambio. Este es precisamente el criterio de replicación intrasujeto que León y Montero describen para los diseños N=1.

\noindent\textbf{\color{gris}Por qué son incorrectas las otras opciones:}
\begin{itemize}
  \item {\small A — Incorrecto: si la disminución fuera por maduración, la conducta no volvería a aumentar al retirar la intervención; el patrón de reversión descarta la maduración.}
  \item {\small C — Incorrecto: la generalización es una cuestión de validez externa, no de validez interna; el diseño de caso único sí permite inferir causalidad aunque con un solo sujeto.}
  \item {\small D — Incorrecto: el patrón de reversión sistemática (mejora-empeora-mejora) va más allá de la mera correlación y permite inferir causalidad dentro del caso.}
\end{itemize}
\vspace{8pt}\hrule\vspace{10pt}

\begin{mdframed}[backgroundcolor=fondogris,linecolor=azul,linewidth=1pt,innerleftmargin=10pt,innerrightmargin=10pt,innertopmargin=7pt,innerbottommargin=7pt]
\textbf{Ítem 90 / eg07} \hfill {\small\bfseries\color{azul} Nivel 1 · Básico}
\end{mdframed}

Una investigadora asigna aleatoriamente a 60 niños en dos grupos: el grupo A recibe un programa de habilidades sociales durante 8 semanas y el grupo B realiza actividades de manualidades (condición de control activo). Mide la competencia social antes y después. ¿Por qué es preferible el control activo (manualidades) al control pasivo (sin intervención)?

\begin{enumerate}
  \item[A.] Porque el control activo permite medir la competencia social también en el grupo control
  \item[\textbf{\color{verde}B.}] \textbf{\color{verde} Porque el control activo iguala la cantidad de atención y expectativas entre grupos, controlando el efecto Hawthorne y el efecto de la demanda} \hfill $\checkmark$
  \item[C.] Porque el control activo aumenta el tamaño muestral disponible para el análisis estadístico
  \item[D.] Porque el control activo convierte el diseño en factorial al tener dos tipos de intervención
\end{enumerate}

\noindent\textbf{\color{azul}Explicación:} El grupo de control activo realiza una actividad alternativa del mismo tiempo y formato que el grupo experimental, lo que iguala los efectos inespecíficos: la atención recibida, las expectativas de mejora, la interacción social y el efecto Hawthorne. Sin este control, cualquier mejora en el grupo experimental podría atribuirse a esos factores inespecíficos en lugar de al contenido específico del programa de habilidades sociales.

\noindent\textbf{\color{gris}Por qué son incorrectas las otras opciones:}
\begin{itemize}
  \item {\small A — Incorrecto: en ambas condiciones se mide la competencia social antes y después; la medición no depende del tipo de grupo control.}
  \item {\small C — Incorrecto: el tamaño muestral total (60 niños) es el mismo independientemente del tipo de grupo control; el control activo no añade participantes.}
  \item {\small D — Incorrecto: para que sea un diseño factorial se necesitarían dos variables independientes cruzadas; aquí solo hay una VI (tipo de actividad) con dos niveles.}
\end{itemize}
\vspace{8pt}\hrule\vspace{10pt}

\begin{mdframed}[backgroundcolor=fondogris,linecolor=ambar,linewidth=1pt,innerleftmargin=10pt,innerrightmargin=10pt,innertopmargin=7pt,innerbottommargin=7pt]
\textbf{Ítem 91 / e3} \hfill {\small\bfseries\color{ambar} Nivel 2 · Intermedio}
\end{mdframed}

Un investigador quiere comparar el efecto de dos técnicas de memorización (repetición espaciada vs. recuperación práctica) sobre el recuerdo a largo plazo. Sabe que el nivel de inteligencia previa puede influir en los resultados. Por ello clasifica primero a los participantes en tres bloques según su CI (alto, medio, bajo) y dentro de cada bloque asigna aleatoriamente a las dos técnicas. ¿Qué estrategia de control está aplicando?

\begin{enumerate}
  \item[A.] Contrabalanceo para neutralizar los efectos de orden entre condiciones
  \item[\textbf{\color{verde}B.}] \textbf{\color{verde} Bloqueo, que controla la variabilidad debida al CI introduciendo una variable extraña como factor de agrupación previo} \hfill $\checkmark$
  \item[C.] Diseño factorial 2×2 en el que el CI es la segunda variable independiente
  \item[D.] Igualación postest para comprobar la equivalencia de los grupos tras el estudio
\end{enumerate}

\noindent\textbf{\color{azul}Explicación:} El bloqueo es una estrategia de control que consiste en identificar una variable extraña relevante (aquí, el CI), clasificar a los participantes en bloques según esa variable y luego asignar aleatoriamente dentro de cada bloque. El objetivo es reducir la variabilidad error y aumentar la sensibilidad del diseño, sin que el CI se convierta en una variable independiente manipulada.

\noindent\textbf{\color{gris}Por qué son incorrectas las otras opciones:}
\begin{itemize}
  \item {\small A — Incorrecto: el contrabalanceo se usa en diseños intrasujeto para controlar el orden de presentación de condiciones, no para controlar variables extrañas entre sujetos.}
  \item {\small C — Incorrecto: en el diseño factorial el investigador manipula deliberadamente ambas variables; aquí el CI no se manipula, sino que se controla como variable de bloqueo.}
  \item {\small D — Incorrecto: la igualación postest no existe como técnica de control en diseños experimentales; la equivalencia debe garantizarse antes de la intervención.}
\end{itemize}
\vspace{8pt}\hrule\vspace{10pt}

\begin{mdframed}[backgroundcolor=fondogris,linecolor=ambar,linewidth=1pt,innerleftmargin=10pt,innerrightmargin=10pt,innertopmargin=7pt,innerbottommargin=7pt]
\textbf{Ítem 92 / e4} \hfill {\small\bfseries\color{ambar} Nivel 2 · Intermedio}
\end{mdframed}

Tras un experimento entre grupos, el investigador encuentra que los grupos difieren significativamente en la variable dependiente. Sin embargo, al revisar los datos detecta que el grupo experimental tenía de media 8 años más que el grupo control antes del tratamiento. ¿Qué amenaza a la validez interna explica este hallazgo?

\begin{enumerate}
  \item[A.] Mortalidad experimental, por pérdida diferencial de participantes entre condiciones
  \item[\textbf{\color{verde}B.}] \textbf{\color{verde} Sesgo de selección, porque la aleatorización no generó grupos equivalentes en la variable extraña edad} \hfill $\checkmark$
  \item[C.] Efecto de historia, porque ocurrieron eventos externos entre el pretest y el postest
  \item[D.] Regresión a la media, porque el grupo experimental tenía puntuaciones iniciales extremas
\end{enumerate}

\noindent\textbf{\color{azul}Explicación:} El sesgo de selección ocurre cuando los grupos difieren en variables relevantes antes de la intervención, no como consecuencia de ella. Aunque la aleatorización normalmente previene este sesgo, con muestras pequeñas puede fallar. La solución es verificar la equivalencia inicial con pruebas estadísticas de comparación pretest.

\noindent\textbf{\color{gris}Por qué son incorrectas las otras opciones:}
\begin{itemize}
  \item {\small A — Incorrecto: la mortalidad experimental se refiere a la pérdida de participantes durante el estudio, no a diferencias de edad preexistentes.}
  \item {\small C — Incorrecto: el efecto de historia son eventos externos que ocurren entre medidas; aquí la diferencia existía antes de la intervención.}
  \item {\small D — Incorrecto: la regresión a la media afecta a puntuaciones extremas en la VD, no a diferencias previas en variables extrañas como la edad.}
\end{itemize}
\vspace{8pt}\hrule\vspace{10pt}

\begin{mdframed}[backgroundcolor=fondogris,linecolor=ambar,linewidth=1pt,innerleftmargin=10pt,innerrightmargin=10pt,innertopmargin=7pt,innerbottommargin=7pt]
\textbf{Ítem 93 / v03} \hfill {\small\bfseries\color{ambar} Nivel 2 · Intermedio}
\end{mdframed}

Una investigadora realiza un experimento en el que los participantes, al rellenar el cuestionario pretest, sospechan cuál es el objetivo del estudio y modifican su comportamiento en consecuencia. ¿Qué amenaza a la validez describe esta situación?

\begin{enumerate}
  \item[A.] Efecto de maduración, porque los participantes cambian durante el tiempo del estudio
  \item[\textbf{\color{verde}B.}] \textbf{\color{verde} Efecto del pretestado o reactividad al procedimiento, porque el pretest sensibiliza a los participantes y altera su respuesta posterior} \hfill $\checkmark$
  \item[C.] Mortalidad experimental, porque algunos participantes abandonan el estudio al darse cuenta de su objetivo
  \item[D.] Sesgo de instrumentación, porque el cuestionario no mide correctamente la variable de interés
\end{enumerate}

\noindent\textbf{\color{azul}Explicación:} El efecto del pretestado (pretest sensitization) es una amenaza a la validez externa: cuando los participantes responden a un pretest, pueden volverse sensibles al tratamiento de un modo distinto a quienes no lo reciben. Esto limita la generalización porque en la vida real las personas no suelen recibir "pretests" antes de las experiencias que se estudian.

\noindent\textbf{\color{gris}Por qué son incorrectas las otras opciones:}
\begin{itemize}
  \item {\small A — Incorrecto: la maduración son cambios naturales en los participantes con el tiempo, no el efecto de responder a un cuestionario.}
  \item {\small C — Incorrecto: la mortalidad experimental es la pérdida de participantes; aquí los participantes permanecen pero cambian su comportamiento.}
  \item {\small D — Incorrecto: el sesgo de instrumentación es un problema con el instrumento de medida, no con la reactividad del participante al procedimiento.}
\end{itemize}
\vspace{8pt}\hrule\vspace{10pt}

\begin{mdframed}[backgroundcolor=fondogris,linecolor=ambar,linewidth=1pt,innerleftmargin=10pt,innerrightmargin=10pt,innertopmargin=7pt,innerbottommargin=7pt]
\textbf{Ítem 94 / cu03} \hfill {\small\bfseries\color{ambar} Nivel 2 · Intermedio}
\end{mdframed}

Un terapeuta quiere usar un diseño de caso único para evaluar una intervención de control de impulsos en un adolescente, pero le preocupa que retirar la intervención (pasar de B a A de nuevo) pueda ser éticamente problemático si el adolescente ha mejorado. ¿Qué alternativa metodológica existe al diseño A-B-A-B?

\begin{enumerate}
  \item[A.] El diseño factorial 2×2, aplicando dos intervenciones distintas al mismo adolescente
  \item[\textbf{\color{verde}B.}] \textbf{\color{verde} El diseño de línea de base múltiple, que aplica la intervención de forma escalonada a distintas conductas o contextos sin necesidad de retirarla} \hfill $\checkmark$
  \item[C.] El diseño entre grupos, asignando al adolescente a una condición experimental
  \item[D.] El diseño ex post facto prospectivo, siguiendo al adolescente durante varios años
\end{enumerate}

\noindent\textbf{\color{azul}Explicación:} El diseño de línea de base múltiple (multiple baseline design) aplica la intervención de forma escalonada a distintas conductas, contextos o sujetos. Si cada conducta o contexto mejora solo cuando la intervención se introduce para esa conducta o contexto, se infiere causalidad sin necesidad de retirar la intervención. Es la alternativa estándar cuando la retirada es ética o prácticamente inviable.

\noindent\textbf{\color{gris}Por qué son incorrectas las otras opciones:}
\begin{itemize}
  \item {\small A — Incorrecto: el diseño factorial requiere múltiples participantes y grupos; no es la alternativa para el caso único.}
  \item {\small C — Incorrecto: el diseño entre grupos requiere asignar aleatoriamente a múltiples participantes; no es aplicable a un solo sujeto.}
  \item {\small D — Incorrecto: el diseño prospectivo es un diseño observacional sin intervención; no permite evaluar el efecto de una intervención.}
\end{itemize}
\vspace{8pt}\hrule\vspace{10pt}

\begin{mdframed}[backgroundcolor=fondogris,linecolor=ambar,linewidth=1pt,innerleftmargin=10pt,innerrightmargin=10pt,innertopmargin=7pt,innerbottommargin=7pt]
\textbf{Ítem 95 / cu04} \hfill {\small\bfseries\color{ambar} Nivel 2 · Intermedio}
\end{mdframed}

Una terapeuta usa un diseño A-B con un único paciente: mide la ansiedad durante una semana sin intervención (A) y luego durante cuatro semanas con la terapia (B). Observa una mejora clara en la fase B. Un colega le dice que este diseño tiene una limitación importante. ¿Cuál es?

\begin{enumerate}
  \item[A.] Que no se puede calcular ningún estadístico con solo un participante
  \item[\textbf{\color{verde}B.}] \textbf{\color{verde} Que sin una segunda fase A (retirada) no puede descartarse que la mejora se deba a la maduración u otros factores que no son la terapia} \hfill $\checkmark$
  \item[C.] Que la fase de línea de base (A) de una semana es demasiado larga y produce efectos de práctica
  \item[D.] Que el diseño A-B es exclusivo de investigación de laboratorio y no puede usarse en clínica
\end{enumerate}

\noindent\textbf{\color{azul}Explicación:} El diseño A-B sin retirada tiene una validez interna limitada: sin volver a una fase A, no puede descartarse que la mejora se deba a la maduración natural, al paso del tiempo, al efecto de la relación terapéutica u otros factores ajenos a la intervención específica. La segunda fase A (retirada) es lo que permite descartar estas explicaciones alternativas.

\noindent\textbf{\color{gris}Por qué son incorrectas las otras opciones:}
\begin{itemize}
  \item {\small A — Incorrecto: en los diseños de caso único existen métodos estadísticos específicos (análisis visual, C-statistic, tau-U) para evaluar los datos de un solo participante.}
  \item {\small C — Incorrecto: una línea de base corta puede ser un problema de estabilidad, pero el colega señala un problema de validez interna más fundamental: la ausencia de retirada.}
  \item {\small D — Incorrecto: el diseño A-B (y sus variantes) se usa ampliamente en la práctica clínica y en intervención conductual aplicada.}
\end{itemize}
\vspace{8pt}\hrule\vspace{10pt}

\begin{mdframed}[backgroundcolor=fondogris,linecolor=ambar,linewidth=1pt,innerleftmargin=10pt,innerrightmargin=10pt,innertopmargin=7pt,innerbottommargin=7pt]
\textbf{Ítem 96 / eg08} \hfill {\small\bfseries\color{ambar} Nivel 2 · Intermedio}
\end{mdframed}

Una investigadora realiza un experimento aleatorizado con tres grupos (terapia A, terapia B y lista de espera). Al finalizar, los tres grupos difieren significativamente en los síntomas (F(2,87)=14.2, p<.001). ¿Qué análisis adicional es imprescindible para saber si terapia A es superior a terapia B?

\begin{enumerate}
  \item[A.] Calcular la diferencia de medias entre los tres grupos directamente con un test de chi-cuadrado
  \item[\textbf{\color{verde}B.}] \textbf{\color{verde} Realizar pruebas post hoc (como Tukey HSD o Bonferroni) que comparen sistemáticamente todos los pares de grupos controlando el error tipo I} \hfill $\checkmark$
  \item[C.] Repetir el ANOVA solo con los grupos de terapia A y B, excluyendo el grupo control
  \item[D.] Calcular el tamaño del efecto eta-cuadrado para determinar si la diferencia entre A y B es clínicamente relevante
\end{enumerate}

\noindent\textbf{\color{azul}Explicación:} El ANOVA omnibus solo indica que existe al menos una diferencia significativa entre los grupos, pero no señala cuáles son los pares que difieren. Las pruebas post hoc (Tukey, Bonferroni, Scheffé) comparan cada par de grupos controlando la inflación del error tipo I que resultaría de realizar múltiples pruebas t sin corrección.

\noindent\textbf{\color{gris}Por qué son incorrectas las otras opciones:}
\begin{itemize}
  \item {\small A — Incorrecto: el chi-cuadrado analiza frecuencias de variables categóricas; aquí la VD es continua (síntomas) y se requieren comparaciones de medias.}
  \item {\small C — Incorrecto: repetir el ANOVA solo con A y B sin corrección múltiple es estadísticamente inapropiado y equivale a realizar una prueba t sin control del error tipo I acumulado.}
  \item {\small D — Incorrecto: el eta-cuadrado cuantifica la magnitud del efecto global del factor, pero no indica específicamente si A difiere de B; para eso se necesitan las comparaciones por pares.}
\end{itemize}
\vspace{8pt}\hrule\vspace{10pt}

\begin{mdframed}[backgroundcolor=fondogris,linecolor=ambar,linewidth=1pt,innerleftmargin=10pt,innerrightmargin=10pt,innertopmargin=7pt,innerbottommargin=7pt]
\textbf{Ítem 97 / eg09} \hfill {\small\bfseries\color{ambar} Nivel 2 · Intermedio}
\end{mdframed}

Una investigadora quiere comparar el efecto de dos tipos de retroalimentación (correctiva vs. positiva) sobre el aprendizaje. Asigna aleatoriamente a 80 participantes (40 por condición). Antes del experimento comprueba que los dos grupos son equivalentes en conocimientos previos mediante una prueba t. ¿Qué aporta esta comprobación pretest al diseño?

\begin{enumerate}
  \item[A.] Convierte la asignación aleatoria en innecesaria para el resto del experimento
  \item[\textbf{\color{verde}B.}] \textbf{\color{verde} Proporciona evidencia de que la aleatorización funcionó correctamente y los grupos son comparables al inicio, fortaleciendo la validez interna} \hfill $\checkmark$
  \item[C.] Aumenta automáticamente el poder estadístico para detectar diferencias en la variable dependiente
  \item[D.] Garantiza que no habrá mortalidad experimental diferencial entre los dos grupos
\end{enumerate}

\noindent\textbf{\color{azul}Explicación:} Verificar la equivalencia pretest es una práctica habitual en experimentos para demostrar que la aleatorización ha funcionado. Aunque la aleatorización garantiza la equivalencia en esperanza matemática, con muestras pequeñas pueden producirse desequilibrios por azar. Documentar la equivalencia inicial fortalece la validez interna y hace más creíble la atribución causal de las diferencias postest a la intervención.

\noindent\textbf{\color{gris}Por qué son incorrectas las otras opciones:}
\begin{itemize}
  \item {\small A — Incorrecto: la comprobación pretest no exime de la asignación aleatoria; son procedimientos complementarios con funciones distintas.}
  \item {\small C — Incorrecto: el poder estadístico depende del tamaño muestral, el tamaño del efecto y el nivel alfa; no aumenta por comprobar la equivalencia pretest.}
  \item {\small D — Incorrecto: la mortalidad diferencial depende de factores como la carga de la intervención o la satisfacción de los participantes; la equivalencia en conocimientos previos no la evita.}
\end{itemize}
\vspace{8pt}\hrule\vspace{10pt}

\begin{mdframed}[backgroundcolor=fondogris,linecolor=rojo,linewidth=1pt,innerleftmargin=10pt,innerrightmargin=10pt,innertopmargin=7pt,innerbottommargin=7pt]
\textbf{Ítem 98 / e5} \hfill {\small\bfseries\color{rojo} Nivel 3 · Avanzado}
\end{mdframed}

Un investigador aplica un diseño experimental entre grupos (tratamiento vs. control) con asignación aleatoria. Quiere saber si hay diferencias significativas en la VD entre los dos grupos. Tiene una VD continua y distribución normal. ¿Cuál es el estadístico más adecuado para analizar sus datos?

\begin{enumerate}
  \item[A.] Coeficiente de correlación de Pearson, porque mide la relación entre la VI y la VD
  \item[\textbf{\color{verde}B.}] \textbf{\color{verde} Prueba t de Student para muestras independientes, porque compara medias de dos grupos no relacionados} \hfill $\checkmark$
  \item[C.] ANOVA de medidas repetidas, porque los mismos sujetos están en ambas condiciones
  \item[D.] Chi-cuadrado de Pearson, porque compara frecuencias entre condiciones
\end{enumerate}

\noindent\textbf{\color{azul}Explicación:} La prueba t de Student para muestras independientes es el estadístico adecuado cuando se comparan las medias de dos grupos independientes (asignación aleatoria a condiciones distintas) en una VD continua con distribución normal. Si hubiera más de dos grupos, el análisis sería ANOVA unifactorial entre sujetos.

\noindent\textbf{\color{gris}Por qué son incorrectas las otras opciones:}
\begin{itemize}
  \item {\small A — Incorrecto: la correlación de Pearson mide la fuerza de la relación entre dos variables continuas, no compara medias entre grupos.}
  \item {\small C — Incorrecto: el ANOVA de medidas repetidas se usa cuando los mismos sujetos pasan por todas las condiciones (diseño intrasujeto); aquí cada participante está en solo un grupo.}
  \item {\small D — Incorrecto: el chi-cuadrado analiza frecuencias de variables categóricas; aquí la VD es continua y la pregunta es sobre diferencias de medias.}
\end{itemize}
\vspace{8pt}\hrule\vspace{10pt}

\begin{mdframed}[backgroundcolor=fondogris,linecolor=rojo,linewidth=1pt,innerleftmargin=10pt,innerrightmargin=10pt,innertopmargin=7pt,innerbottommargin=7pt]
\textbf{Ítem 99 / e6} \hfill {\small\bfseries\color{rojo} Nivel 3 · Avanzado}
\end{mdframed}

Una investigadora usa un diseño experimental para evaluar tres niveles de dosis de un fármaco ansiolítico (baja, media, alta) frente a placebo. Asigna aleatoriamente a 80 participantes (20 por condición). Después del tratamiento mide ansiedad (escala continua). ¿Qué análisis permite comparar las cuatro condiciones simultáneamente y, si el resultado es significativo, identificar entre qué pares de condiciones existen diferencias?

\begin{enumerate}
  \item[A.] Cuatro pruebas t independientes, una por cada par de condiciones posible
  \item[\textbf{\color{verde}B.}] \textbf{\color{verde} ANOVA unifactorial entre sujetos seguido de pruebas post hoc (p. ej. Tukey HSD)} \hfill $\checkmark$
  \item[C.] Regresión lineal múltiple con las cuatro condiciones como predictores continuos
  \item[D.] Prueba de Kruskal-Wallis sin análisis posterior
\end{enumerate}

\noindent\textbf{\color{azul}Explicación:} Cuando se comparan más de dos grupos en una VD continua, el ANOVA unifactorial es el procedimiento adecuado porque controla la tasa de error tipo I (falsos positivos) que se inflaría con múltiples pruebas t. Si el ANOVA es significativo, las pruebas post hoc (Tukey, Bonferroni, Scheffé) determinan entre qué pares de grupos existen diferencias.

\noindent\textbf{\color{gris}Por qué son incorrectas las otras opciones:}
\begin{itemize}
  \item {\small A — Incorrecto: realizar múltiples pruebas t sin corrección infla la tasa de error tipo I. Con 4 grupos hay 6 pares posibles; la probabilidad de al menos un falso positivo supera el 25\%.}
  \item {\small C — Incorrecto: la regresión lineal múltiple con cuatro predictores continuos no es la forma estándar de analizar cuatro grupos categóricos independientes.}
  \item {\small D — Incorrecto: la prueba de Kruskal-Wallis es la alternativa no paramétrica; no indica entre qué grupos están las diferencias sin análisis post hoc adicional.}
\end{itemize}
\vspace{8pt}\hrule\vspace{10pt}

\begin{mdframed}[backgroundcolor=fondogris,linecolor=rojo,linewidth=1pt,innerleftmargin=10pt,innerrightmargin=10pt,innertopmargin=7pt,innerbottommargin=7pt]
\textbf{Ítem 100 / eg10} \hfill {\small\bfseries\color{rojo} Nivel 3 · Avanzado}
\end{mdframed}

Una investigadora realiza un ensayo clínico aleatorizado (ECA) para evaluar una nueva terapia para el insomnio. Los participantes se asignan al azar a terapia nueva o terapia estándar. Tanto los participantes como los terapeutas desconocen a qué condición pertenece cada participante (doble ciego). ¿Qué amenaza a la validez interna controla específicamente el procedimiento de doble ciego?

\begin{enumerate}
  \item[A.] La mortalidad experimental diferencial entre las condiciones del estudio
  \item[\textbf{\color{verde}B.}] \textbf{\color{verde} El sesgo de expectativas: tanto los participantes como los evaluadores pueden comportarse de forma diferente si saben a qué condición pertenecen, influyendo en los resultados al margen de la terapia} \hfill $\checkmark$
  \item[C.] La regresión a la media de las puntuaciones extremas de insomnio al inicio
  \item[D.] El efecto de historia producido por eventos externos al experimento
\end{enumerate}

\noindent\textbf{\color{azul}Explicación:} El doble ciego controla el sesgo de expectativas o placebo: los participantes que saben que reciben el tratamiento nuevo pueden mejorar solo por esa expectativa, y los terapeutas que saben qué condición aplican pueden (sin pretenderlo) transmitir distintas expectativas o prestar diferente atención. Al cegar a ambos, estos sesgos quedan neutralizados y los resultados son más atribuibles a la terapia en sí.

\noindent\textbf{\color{gris}Por qué son incorrectas las otras opciones:}
\begin{itemize}
  \item {\small A — Incorrecto: la mortalidad experimental se controla mediante el seguimiento activo de todos los participantes y el análisis por intención de tratar; el ciego no evita que alguien abandone.}
  \item {\small C — Incorrecto: la regresión a la media se controla mediante la aleatorización y el uso de grupos de control; el procedimiento de ciego no actúa sobre este fenómeno estadístico.}
  \item {\small D — Incorrecto: el efecto de historia son eventos externos que ocurren fuera del estudio; el ciego no puede controlar lo que ocurre en el entorno de los participantes.}
\end{itemize}
\vspace{8pt}\hrule\vspace{10pt}

\end{document}